% Options for packages loaded elsewhere
\PassOptionsToPackage{unicode}{hyperref}
\PassOptionsToPackage{hyphens}{url}
\PassOptionsToPackage{dvipsnames,svgnames*,x11names*}{xcolor}
%
\documentclass[
  a4paper,
]{ctexbook}
\usepackage{lmodern}
\usepackage{amssymb,amsmath}
\usepackage{ifxetex,ifluatex}
\ifnum 0\ifxetex 1\fi\ifluatex 1\fi=0 % if pdftex
  \usepackage[T1]{fontenc}
  \usepackage[utf8]{inputenc}
  \usepackage{textcomp} % provide euro and other symbols
\else % if luatex or xetex
  \usepackage{unicode-math}
  \defaultfontfeatures{Scale=MatchLowercase}
  \defaultfontfeatures[\rmfamily]{Ligatures=TeX,Scale=1}
  \setmainfont[]{Noto Serif CJK SC}
  \setsansfont[]{Noto Sans CJK SC}
  \setmonofont[]{DejaVu Sans Mono}
\fi
% Use upquote if available, for straight quotes in verbatim environments
\IfFileExists{upquote.sty}{\usepackage{upquote}}{}
\IfFileExists{microtype.sty}{% use microtype if available
  \usepackage[]{microtype}
  \UseMicrotypeSet[protrusion]{basicmath} % disable protrusion for tt fonts
}{}
\makeatletter
\@ifundefined{KOMAClassName}{% if non-KOMA class
  \IfFileExists{parskip.sty}{%
    \usepackage{parskip}
  }{% else
    \setlength{\parindent}{0pt}
    \setlength{\parskip}{6pt plus 2pt minus 1pt}}
}{% if KOMA class
  \KOMAoptions{parskip=half}}
\makeatother
\usepackage{xcolor}
\IfFileExists{xurl.sty}{\usepackage{xurl}}{} % add URL line breaks if available
\IfFileExists{bookmark.sty}{\usepackage{bookmark}}{\usepackage{hyperref}}
\hypersetup{
  pdftitle={量子动力学数值计算},
  pdfauthor={李皓},
  colorlinks=true,
  linkcolor=teal,
  filecolor=Maroon,
  citecolor=Blue,
  urlcolor=teal,
  pdfcreator={LaTeX via pandoc}}
\urlstyle{same} % disable monospaced font for URLs
\usepackage[margin=2.4cm]{geometry}
\usepackage{color}
\usepackage{fancyvrb}
\newcommand{\VerbBar}{|}
\newcommand{\VERB}{\Verb[commandchars=\\\{\}]}
\DefineVerbatimEnvironment{Highlighting}{Verbatim}{commandchars=\\\{\}}
% Add ',fontsize=\small' for more characters per line
\newenvironment{Shaded}{}{}
\newcommand{\AlertTok}[1]{\textcolor[rgb]{1.00,0.00,0.00}{\textbf{#1}}}
\newcommand{\AnnotationTok}[1]{\textcolor[rgb]{0.38,0.63,0.69}{\textbf{\textit{#1}}}}
\newcommand{\AttributeTok}[1]{\textcolor[rgb]{0.49,0.56,0.16}{#1}}
\newcommand{\BaseNTok}[1]{\textcolor[rgb]{0.25,0.63,0.44}{#1}}
\newcommand{\BuiltInTok}[1]{#1}
\newcommand{\CharTok}[1]{\textcolor[rgb]{0.25,0.44,0.63}{#1}}
\newcommand{\CommentTok}[1]{\textcolor[rgb]{0.38,0.63,0.69}{\textit{#1}}}
\newcommand{\CommentVarTok}[1]{\textcolor[rgb]{0.38,0.63,0.69}{\textbf{\textit{#1}}}}
\newcommand{\ConstantTok}[1]{\textcolor[rgb]{0.53,0.00,0.00}{#1}}
\newcommand{\ControlFlowTok}[1]{\textcolor[rgb]{0.00,0.44,0.13}{\textbf{#1}}}
\newcommand{\DataTypeTok}[1]{\textcolor[rgb]{0.56,0.13,0.00}{#1}}
\newcommand{\DecValTok}[1]{\textcolor[rgb]{0.25,0.63,0.44}{#1}}
\newcommand{\DocumentationTok}[1]{\textcolor[rgb]{0.73,0.13,0.13}{\textit{#1}}}
\newcommand{\ErrorTok}[1]{\textcolor[rgb]{1.00,0.00,0.00}{\textbf{#1}}}
\newcommand{\ExtensionTok}[1]{#1}
\newcommand{\FloatTok}[1]{\textcolor[rgb]{0.25,0.63,0.44}{#1}}
\newcommand{\FunctionTok}[1]{\textcolor[rgb]{0.02,0.16,0.49}{#1}}
\newcommand{\ImportTok}[1]{#1}
\newcommand{\InformationTok}[1]{\textcolor[rgb]{0.38,0.63,0.69}{\textbf{\textit{#1}}}}
\newcommand{\KeywordTok}[1]{\textcolor[rgb]{0.00,0.44,0.13}{\textbf{#1}}}
\newcommand{\NormalTok}[1]{#1}
\newcommand{\OperatorTok}[1]{\textcolor[rgb]{0.40,0.40,0.40}{#1}}
\newcommand{\OtherTok}[1]{\textcolor[rgb]{0.00,0.44,0.13}{#1}}
\newcommand{\PreprocessorTok}[1]{\textcolor[rgb]{0.74,0.48,0.00}{#1}}
\newcommand{\RegionMarkerTok}[1]{#1}
\newcommand{\SpecialCharTok}[1]{\textcolor[rgb]{0.25,0.44,0.63}{#1}}
\newcommand{\SpecialStringTok}[1]{\textcolor[rgb]{0.73,0.40,0.53}{#1}}
\newcommand{\StringTok}[1]{\textcolor[rgb]{0.25,0.44,0.63}{#1}}
\newcommand{\VariableTok}[1]{\textcolor[rgb]{0.10,0.09,0.49}{#1}}
\newcommand{\VerbatimStringTok}[1]{\textcolor[rgb]{0.25,0.44,0.63}{#1}}
\newcommand{\WarningTok}[1]{\textcolor[rgb]{0.38,0.63,0.69}{\textbf{\textit{#1}}}}
\usepackage{longtable,booktabs}
% Correct order of tables after \paragraph or \subparagraph
\usepackage{etoolbox}
\makeatletter
\patchcmd\longtable{\par}{\if@noskipsec\mbox{}\fi\par}{}{}
\makeatother
% Allow footnotes in longtable head/foot
\IfFileExists{footnotehyper.sty}{\usepackage{footnotehyper}}{\usepackage{footnote}}
\makesavenoteenv{longtable}
\usepackage{graphicx}
\makeatletter
\def\maxwidth{\ifdim\Gin@nat@width>\linewidth\linewidth\else\Gin@nat@width\fi}
\def\maxheight{\ifdim\Gin@nat@height>\textheight\textheight\else\Gin@nat@height\fi}
\makeatother
% Scale images if necessary, so that they will not overflow the page
% margins by default, and it is still possible to overwrite the defaults
% using explicit options in \includegraphics[width, height, ...]{}
\setkeys{Gin}{width=\maxwidth,height=\maxheight,keepaspectratio}
% Set default figure placement to htbp
\makeatletter
\def\fps@figure{htbp}
\makeatother
\setlength{\emergencystretch}{3em} % prevent overfull lines
\providecommand{\tightlist}{%
  \setlength{\itemsep}{0pt}\setlength{\parskip}{0pt}}
\setcounter{secnumdepth}{-\maxdimen} % remove section numbering

\title{量子动力学数值计算}
\usepackage{etoolbox}
\makeatletter
\providecommand{\subtitle}[1]{% add subtitle to \maketitle
  \apptocmd{\@title}{\par {\large #1 \par}}{}{}
}
\makeatother
\subtitle{基于 GeneralModule 的 Fortran 2008 理论与实践}
\author{李皓}
\date{2026-09-25}

\begin{document}
\maketitle

{
\hypersetup{linkcolor=}
\setcounter{tocdepth}{3}
\tableofcontents
}
\hypertarget{ux5173ux4e8eux672cux4e66}{%
\section{关于本书}\label{ux5173ux4e8eux672cux4e66}}

本书围绕 GeneralModule
开源代码库，系统讨论量子动力学问题的数学表述、离散化方法、数值算法和工程实现。本书强调理论、算法与代码验证的完整链条。

本书对应的代码仓库为：

\begin{itemize}
\item
  \url{https://github.com/l1Ha/QuantumGeneralModule}
\item
  仓库二维码：
\end{itemize}

\includegraphics[width=2cm,height=\textheight]{qr/repo_root.png}

\hypertarget{ux7b26ux53f7ux4e0eux7ea6ux5b9a}{%
\section{符号与约定}\label{ux7b26ux53f7ux4e0eux7ea6ux5b9a}}

除特别说明外，本书使用原子单位：

\[
\hbar=1,\qquad m_e=1,\qquad e=1,\qquad 4\pi\varepsilon_0=1.
\]

能量单位为 Hartree，长度单位为 Bohr。

\hypertarget{generalmodule-ux7b2cux56dbux7248-ux6269ux5c55ux7ae0ux8282ux7b2cux4e00ux90e8ux5206}{%
\section{GeneralModule 第四版 ·
扩展章节（第一部分）}\label{generalmodule-ux7b2cux56dbux7248-ux6269ux5c55ux7ae0ux8282ux7b2cux4e00ux90e8ux5206}}

\hypertarget{ux91cfux5b50ux529bux5b66ux6570ux5b66ux7ed3ux6784ux89d2ux52a8ux91cfux4ee3ux6570ux5206ux5b50ux5750ux6807ux7406ux8bbaux4e0eux6570ux503cux65b9ux6cd5ux57faux7840}{%
\subsection{量子力学数学结构、角动量代数、分子坐标理论与数值方法基础}\label{ux91cfux5b50ux529bux5b66ux6570ux5b66ux7ed3ux6784ux89d2ux52a8ux91cfux4ee3ux6570ux5206ux5b50ux5750ux6807ux7406ux8bbaux4e0eux6570ux503cux65b9ux6cd5ux57faux7840}}

本章为 GeneralModule
第四版新增的理论扩展章节，面向计算量子动力学、分子光谱学与超冷碰撞领域的研究生读者及算法开发者。全章沿四条主线递进展开：第一章论述希尔伯特空间与算符谱理论，这是有限基组与离散变量表象方法的数学根基；第二章系统建立角动量代数、Clebsch--Gordan
系数与 Wigner 3j/6j/9j
符号的重耦合理论；第三章讨论空间固定坐标系与体固定坐标系、Euler
角、Jacobi 坐标、超球坐标，以及 Watson 转振哈密顿量、Coriolis
耦合、离心畸变与振动角动量；第四章阐述单位制、实对称本征问题的
Householder--QL 算法、条件数、误差传播与收敛阶。

每一节均按统一体例组织：\textbf{定义}、\textbf{公式与推导}、\textbf{物理含义}、\textbf{数值陷阱}、\textbf{GeneralModule
实现映射}（源码路径、函数、GitHub
直链与二维码）、\textbf{小型数值实验}与\textbf{练习}。所有数值实验结果均由本库源码直接编译运行获得（GNU
Fortran 11.4，\texttt{real64}
双精度，随机数种子取默认值），可复现、可作为读者自建验收测试的基线。行文约定：公式采用
Pandoc 兼容的 LaTeX 记号，行内公式以 \(...\) 界定，独立公式以 \[...\]
界定。

\begin{center}\rule{0.5\linewidth}{0.5pt}\end{center}

\hypertarget{ux7b2cux4e00ux7ae0-ux5e0cux5c14ux4f2fux7279ux7a7aux95f4ux4e0eux7b97ux7b26ux8c31ux7406ux8bba}{%
\subsection{第一章
希尔伯特空间与算符谱理论}\label{ux7b2cux4e00ux7ae0-ux5e0cux5c14ux4f2fux7279ux7a7aux95f4ux4e0eux7b97ux7b26ux8c31ux7406ux8bba}}

量子力学的全部数学结构建立在内积空间及其完备化之上。计算量子动力学的每一个数值算法------无论是有限基组展开、离散变量表象（DVR）还是波包传播------都是在希尔伯特空间的某个有限维截断子空间内工作。因此，严格理解无限维理论与有限维近似之间的对应与偏差，是判别数值结果可靠性的前提。本节内容由浅入深：先给出内积与完备性的基本定义，继而讨论基组展开与表象变换，然后陈述自伴算符的谱定理，最后落实到幺正演化的数值实现。

\hypertarget{ux5185ux79efux8303ux6570ux4e0eux5e0cux5c14ux4f2fux7279ux7a7aux95f4}{%
\subsubsection{1.1
内积、范数与希尔伯特空间}\label{ux5185ux79efux8303ux6570ux4e0eux5e0cux5c14ux4f2fux7279ux7a7aux95f4}}

\textbf{定义。} 设 \(\mathcal{H}\) 为复数域 \(\mathbb{C}\)
上的线性空间。若映射
\(\langle\,\cdot\,|\,\cdot\,\rangle:\mathcal{H}\times\mathcal{H}\to\mathbb{C}\)
满足对第一变元共轭线性、对第二变元线性、共轭对称
\(\langle\psi|\phi\rangle=\langle\phi|\psi\rangle^{*}\) 以及正定性
\(\langle\psi|\psi\rangle\ge 0\)（等号当且仅当 \(\psi=0\)），则称其为
\(\mathcal{H}\) 上的内积。由内积诱导范数
\(\|\psi\|=\sqrt{\langle\psi|\psi\rangle}\)，并满足 Cauchy--Schwarz
不等式 \(|\langle\psi|\phi\rangle|\le\|\psi\|\,\|\phi\|\)
与三角不等式。若 \(\mathcal{H}\) 在该范数下完备，即任意 Cauchy 序列
\(\{\psi_n\}\)（满足 \(\|\psi_n-\psi_m\|\to 0\)）都收敛于
\(\mathcal{H}\) 中元素，则称 \(\mathcal{H}\)
为希尔伯特空间。量子力学标准模型进一步要求 \(\mathcal{H}\)
可分，即存在可数的稠密子集。

\textbf{公式与推导。} 微观体系中最常见的两类实现为：有限维空间
\(\mathbb{C}^{n}\)，内积
\(\langle u|v\rangle=\sum_i u_i^{*}v_i\)；平方可积函数空间
\(L^{2}(\mathbb{R}^{3})\)，内积
\(\langle\psi|\phi\rangle=\int_{\mathbb{R}^{3}}\psi^{*}(\mathbf{r})\phi(\mathbf{r})\,d^{3}\mathbf{r}\)，其元素满足
\(\int|\psi|^{2}\,d^{3}\mathbf{r}<\infty\)。Cauchy--Schwarz
不等式的证明取 \(\lambda\in\mathbb{C}\) 并利用
\(0\le\|\psi+\lambda\phi\|^{2}=\|\psi\|^{2}+\lambda\langle\psi|\phi\rangle+\lambda^{*}\langle\phi|\psi\rangle+|\lambda|^{2}\|\phi\|^{2}\)：当
\(\phi\neq 0\) 时取
\(\lambda=-\langle\phi|\psi\rangle/\|\phi\|^{2}\)，判别式条件立即给出
\(|\langle\psi|\phi\rangle|^{2}\le\|\psi\|^{2}\|\phi\|^{2}\)。对散射问题，严格而言连续谱态（平面波、散射波）不属于
\(L^{2}\)，须借助装备希尔伯特空间
\(\Phi\subset\mathcal{H}\subset\Phi^{\times}\)（rigged Hilbert
space）处理；数值计算中则以大盒子边界条件将连续谱离散化，使其纳入可分希尔伯特空间框架。

\textbf{物理含义。} 量子态是 \(\mathcal{H}\) 中的矢量，测量概率由 Born
规则 \(p(a)=\langle\psi|P_{a}|\psi\rangle\) 给出；两态之间的交叠
\(\langle\psi|\phi\rangle\)
度量其相干叠加的强度，其模方在光谱学中对应跃迁线强度中的重叠因子。范数平方
\(\|\psi\|^{2}=1\)
表述总概率守恒，是任何数值传播格式必须保持的第一守恒量。

\textbf{数值陷阱。} 其一，离散化后内积的度量约定不一：将 \(L^{2}\) 内积
\(\int\psi^{*}\phi\,dx\) 以梯形或矩形公式离散，得到
\(\sum_i\psi_i^{*}\phi_i\,\Delta x\)；而本征求解器（如本库
\texttt{fgh\_solve\_bound\_states}）返回的本征矢量在``格点 Kronecker
度量'' \(\sum_i z_i^{*}z_i'=1\) 下归一化，物理波函数须按
\(\psi(x_i)=z_i/\sqrt{\Delta x}\)
还原。混淆两种度量是转振计算中出现量级为 \(\Delta x\)
的系统性错误的头号来源，3.4
节的数值实验将实际演示该错误（转动常数被缩小约 90
倍）。其二，长时间传播中归一化以 \(\mathcal{O}(\varepsilon)\)
每步的速度漂移，须定期监测 \(\langle\psi|\psi\rangle\)。

\textbf{GeneralModule 实现映射。} 源码
\texttt{src/mod\_dvr\_grid.f90}：派生类型 \texttt{dvr\_1d\_t}（格点坐标
\texttt{x}、动能矩阵
\texttt{t\_mat}）、\texttt{dvr\_expectation\_value}（以
\(\sum_i|\psi_i|^{2}\,O(x_i)\,\Delta x\)
计算期望值）、\texttt{dvr\_matrix\_element}（两态跃迁矩阵元）。GitHub
直链：\href{https://github.com/l1Ha/QuantumGeneralModule/blob/main/src/mod_dvr_grid.f90}{src/mod\_dvr\_grid.f90}。二维码：\texttt{qr/src\_\_mod\_dvr\_grid.f90.png}。

\begin{figure}
\centering
\includegraphics{qr/src__mod_dvr_grid.f90.png}
\caption{DVR 网格模块二维码}
\end{figure}

\textbf{小型数值实验。} 取 Gauss 波包
\(\psi(x)=\pi^{-1/4}\exp(-x^{2}/2)\)，在区间 \([-10,10]\) 上以 Sinc-DVR
网格（\(n=256\)）离散并计算 \(\sum_i|\psi(x_i)|^{2}\Delta x\)，所得值为
\(0.999999999997\)，偏差 \(3\times 10^{-12}\)，与 \(L^{2}\)
内积的求积误差一致；若省略权重 \(\Delta x\)，结果变为
\(0.0323\)，即偏差达两个数量级以上，直观演示度量约定的重要性。

\textbf{练习。} (1) 证明平行四边形定律
\(\|\psi+\phi\|^{2}+\|\psi-\phi\|^{2}=2\|\psi\|^{2}+2\|\phi\|^{2}\)，并说明内积空间的范数必然满足它。(2)
对三维谐振子基态波函数
\(\psi(\mathbf{r})=(\alpha/\pi)^{3/4}e^{-\alpha r^{2}/2}\)
计算归一化常数并验证 \(\langle\psi|\psi\rangle=1\)。(3)
说明散射态为何不属于 \(L^{2}\)，并解释盒子归一化 \(\psi\propto\sin(kr)\)
如何在 \(L\to\infty\) 时恢复连续谱。

\hypertarget{ux5b8cux5907ux6027ux57faux7ec4ux5c55ux5f00ux4e0eux8868ux8c61ux53d8ux6362}{%
\subsubsection{1.2
完备性、基组展开与表象变换}\label{ux5b8cux5907ux6027ux57faux7ec4ux5c55ux5f00ux4e0eux8868ux8c61ux53d8ux6362}}

\textbf{定义。} 设 \(\{\varphi_n\}_{n=1}^{\infty}\subset\mathcal{H}\)
为正交归一序列，即
\(\langle\varphi_m|\varphi_n\rangle=\delta_{mn}\)。若对任意
\(\psi\in\mathcal{H}\) 均有
\(\psi=\sum_n c_n\varphi_n\)（按范数收敛），\(c_n=\langle\varphi_n|\psi\rangle\)，则称
\(\{\varphi_n\}\) 为完备正交归一基。完备性等价于单位分解（resolution of
identity）

\[\sum_{n=1}^{\infty}|\varphi_n\rangle\langle\varphi_n|=\mathbb{1},\]

也等价于 Parseval（Plancherel）等式
\(\|\psi\|^{2}=\sum_n|c_n|^{2}\)。给定基组即选定一个表象；基组间的幺正矩阵
\(U_{mn}=\langle\varphi_m|\chi_n\rangle\) 实现表象变换，算符矩阵按
\(A^{(\chi)}=U^{\dagger}A^{(\varphi)}U\) 变换。

\textbf{公式与推导。} 由单位分解立即得到 Parseval
等式：\(\|\psi\|^{2}=\langle\psi|\mathbb{1}|\psi\rangle=\sum_n|\langle\varphi_n|\psi\rangle|^{2}\)。变分原理给出截断近似的理论保障：对任意归一化
\(\psi\)，\(\langle\psi|\hat{H}|\psi\rangle\ge E_0\)（Ritz
变分原理），因此在 \(N\) 维截断基中对角化 \(\hat{H}\) 得到的最低本征值
\(E_0^{(N)}\)
单调不增地自上方逼近精确基态能量；Hylleraas--Undheim--MacDonald
定理进一步保证第 \(k\) 条能级 \(E_k^{(N)}\)
随基组增大单调不增且与精确谱交错。对于 Sinc-DVR，Shannon
采样定理保证带限函数的展开是精确的，动能矩阵元具有解析形式

\[T_{ij}=\frac{\hbar^{2}}{2\mu\,\Delta x^{2}}
\begin{cases}
\dfrac{\pi^{2}}{3}, & i=j,\\[2mm]
\dfrac{2(-1)^{i-j}}{(i-j)^{2}}, & i\neq j,
\end{cases}\]

这正是本库 \texttt{dvr\_sinc\_init} 所实现的公式（Colbert--Miller,
1992）。

\textbf{物理含义。}
选取表象对应实验上的测量基选择：坐标表象波函数描述空间概率幅，能量表象展开系数描述定态布居，角动量表象描述取向分布。Fourier
网格哈密顿量（FGH）与 DVR
方法的物理内涵是：在正交格点上以``离散坐标本征态''近似连续坐标本征态，势能算符严格对角，动能算符解析非对角，二者相加后对角化即得束缚能级。

\textbf{数值陷阱。}
其一，截断基组的完备性只在变分子空间内成立：若势能面在网格外仍有可观振幅，能级不随网格增大收敛（基组套设错误，box
error）；近解离能级尤其敏感，3.4 节实验中 Morse 势在盒子
\([0.4,6.0]\,a_0\) 内只得 15 条束缚能级，而半经典估计为 17
条，即为例证。其二，Fourier 网格存在 Nyquist 混叠：波数
\(|k|>\pi/\Delta x\) 的分量被折叠为低频伪影，故网格间距必须满足
\(\Delta x\le\pi/k_{\max}\)。其三，若动能矩阵元以数值微分构造而非解析式，破坏了厄米性（\(T_{ij}\neq T_{ji}^{*}\)），变分上界性质随之失效。

\textbf{GeneralModule 实现映射。} 源码
\texttt{src/mod\_dvr\_grid.f90}：\texttt{dvr\_sinc\_init}（Colbert--Miller
Sinc-DVR 动能矩阵）、\texttt{fgh\_solve\_bound\_states}（FGH
束缚态求解，组装 \(H_{ij}=T_{ij}+V(x_i)\delta_{ij}\)
并调用对称对角化）、\texttt{dvr\_legendre\_init}（Gauss--Legendre 角向
DVR）。文献依据见 Colbert 与 Miller（J. Chem. Phys. 96, 1982 (1992),
DOI: 10.1063/1.462125）以及 Marston 与 Balint-Kurti（J. Chem. Phys. 91,
3571 (1989), DOI: 10.1063/1.456888）。GitHub
直链：\href{https://github.com/l1Ha/QuantumGeneralModule/blob/main/src/mod_dvr_grid.f90}{src/mod\_dvr\_grid.f90}。二维码：\texttt{qr/src\_\_mod\_dvr\_grid.f90.png}。

\textbf{小型数值实验。} 以 1.1 节的同一线型分子势（见 3.4
节参数）考察网格收敛性：将网格点数从 \(n=250\) 增至 \(n=500\)，基态能量
\(E_0\) 的变化小于 \(10^{-9}\ E_h\)，表明低能级已收敛；而第 15
条能级随盒子从 \(6.0\ a_0\) 扩至 \(8.0\ a_0\) 下移约
\(4\ \mathrm{cm^{-1}}\)，显示近解离能级的盒子误差。此实验说明``逐条能级、逐个参数地检验收敛''是变分计算不可省略的工序。

\textbf{练习。} (1) 由单位分解证明 Parseval 等式。(2) 证明 Sinc 函数族
\(\mathrm{sinc}[\pi(x-x_i)/\Delta x]\) 在带限函数空间内正交归一。(3)
对宽度为 \(L\) 的无限深势阱，估计用 Sinc-DVR 重现前 10
条能级所需的最小格距。

\hypertarget{ux81eaux4f34ux7b97ux7b26ux4e0eux8c31ux5b9aux7406}{%
\subsubsection{1.3
自伴算符与谱定理}\label{ux81eaux4f34ux7b97ux7b26ux4e0eux8c31ux5b9aux7406}}

\textbf{定义。} 稠密定义域 \(\mathcal{D}(A)\) 上的线性算符 \(A\)
称为对称的（厄米的），若
\(\langle\psi|A\phi\rangle=\langle A\psi|\phi\rangle\) 对所有
\(\psi,\phi\in\mathcal{D}(A)\) 成立；若进一步满足
\(\mathcal{D}(A)=\mathcal{D}(A^{\dagger})\)，则称 \(A\)
为自伴的（self-adjoint）。在有限维空间中二者无区别，但在无限维空间中定义域之差是本质的。自伴算符的谱由离散谱（孤立有限重数本征值）、连续谱与残余谱组成；谱定理断言：自伴算符
\(A\) 幺正等价于其谱测度，即

\[A=\int_{\sigma(A)}\lambda\,dE_{\lambda},\qquad \mathbb{1}=\int_{\sigma(A)}dE_{\lambda},\]

其中 \(\{E_{\lambda}\}\) 为投影值测度；对纯离散谱退化为
\(A=\sum_k\lambda_k|k\rangle\langle k|\)。

\textbf{公式与推导。} 有限维谱定理可用 Rayleigh
商的极值刻画（Courant--Fischer
极小极大原理）：\(E_0=\min_{\psi\neq0}\langle\psi|A|\psi\rangle/\langle\psi|\psi\rangle\)，取极小的
\(\psi\) 满足变分方程
\(\delta\{\langle\psi|A|\psi\rangle-\lambda\langle\psi|\psi\rangle\}=0\)，即
\(A\psi=\lambda\psi\)；对子空间逐级重复即得全部本征对。自伴性给出三条直接推论：本征值全为实数（\(\lambda\langle\psi|\psi\rangle=\langle\psi|A|\psi\rangle\in\mathbb{R}\)）；属于不同本征值的本征矢正交；任意对称实矩阵可由正交矩阵对角化
\(A=Z\Lambda Z^{T}\)。反例说明定义域的重要性：半直线 \([0,\infty)\)
上的动量算符 \(-i\hbar\,d/dx\)
在满足边界条件的定义域上对称但不自伴，其本征值可任意复化，不存在物理可观测对应的谱分解。

\textbf{物理含义。}
可观测量与自伴算符一一对应：能量、角动量分量、偶极矩均为自伴算符，故其测量值实、本征态可构成正交完备基、含时演化保概率。数值上，DVR
或基组方法组装的哈密顿矩阵正是自伴算符在有限维子空间的投影，其对角化结果（能级与波函数）构成一切后续光谱与动力学分析的原材料。复吸收势（CAP）则刻意将哈密顿变为非厄米以吸收出射流，此时谱不再全实，须以非厄米谱理论解释（参见本库
\texttt{mod\_absorbing\_boundary.f90}）。

\textbf{数值陷阱。}
其一，矩阵组装代码只填上三角或只填下三角而忘记镜像，得到形式上``对称''实则不对称的输入，\texttt{tred2}/\texttt{tql2}
假定对称性后给出的结果不可预测；验收办法是显式检验
\(\max_{ij}|A_{ij}-A_{ji}|=0\)。其二，简并本征值对应的本征矢在对角化后是简并子空间内任意正交基，依赖于舍入噪声，跨程序比对本征矢无意义，比对投影算符
\(\sum_{k\in\text{简并}}|k\rangle\langle k|\) 才有意义。其三，投影矩阵
\(A=PAP^{\dagger}\)
仍是自伴的，但若投影后基组不正交（如重叠矩阵未同时处理），广义本征问题
\(Hc=ESc\) 被误当作标准问题对角化，能级系统性偏离。

\textbf{GeneralModule 实现映射。} 源码
\texttt{src/mod\_linear\_algebra.f90}：\texttt{diag\_symmetric\_matrix}（三对角化
\texttt{tred2} 加 QL 迭代 \texttt{tql2} 的统一入口，输出升序本征值 \(d\)
与正交本征矢矩阵
\(z\)）、\texttt{inv\_real\_matrix}/\texttt{inv\_complex\_matrix}（全主元
Gauss--Jordan 求逆，\texttt{stat=-1} 指示奇异矩阵）。GitHub
直链：\href{https://github.com/l1Ha/QuantumGeneralModule/blob/main/src/mod_linear_algebra.f90}{src/mod\_linear\_algebra.f90}。二维码：\texttt{qr/src\_\_mod\_linear\_algebra.f90.png}。

\begin{figure}
\centering
\includegraphics{qr/src__mod_linear_algebra.f90.png}
\caption{线性代数模块二维码}
\end{figure}

\textbf{小型数值实验。} 取 \(n=300\) 的随机实对称矩阵（元素均匀分布于
\([-1,1]\)），调用 \texttt{diag\_symmetric\_matrix} 后逐列计算残余
\(\|A z_k-\lambda_k z_k\|_{\infty}\) 与正交性偏差
\(\max_{ij}|(Z^{T}Z-\mathbb{1})_{ij}|\)。实测结果：最大残余
\(4.4\times10^{-14}\)，最大正交性偏差 \(2.2\times10^{-14}\)。二者均为
\(\mathcal{O}(n\varepsilon\|A\|)\)
量级（\(\varepsilon=2.2\times10^{-16}\)，\(n\varepsilon\approx6.6\times10^{-14}\)），符合舍入误差的预期标度，说明算法达到后向稳定。

\textbf{练习。} (1) 证明有限维空间中对称算符必然自伴。(2)
验证半直线上动量算符的对称性并找出使其非自伴的边界条件缺口。(3) 对
\(n=10\) 的随机对称矩阵计算投影算符残差
\(\|\sum_k|k\rangle\langle k|-\mathbb{1}\|_{F}\) 并解释其量级。

\hypertarget{ux5e7aux6b63ux6f14ux5316ux4e0eux6ce2ux5305ux4f20ux64ad}{%
\subsubsection{1.4
幺正演化与波包传播}\label{ux5e7aux6b63ux6f14ux5316ux4e0eux6ce2ux5305ux4f20ux64ad}}

\textbf{定义。} 算符 \(U\) 称为幺正的，若
\(U^{\dagger}U=U U^{\dagger}=\mathbb{1}\)。含时薛定谔方程

\[i\hbar\frac{\partial}{\partial t}|\psi(t)\rangle=\hat{H}(t)|\psi(t)\rangle\]

在 \(\hat{H}\) 不含时时的形式解为
\(|\psi(t)\rangle=U(t)|\psi(0)\rangle\)，\(U(t)=\exp(-i\hat{H}t/\hbar)\)；Stone
定理断言：强连续单参数幺正群 \(U(t)\)
的生成元必为自伴算符，反之亦然。幺正性使 \(\|\psi(t)\|\)
守恒，对应概率守恒。

\textbf{公式与推导。} 数值传播的核心是对 \(U(\Delta t)\)
的高精度近似。Baker--Campbell--Hausdorff（BCH）展开给出对称分裂（Trotter--Strang
分裂）

\[e^{-i(\hat{T}+\hat{V})\Delta t/\hbar}
=e^{-i\hat{V}\Delta t/2\hbar}\,e^{-i\hat{T}\Delta t/\hbar}\,e^{-i\hat{V}\Delta t/2\hbar}
+\mathcal{O}(\Delta t^{3}),\]

误差首项为
\(-\frac{i\Delta t^{3}}{24\hbar^{3}}[\hat{V},[\hat{V},\hat{T}]]+\frac{i\Delta t^{3}}{24\hbar^{3}}[\hat{T},[\hat{T},\hat{V}]]\)
型双对易子。由于三个因子均为幺正，对称分裂格式对任意步长都严格幺正------这不是近似而是精确性质，因此范数在机器精度内守恒；其三阶局部误差对束缚体系表现为相位误差而非振幅误差，长时间不产生久期增长。在坐标网格上，\(\hat{V}\)
对角、\(\hat{T}\)
在动量表象对角，故每次推进只需两次快速傅里叶变换；对大步长与长时间演化，Chebyshev
多项式展开
\(e^{-i\hat{H}t/\hbar}\approx\sum_{n}a_n T_n(\hat{H}_{\mathrm{norm}})\)
提供谱收敛的替代方案（Kosloff, 1988）。

\textbf{物理含义。}
波包传播直接给出光解通量、散射矩阵、含时对齐度等可观测量；幺正性对应封闭体系的概率守恒，任何非幺正格式（如未加修正的显式欧拉法）都会在数千步内累积出物理上无意义的增益或衰减。分裂算符格式的辛对称性还保证了相空间体积守恒的量子对应------能级布居在长时间平均下不漂移。

\textbf{数值陷阱。} 其一，FFT
隐含周期边界：波包到达网格边缘后从另一侧回卷（wrap-around），须以足够大的盒子或复吸收势（本库
\texttt{mod\_absorbing\_boundary.f90}）消除。其二，时间步长须同时满足分裂误差与Nyquist
条件：\(\Delta t\) 过大时高动量分量的相位 \(e^{-iT(p)\Delta t/\hbar}\)
采样不足，产生高频混叠。其三，尽管格式严格幺正，浮点舍入仍以每步
\(\mathcal{O}(\varepsilon)\)
引入微小的范数漂移，万步量级后可观测；宜在传播中周期性检查
\(\langle\psi|\psi\rangle\) 并记录。

\textbf{GeneralModule 实现映射。} 源码
\texttt{src/mod\_wavepacket\_propagator.f90}：\texttt{propagate\_split\_operator\_1d}（坐标半步势能相位、FFT
至动量空间、动能整步相位、逆变换、势能后半步，内部调用
\texttt{mod\_linear\_algebra.f90} 的
\texttt{fft\_1d}）、\texttt{propagate\_split\_operator\_2d}、\texttt{rk4\_step}、\texttt{abm4\_step}；\texttt{src/mod\_chebyshev\_propagator.f90}
提供切比雪夫大步长推进器。文献依据：Feit, Fleck 与 Steiger（J. Comput.
Phys. 47, 412 (1982), DOI: 10.1016/0021-9991(82)90091-2）、Kosloff（J.
Phys. Chem. 92, 2087 (1988), DOI: 10.1021/j100319a003）。GitHub
直链：\href{https://github.com/l1Ha/QuantumGeneralModule/blob/main/src/mod_wavepacket_propagator.f90}{src/mod\_wavepacket\_propagator.f90}。二维码：\texttt{qr/src\_\_mod\_wavepacket\_propagator.f90.png}。

\begin{figure}
\centering
\includegraphics{qr/src__mod_wavepacket_propagator.f90.png}
\caption{波包传播模块二维码}
\end{figure}

\textbf{小型数值实验。}
一维谐振子（\(m=1\)，\(\omega=0.02\ E_h/\hbar\)），网格
\(n=1024\)、\(\Delta x=0.05\ a_0\)、\(\Delta t=0.5\ a.u.\)，以高斯波包为初态连续推进
\(20000\) 步（相当于
\(10000\ a.u.\approx242\ \mathrm{fs}\)）。实测终态范数偏差
\(|\ \|\psi\|^{2}-1\ |=8.8\times10^{-12}\)，确认了格式的严格幺正性与舍入漂移的
\(\mathcal{O}(N_{\mathrm{step}}\varepsilon)\) 标度。

\textbf{练习。} (1) 用 BCH 展开推导对称分裂的
\(\mathcal{O}(\Delta t^{3})\) 误差项并证明其为厄米算符的虚倍数。(2) 证明
\(e^{-i\hat{V}\Delta t/2\hbar}e^{-i\hat{T}\Delta t/\hbar}e^{-i\hat{V}\Delta t/2\hbar}\)
对任意 \(\Delta t\) 幺正。(3) 对谐振子精确解比较分裂格式的相位误差随
\(\Delta t\) 的标度，并估计达到 \(10^{-6}\) 相位精度所需步长。 \#
第二部分　谱方法、离散表示与收敛理论

\hypertarget{ux4eceux57faux7ec4ux5c55ux5f00ux5230ux77e9ux9635ux95eeux9898}{%
\subsection{2.1　从基组展开到矩阵问题}\label{ux4eceux57faux7ec4ux5c55ux5f00ux5230ux77e9ux9635ux95eeux9898}}

设哈密顿量为

\[
\hat H=\hat T+\hat V.
\]

将态展开为

\[
|\psi\rangle=\sum_{n=1}^{N}c_n|\chi_n\rangle.
\]

代入薛定谔方程并在基组上投影，得到广义本征值问题

\[
\sum_j H_{ij}c_j
=
E\sum_j S_{ij}c_j,
\]

其中

\[
H_{ij}=\langle \chi_i|\hat H|\chi_j\rangle,
\qquad
S_{ij}=\langle \chi_i|\chi_j\rangle.
\]

若基组正交，则 \(S_{ij}=\delta_{ij}\)，问题退化为标准本征值问题

\[
H\mathbf c=E\mathbf c.
\]

若基组非正交但线性无关，可以对重叠矩阵作 Cholesky 分解

\[
S=LL^T,
\]

并定义

\[
\tilde H=L^{-1}H L^{-T}.
\]

广义问题变为对称标准问题

\[
\tilde H\mathbf y=E\mathbf y,
\qquad
\mathbf c=L^{-T}\mathbf y.
\]

数值上要避免直接求 \(S^{-1}\)，因为当基组近线性相关时，\(S\)
的最小本征值接近零，\(\kappa(S)\) 会急剧增大。

\hypertarget{dvr-ux7684ux57faux672cux7ed3ux6784}{%
\subsection{2.2　DVR
的基本结构}\label{dvr-ux7684ux57faux672cux7ed3ux6784}}

离散变量表象不是一种具体算法，而是一类满足以下条件的表示：

\begin{enumerate}
\def\labelenumi{\arabic{enumi}.}
\tightlist
\item
  网格点 \(\{x_i\}_{i=1}^{N}\) 对应正交函数或正交多项式；
\item
  离散内积由权重 \(\{w_i\}\) 定义；
\item
  局域算符在网格上近似对角；
\item
  动能矩阵由解析或半解析公式给出。
\end{enumerate}

离散内积为

\[
\langle f|g\rangle_N
=
\sum_{i=1}^{N}w_i f(x_i)^{*}g(x_i).
\]

对局域势能算符，

\[
\langle f|\hat V|g\rangle
=
\int f^{*}(x)V(x)g(x)\,dx
\approx
\sum_i w_i f(x_i)^{*}V(x_i)g(x_i),
\]

因此

\[
V_{ij}=V(x_i)\delta_{ij}.
\]

这是 DVR
的核心优势。它把高维问题中的势能矩阵由稠密矩阵变成对角矩阵，使矩阵存储从
\(O(N^2)\) 降至 \(O(N)\)。

\hypertarget{fourier-ux7f51ux683cux4e0e-fgh}{%
\subsection{2.3　Fourier 网格与
FGH}\label{fourier-ux7f51ux683cux4e0e-fgh}}

均匀周期域 \(x\in[0,L)\) 上的 Fourier 基为

\[
\phi_k(x)=\frac{1}{\sqrt L}e^{ikx},
\qquad
k=\frac{2\pi n}{L},
\]

\[
n=-\frac{N}{2},\ldots,\frac{N}{2}-1.
\]

动量算符满足

\[
\hat p\phi_k(x)=\hbar k\phi_k(x).
\]

因此动能矩阵在动量空间对角：

\[
T_k=\frac{\hbar^2k^2}{2m}.
\]

势能在坐标网格上近似为

\[
V_{ij}=V(x_i)\delta_{ij}.
\]

这构成 Fourier Grid Hamiltonian 方法。Marston 和 Balint-Kurti 给出了 FGH
在分子束缚态计算中的系统表述【Marston \& Balint-Kurti 1989】。

动量分辨率和最大可表示动量为

\[
\Delta k=\frac{2\pi}{L},
\qquad
k_{\max}=\frac{\pi}{\Delta x}.
\]

若波函数动量宽度为 \(\sigma_k\)，为避免混叠通常要求

\[
k_{\max}\gtrsim 6\sigma_k.
\]

空间分辨率应满足

\[
\Delta x
\ll
\frac{1}{k_{\max}^{\mathrm{local}}}.
\]

在数值实验中，最直接的检验是把 \(N\)
增加一倍，观察本征值和期望值是否稳定。

\hypertarget{sinc-dvr-ux7684ux63a8ux5bfcux4e0eux8befux5deeux7ed3ux6784}{%
\subsection{2.4　Sinc-DVR
的推导与误差结构}\label{sinc-dvr-ux7684ux63a8ux5bfcux4e0eux8befux5deeux7ed3ux6784}}

对均匀网格

\[
x_i=x_{\min}+i\Delta x,
\qquad
i=1,\ldots,N,
\]

Colbert 和 Miller 给出的 Sinc-DVR 动能矩阵元为

\[
T_{ij}
=
\frac{\hbar^2}{2m\Delta x^2}
\begin{cases}
\pi^2/3, & i=j,\\[4pt]
2(-1)^{i-j}/(i-j)^2, & i\ne j.
\end{cases}
\]

原子单位下 \(\hbar=1\)，因此

\[
T_{ij}
=
\frac{1}{2m\Delta x^2}
\begin{cases}
\pi^2/3, & i=j,\\[4pt]
2(-1)^{i-j}/(i-j)^2, & i\ne j.
\end{cases}
\]

总哈密顿量为

\[
H_{ij}=T_{ij}+V(x_i)\delta_{ij}.
\]

若离散内积权重为 \(w_i=\Delta x\)，则连续归一化波函数为

\[
\psi(x_i)=\frac{c_i}{\sqrt{\Delta x}}.
\]

\hypertarget{ux8befux5deeux6765ux6e90}{%
\subsubsection{误差来源}\label{ux8befux5deeux6765ux6e90}}

Sinc-DVR 的误差主要来自四类：

\begin{enumerate}
\def\labelenumi{\arabic{enumi}.}
\tightlist
\item
  空间截断：计算域太窄导致边界反射；
\item
  网格欠采样：波函数高频分量超过 \(k_{\max}\)；
\item
  势能奇点：局部梯度过大；
\item
  非厄米离散：动能和势能使用了不一致的内积。
\end{enumerate}

对光滑势函数，谱方法通常呈指数收敛：

\[
E_N-E\sim C_Ne^{-\alpha N}.
\]

若势能不可微或存在奇点，收敛可能退化为代数阶：

\[
E_N-E\sim CN^{-p}.
\]

因此，不能仅凭两次计算结果接近判断收敛，应至少使用三个网格密度做
Richardson 外推。

\hypertarget{legendre-dvr-ux4e0eux89d2ux5411ux8868ux793a}{%
\subsection{2.5　Legendre-DVR
与角向表示}\label{legendre-dvr-ux4e0eux89d2ux5411ux8868ux793a}}

分子取向问题常采用

\[
x=\cos\theta.
\]

Legendre 多项式满足

\[
(1-x^2)P_n''(x)-2xP_n'(x)+n(n+1)P_n(x)=0,
\]

且

\[
\int_{-1}^{1}P_n(x)P_m(x)\,dx
=
\frac{2}{2n+1}\delta_{nm}.
\]

Gauss-Legendre 求积为

\[
\int_{-1}^{1}f(x)\,dx
\approx
\sum_{i=1}^{N}w_if(x_i).
\]

在角坐标中，

\[
\int_0^\pi f(\theta)\sin\theta\,d\theta
=
\int_{-1}^{1}f(\arccos x)\,dx
\approx
\sum_iw_if(\arccos x_i).
\]

转子动能算符为

\[
\hat J^2
=
-\hbar^2
\left[
\frac{1}{\sin\theta}
\frac{\partial}{\partial\theta}
\left(
\sin\theta
\frac{\partial}{\partial\theta}
\right)
\right].
\]

在球谐基 \(|lm\rangle\) 中，

\[
\hat J^2|lm\rangle
=
\hbar^2l(l+1)|lm\rangle.
\]

在 DVR 网格上，它变为矩阵 \(J^2_{ij}\)。GeneralModule 的
\texttt{dvr\_legendre\_init} 计算 Gauss-Legendre 零点、权重和 \(J^2\)
矩阵。

\hypertarget{ux5bf9ux5e94ux6e90ux7801}{%
\subsubsection{对应源码}\label{ux5bf9ux5e94ux6e90ux7801}}

\begin{itemize}
\tightlist
\item
  \href{https://github.com/l1Ha/QuantumGeneralModule/blob/main/src/mod_dvr_grid.f90}{src/mod\_dvr\_grid.f90}
\item
  二维码：\texttt{qr/src\_\_mod\_dvr\_grid.f90.png}
\end{itemize}

\includegraphics[width=2cm,height=\textheight]{qr/src__mod_dvr_grid.f90.png}

\hypertarget{dvr-ux7c7bux578bux7684ux9002ux7528ux8303ux56f4}{%
\subsection{2.6　DVR
类型的适用范围}\label{dvr-ux7c7bux578bux7684ux9002ux7528ux8303ux56f4}}

DVR 不是单一算法。不同 DVR 的差异在于权重、动能矩阵和边界条件。

\begin{longtable}[]{@{}lllll@{}}
\toprule
DVR 类型 & 坐标 & 动能矩阵 & 适用问题 & 主要风险\tabularnewline
\midrule
\endhead
Sinc-DVR & 等距 \(x\) & 解析长程衰减矩阵 & 一维局域势能面 &
边界反射\tabularnewline
Fourier/FGH & 周期 \(x\) & 动量空间对角 & 周期或大范围波包 &
周期性假设\tabularnewline
Gauss-Legendre DVR & \(x=\cos\theta\) & 角动量矩阵 & 分子取向与转子 &
端点处理\tabularnewline
Gauss-Hermite DVR & 无限域 & 谐振子结构 & 振动模式 &
势能范围超出权重衰减\tabularnewline
Laguerre-DVR & 半无限径向域 & 径向结构 & 长程势 &
奇点处理\tabularnewline
Distributed Gaussian DVR & 可调局域高斯 & 动能重叠矩阵 &
多维非均匀势能面 & 线性相关\tabularnewline
\bottomrule
\end{longtable}

因此，将 DVR 等同于 Sinc-DVR 是不完整的。GeneralModule 提供 Sinc-DVR 与
Gauss-Legendre DVR，并可通过同一框架扩展到其他正交多项式网格。

\hypertarget{ux591aux7ef4ux5f20ux91cfux79ef-dvr}{%
\subsection{2.7　多维张量积
DVR}\label{ux591aux7ef4ux5f20ux91cfux79ef-dvr}}

若体系有 \(d\) 个自由度，则张量积网格为

\[
(q_{1,i_1},q_{2,i_2},\ldots,q_{d,i_d}).
\]

总网格数为

\[
N_{\mathrm{tot}}=\prod_{\alpha=1}^{d}N_\alpha.
\]

若动能可分离，

\[
\hat H=\sum_{\alpha}\hat T_\alpha+\hat V(q_1,\ldots,q_d),
\]

则

\[
\hat T_1\rightarrow T_1\otimes I_2,
\qquad
\hat T_2\rightarrow I_1\otimes T_2.
\]

可分离势能项为

\[
V_1(q_1)
\rightarrow
\operatorname{diag}[V_1(q_{1,i})]\otimes I_2.
\]

耦合势能项在网格点集上仍是对角的，但以全维度索引展开。

\hypertarget{ux590dux6742ux5ea6}{%
\subsubsection{复杂度}\label{ux590dux6742ux5ea6}}

若每个维度有 \(N\) 个点，\(d\) 维张量积网格大小为
\(N^d\)。稠密对角化复杂度为

\[
O(N^{3d}),
\]

存储复杂度为

\[
O(N^{2d}).
\]

因此，直接对角化只适用于低维或截断后的小问题。高维问题需要稀疏矩阵、Krylov
方法、Chebyshev 传播、张量分解或缩减基方法。

\hypertarget{ux6536ux655bux6027ux4e0e-richardson-ux5916ux63a8}{%
\subsection{2.8　收敛性与 Richardson
外推}\label{ux6536ux655bux6027ux4e0e-richardson-ux5916ux63a8}}

设精确本征值为 \(E\)，数值本征值为 \(E_N\)。若误差服从

\[
E_N-E
=
CN^{-p}+O(N^{-(p+1)}),
\]

则三种网格 \(N_1,N_2,N_3\) 可估计阶数 \(p\)：

\[
p
=
\frac{
\ln\left[
\frac{E_{N_1}-E_{N_2}}
{E_{N_2}-E_{N_3}}
\right]
}
{\ln(N_2/N_1)}.
\]

若 \(p\)
不稳定，则说明尚未进入渐近区，或存在边界、奇点、混叠等系统误差。

\hypertarget{ux6570ux503cux5b9eux9a8c}{%
\subsubsection{数值实验}\label{ux6570ux503cux5b9eux9a8c}}

\begin{enumerate}
\def\labelenumi{\arabic{enumi}.}
\tightlist
\item
  对 Morse 势分别取 \(N=128,256,512,1024\)，绘制前六个能级误差。
\item
  固定 \(\Delta x\)，改变区间长度，观察高激发态边界反射。
\item
  将 Morse 势替换为谐振子势，比较 Sinc-DVR 和解析谱。
\item
  用 Legendre-DVR 计算刚性转子 \(J=0,1,2\) 能级，并与 \(BJ(J+1)\) 比较。
\end{enumerate}

\hypertarget{ux4eceux77e9ux9635ux7ed3ux6784ux5230ux7b97ux6cd5ux9009ux62e9}{%
\subsection{2.9　从矩阵结构到算法选择}\label{ux4eceux77e9ux9635ux7ed3ux6784ux5230ux7b97ux6cd5ux9009ux62e9}}

在 DVR 中，哈密顿量矩阵的结构决定算法：

\begin{longtable}[]{@{}lll@{}}
\toprule
矩阵结构 & 算法 & 复杂度\tabularnewline
\midrule
\endhead
中小型稠密对称矩阵 & Householder + QL & \(O(N^3)\)\tabularnewline
大型稀疏矩阵 & Lanczos / Arnoldi & 接近 \(O(kN)\)\tabularnewline
FFT 可分离动能 & Split-Operator & \(O(N\log N)\) 每步\tabularnewline
大步长谱传播 & Chebyshev 展开 & \(O(kN)\) 每步\tabularnewline
高维张量结构 & Tensor train / MCTDH & 问题相关\tabularnewline
\bottomrule
\end{longtable}

GeneralModule 的 \texttt{mod\_linear\_algebra}
提供自包含稠密对称本征求解和 FFT；\texttt{mod\_dvr\_grid}
提供网格构造；\texttt{mod\_wavepacket\_propagator} 和
\texttt{mod\_chebyshev\_propagator}
提供时间演化。三者构成一个完整的谱方法教学与验证链路。

\hypertarget{python-ux4ea4ux53c9ux9a8cux8bc1}{%
\subsection{2.10　Python
交叉验证}\label{python-ux4ea4ux53c9ux9a8cux8bc1}}

\texttt{pygenmod} 使用 NumPy 和 Matplotlib 对
DVR、散射和谱方法进行独立实现与可视化。这有两个作用：

\begin{enumerate}
\def\labelenumi{\arabic{enumi}.}
\tightlist
\item
  用另一种语言检验 Fortran 数值实现；
\item
  将 Fortran 输出转化为发表级图形。
\end{enumerate}

NumPy 的数组操作和 FFT 实现可作为参考算法，Matplotlib
用于绘制能级、波函数和误差曲线【Harris et al.~2020；Hunter
2007】。此类交叉验证不是替代 Fortran 主计算，而是数值验证工具。

\hypertarget{ux5bf9ux5e94ux6e90ux7801-1}{%
\subsubsection{对应源码}\label{ux5bf9ux5e94ux6e90ux7801-1}}

\begin{itemize}
\tightlist
\item
  \href{https://github.com/l1Ha/QuantumGeneralModule/blob/main/python/pygenmod/visualizer.py}{python/pygenmod/visualizer.py}
\item
  二维码：\texttt{qr/python\_\_pygenmod\_\_visualizer.py.png}
\end{itemize}

\includegraphics[width=2cm,height=\textheight]{qr/python__pygenmod__visualizer.py.png}

\hypertarget{ux53c2ux8003ux6587ux732e}{%
\subsection{2.11　参考文献}\label{ux53c2ux8003ux6587ux732e}}

\begin{enumerate}
\def\labelenumi{\arabic{enumi}.}
\tightlist
\item
  D. T. Colbert and W. H. Miller, ``A novel discrete variable
  representation for quantum mechanical reactive scattering via the
  S-matrix Kohn method'', \emph{J. Chem. Phys.} \textbf{96}, 1982
  (1992). DOI: 10.1063/1.462125.
\item
  C. C. Marston and G. G. Balint-Kurti, ``The Fourier grid Hamiltonian
  method for bound state eigenvalues and eigenfunctions'', \emph{J.
  Chem. Phys.} \textbf{91}, 3571 (1989). DOI: 10.1063/1.456888.
\item
  D. Baye and P.-H. Heenen, ``Generalised meshes for quantum mechanical
  problems'', \emph{J. Phys. B: At. Mol. Opt. Phys.} \textbf{19}, 1991
  (1986). DOI: 10.1088/0022-3700/19/14/006.
\item
  J. C. Light and T. Carrington, Jr., ``Discrete-variable
  representations and their utilization in quantum dynamics
  calculations'', \emph{Adv. Chem. Phys.} \textbf{114}, 263 (2000). DOI:
  10.1002/9780470141731.ch5.
\item
  C. R. Harris et al., ``Array programming with NumPy'', \emph{Nature}
  \textbf{585}, 357 (2020). DOI: 10.1038/s41586-020-2649-2.
\item
  J. D. Hunter, ``Matplotlib: A 2D graphics environment'', \emph{Comput.
  Sci. Eng.} \textbf{9}, 90 (2007). DOI: 10.1109/MCSE.2007.55.
\item
  GeneralModule
  源码与文档：\url{https://github.com/l1Ha/QuantumGeneralModule} \#
  第三部分　含时演化、开放系统与量子控制
\end{enumerate}

\hypertarget{ux65f6ux95f4ux6f14ux5316ux4e0eux5b88ux6052ux7ed3ux6784}{%
\subsection{3.1　时间演化与守恒结构}\label{ux65f6ux95f4ux6f14ux5316ux4e0eux5b88ux6052ux7ed3ux6784}}

含时薛定谔方程为

\[
i\hbar\frac{\partial}{\partial t}|\psi(t)\rangle
=
\hat H(t)|\psi(t)\rangle.
\]

若哈密顿量不含时，则形式解为

\[
|\psi(t)\rangle
=
\hat U(t,t_0)|\psi(t_0)\rangle,
\]

\[
\hat U(t,t_0)
=
\exp\left[-\frac{i}{\hbar}\hat H(t-t_0)\right].
\]

若哈密顿量含时，则

\[
\hat U(t,t_0)
=
\mathcal{T}
\exp\left[
-\frac{i}{\hbar}
\int_{t_0}^{t}\hat H(t')\,dt'
\right],
\]

其中 \(\mathcal{T}\)
是时间排序算符。数值方法的任务是用有限步近似该时间排序指数。

对不含时厄米哈密顿量，

\[
\hat U^{\dagger}\hat U=I,
\]

即

\[
\langle\psi(t)|\psi(t)\rangle
=
\langle\psi(t_0)|\psi(t_0)\rangle.
\]

数值传播器必须保持或至少诊断该性质。

\hypertarget{ux53efux89c2ux6d4bux91cf}{%
\subsubsection{可观测量}\label{ux53efux89c2ux6d4bux91cf}}

坐标期望值为

\[
\langle x\rangle(t)
=
\int x|\psi(x,t)|^2\,dx.
\]

动量期望值为

\[
\langle p\rangle(t)
=
\int \psi^{*}(x,t)
\left(-i\hbar\frac{\partial}{\partial x}\right)
\psi(x,t)\,dx.
\]

若使用动量空间波函数 \(\tilde\psi(k,t)\)，则

\[
\langle p\rangle(t)
=
\int \hbar k\,|\tilde\psi(k,t)|^2\,dk.
\]

\hypertarget{split-operator-ux65b9ux6cd5}{%
\subsection{3.2　Split-Operator
方法}\label{split-operator-ux65b9ux6cd5}}

对于不含时哈密顿量

\[
\hat H=\hat T+\hat V,
\]

若 \(\hat T\) 与 \(\hat V\) 不对易，则

\[
e^{-i(\hat T+\hat V)\Delta t/\hbar}
\ne
e^{-i\hat T\Delta t/\hbar}
e^{-i\hat V\Delta t/\hbar}.
\]

利用 Baker--Campbell--Hausdorff 展开可得一阶分裂

\[
e^{-i(\hat T+\hat V)\Delta t/\hbar}
=
e^{-i\hat T\Delta t/\hbar}
e^{-i\hat V\Delta t/\hbar}
+
O(\Delta t^2).
\]

二阶对称 Strang 分裂为

\[
\hat U(\Delta t)
=
e^{-i\hat V\Delta t/(2\hbar)}
e^{-i\hat T\Delta t/\hbar}
e^{-i\hat V\Delta t/(2\hbar)}
+
O(\Delta t^3).
\]

Feit、Fleck 和 Steiger 将该谱方法系统用于薛定谔方程求解【Feit, Fleck \&
Steiger 1982】。Kosloff
对含时量子动力学中的谱传播方法作了系统综述【Kosloff 1988】。

在坐标表象中，\(\hat V\) 是局域算符：

\[
[\hat V\psi](x)=V(x)\psi(x).
\]

在动量表象中，\(\hat T=p^2/(2m)\) 是对角算符：

\[
[\hat T\tilde\psi](k)=\frac{\hbar^2 k^2}{2m}\tilde\psi(k).
\]

因此，一步传播为

\[
\psi(x,t+\Delta t)
=
e^{-iV(x)\Delta t/(2\hbar)}
\mathcal{F}^{-1}
\left[
e^{-i\hbar k^2\Delta t/(2m)}
\mathcal{F}
\left[
e^{-iV(x)\Delta t/(2\hbar)}
\psi(x,t)
\right]
\right].
\]

在原子单位下，\(\hbar=1\)，公式简化为

\[
\psi(x,t+\Delta t)
=
e^{-iV(x)\Delta t/2}
\mathcal{F}^{-1}
\left[
e^{-ik^2\Delta t/(2m)}
\mathcal{F}
\left[
e^{-iV(x)\Delta t/2}
\psi(x,t)
\right]
\right].
\]

\hypertarget{ux6570ux503cux6b65ux9aa4}{%
\subsubsection{数值步骤}\label{ux6570ux503cux6b65ux9aa4}}

\begin{enumerate}
\def\labelenumi{\arabic{enumi}.}
\tightlist
\item
  在坐标网格上计算势能半步相位 \(e^{-iV(x_i)\Delta t/2}\)；
\item
  将波函数乘以该相位；
\item
  使用 FFT 变换到动量空间；
\item
  乘以动能相位 \(e^{-ik^2\Delta t/(2m)}\)；
\item
  使用逆 FFT 回到坐标空间；
\item
  再次乘以势能半步相位。
\end{enumerate}

每步需要两次 FFT，复杂度为 \(O(N\log N)\)，显著优于稠密矩阵指数的
\(O(N^3)\)。

\hypertarget{ux8303ux6570ux4e0eux7a33ux5b9aux6027}{%
\subsubsection{范数与稳定性}\label{ux8303ux6570ux4e0eux7a33ux5b9aux6027}}

每个指数相位因子都是幺正算符。因此，在无穷精度下 Split-Operator
严格保持范数：

\[
\|\psi(t+\Delta t)\|_2=\|\psi(t)\|_2.
\]

有限精度误差来源包括：

\begin{itemize}
\tightlist
\item
  FFT 归一化不一致；
\item
  网格混叠；
\item
  复势导致非幺正演化；
\item
  吸收边界；
\item
  长时间积累的舍入误差；
\item
  不当归一化掩盖真实误差。
\end{itemize}

因此，应监测

\[
\epsilon_{\mathrm{norm}}(t)
=
\left|
\|\psi(t)\|_2-1
\right|,
\]

而不是在每步强制归一化。

\hypertarget{ux5bf9ux5e94ux6e90ux7801-2}{%
\subsubsection{对应源码}\label{ux5bf9ux5e94ux6e90ux7801-2}}

\begin{itemize}
\tightlist
\item
  \href{https://github.com/l1Ha/QuantumGeneralModule/blob/main/src/mod_wavepacket_propagator.f90}{src/mod\_wavepacket\_propagator.f90}
\item
  二维码：\texttt{qr/src\_\_mod\_wavepacket\_propagator.f90.png}
\end{itemize}

\includegraphics[width=2cm,height=\textheight]{qr/src__mod_wavepacket_propagator.f90.png}

\hypertarget{fftux52a8ux91cfux7f51ux683cux4e0eux6df7ux53e0}{%
\subsection{3.3　FFT、动量网格与混叠}\label{fftux52a8ux91cfux7f51ux683cux4e0eux6df7ux53e0}}

均匀网格 \(x_i=x_0+i\Delta x\)，\(i=0,\ldots,N-1\) 对应动量网格

\[
k_j
=
\frac{2\pi}{L}
\begin{cases}
j, & 0\le j<N/2,\\
j-N, & N/2\le j<N,
\end{cases}
\]

其中 \(L=N\Delta x\)。这样，\(k\) 覆盖

\[
-\frac{\pi}{\Delta x}
\le
k
<
\frac{\pi}{\Delta x}.
\]

离散傅里叶变换为

\[
\tilde\psi(k_j)
=
\sum_{n=0}^{N-1}
\psi(x_n)e^{-ik_jx_n}\Delta x,
\]

逆变换为

\[
\psi(x_n)
=
\frac{1}{L}
\sum_{j=0}^{N-1}
\tilde\psi(k_j)e^{ik_jx_n}.
\]

若波函数在动量空间的有效宽度超过 \(|k|<\pi/\Delta x\)，则发生混叠：

\[
k_{\mathrm{alias}}
=
k-2m\pi/\Delta x,
\qquad
m\in\mathbb{Z}.
\]

避免混叠的条件是

\[
\Delta x
<
\frac{\pi}{k_{\max}}.
\]

在数值实验中应通过逐步减小 \(\Delta x\) 检查可观测量是否收敛。

\hypertarget{rungekutta-ux4e0eux591aux6b65ux65b9ux6cd5}{%
\subsection{3.4　Runge--Kutta
与多步方法}\label{rungekutta-ux4e0eux591aux6b65ux65b9ux6cd5}}

将波函数或密度矩阵写成常微分方程

\[
\frac{d\mathbf{y}}{dt}=f(t,\mathbf{y}),
\]

经典四阶 Runge--Kutta 方法为

\[
\begin{aligned}
k_1&=f(t_n,y_n),\\
k_2&=f\left(t_n+\frac{\Delta t}{2},y_n+\frac{\Delta t}{2}k_1\right),\\
k_3&=f\left(t_n+\frac{\Delta t}{2},y_n+\frac{\Delta t}{2}k_2\right),\\
k_4&=f(t_n+\Delta t,y_n+\Delta t k_3),\\
y_{n+1}
&=
y_n+
\frac{\Delta t}{6}
(k_1+2k_2+2k_3+k_4).
\end{aligned}
\]

其局部截断误差为 \(O(\Delta t^5)\)，全局误差为 \(O(\Delta t^4)\)。RK4
对非幺正演化或非线性常微分方程非常通用，但对长时幺正演化可能产生范数漂移。

Adams--Bashforth--Moulton
四阶预估---校正方法使用历史导数。Adams--Bashforth 预估式为

\[
y_{n+1}^{(p)}
=
y_n
+
\frac{\Delta t}{24}
(55f_n-59f_{n-1}+37f_{n-2}-9f_{n-3}).
\]

在预估点计算

\[
f_{n+1}^{(p)}
=
f(t_{n+1},y_{n+1}^{(p)}),
\]

然后使用 Adams--Moulton 校正式

\[
y_{n+1}
=
y_n
+
\frac{\Delta t}{24}
(9f_{n+1}^{(p)}+19f_n-5f_{n-1}+f_{n-2}).
\]

多步方法依赖历史点的一致性和启动策略。若历史导数来自不同时间步长或不同坐标系，阶数会退化。

\hypertarget{ux5bf9ux5e94ux6e90ux7801-3}{%
\subsubsection{对应源码}\label{ux5bf9ux5e94ux6e90ux7801-3}}

\begin{itemize}
\tightlist
\item
  \href{https://github.com/l1Ha/QuantumGeneralModule/blob/main/src/mod_wavepacket_propagator.f90}{src/mod\_wavepacket\_propagator.f90}
\item
  二维码：\texttt{qr/src\_\_mod\_wavepacket\_propagator.f90.png}
\end{itemize}

\includegraphics[width=2cm,height=\textheight]{qr/src__mod_wavepacket_propagator.f90.png}

\hypertarget{ux7ec3ux4e60}{%
\subsubsection{练习}\label{ux7ec3ux4e60}}

\begin{enumerate}
\def\labelenumi{\arabic{enumi}.}
\tightlist
\item
  将谐振子波包分别用 Split-Operator 和 RK4 传播，比较范数误差随
  \(\Delta t\) 的标度。
\item
  构造非光滑右端项，比较 RK4 和 ABM4 的稳定性。
\item
  改变 FFT 网格密度，观察动量空间尾部截断对
  \(\langle x\rangle\)、\(\langle p\rangle\) 的影响。
\end{enumerate}

\hypertarget{ux4e8cux80fdux7ea7ux7cfbux7edfux4e0e-bloch-ux65b9ux7a0b}{%
\subsection{3.5　二能级系统与 Bloch
方程}\label{ux4e8cux80fdux7ea7ux7cfbux7edfux4e0e-bloch-ux65b9ux7a0b}}

二能级系统的态可写为

\[
|\psi(t)\rangle
=
c_1(t)|1\rangle+c_2(t)|2\rangle.
\]

在旋转波近似下，拉比频率为 \(\Omega_R(t)\)，失谐为 \(\Delta(t)\)，Bloch
矢量为

\[
\mathbf{R}
=
(u,v,w)^T.
\]

光学 Bloch 方程的无耗散部分为

\[
\frac{d\mathbf{R}}{dt}
=
\boldsymbol{\Omega}\times\mathbf{R},
\]

其中

\[
\boldsymbol{\Omega}
=
(\Omega_R,0,\Delta)^T.
\]

展开为

\[
\dot u=-\Delta v,
\]

\[
\dot v=\Delta u-\Omega_R w,
\]

\[
\dot w=\Omega_R v.
\]

\(w=\rho_{22}-\rho_{11}\) 表示布居反转。若 \(\Delta=0\) 且脉冲面积为

\[
\Theta
=
\int \Omega_R(t)\,dt,
\]

则完全反转为

\[
\Theta=\pi.
\]

该方程是检验 RK4、保范数传播和激光脉冲合成的基础模型。

\hypertarget{ux5bf9ux5e94ux6e90ux7801-4}{%
\subsubsection{对应源码}\label{ux5bf9ux5e94ux6e90ux7801-4}}

\begin{itemize}
\tightlist
\item
  \href{https://github.com/l1Ha/QuantumGeneralModule/blob/main/src/mod_wavepacket_propagator.f90}{src/mod\_wavepacket\_propagator.f90}
\item
  二维码：\texttt{qr/src\_\_mod\_wavepacket\_propagator.f90.png}
\end{itemize}

\includegraphics[width=2cm,height=\textheight]{qr/src__mod_wavepacket_propagator.f90.png}

\hypertarget{ux590dux5438ux6536ux8fb9ux754c}{%
\subsection{3.6　复吸收边界}\label{ux590dux5438ux6536ux8fb9ux754c}}

当波包离开有限计算域时，硬边界会产生非物理反射。复吸收势（Complex
Absorbing Potential, CAP）在边界区域加入虚部势：

\[
V_{\mathrm{eff}}(x)
=
V(x)-iW(x),
\]

\[
W(x)\ge0.
\]

常用 \(\sin^2\) 形式为

\[
W(x)
=
\eta
\sin^2
\left[
\frac{\pi}{2}
\frac{x-x_a}{L_a}
\right],
\qquad
x_a\le x\le x_a+L_a.
\]

加入 CAP 后，哈密顿量变为非厄米：

\[
\hat H_{\mathrm{eff}}
=
\hat T+\hat V-i\hat W.
\]

内部区域波函数的模平方不再守恒，其损失应与流出概率一致：

\[
\frac{d}{dt}
\int_{\Omega_{\mathrm{in}}}
|\psi|^2\,dx
=
-
\int_{\partial\Omega_{\mathrm{in}}}
J_n\,dS
-
\int_{\Omega_{\mathrm{CAP}}}
2W(x)|\psi|^2\,dx.
\]

概率流为

\[
J(x,t)
=
\frac{\hbar}{m}
\operatorname{Im}
\left[
\psi^{*}(x,t)
\frac{\partial\psi(x,t)}{\partial x}
\right].
\]

好的 CAP 参数应满足：

\begin{enumerate}
\def\labelenumi{\arabic{enumi}.}
\tightlist
\item
  内部物理区域几乎不受 CAP 影响；
\item
  反射振幅足够小；
\item
  吸收区长度和强度不过度增大计算域；
\item
  吸收量与通量诊断一致。
\end{enumerate}

\hypertarget{ux5bf9ux5e94ux6e90ux7801-5}{%
\subsubsection{对应源码}\label{ux5bf9ux5e94ux6e90ux7801-5}}

\begin{itemize}
\tightlist
\item
  \href{https://github.com/l1Ha/QuantumGeneralModule/blob/main/src/mod_absorbing_boundary.f90}{src/mod\_absorbing\_boundary.f90}
\item
  二维码：\texttt{qr/src\_\_mod\_absorbing\_boundary.f90.png}
\end{itemize}

\includegraphics[width=2cm,height=\textheight]{qr/src__mod_absorbing_boundary.f90.png}

\hypertarget{chebyshev-ux4f20ux64adux4e0eux8c31ux65b9ux6cd5ux5927ux6b65ux957f}{%
\subsection{3.7　Chebyshev
传播与谱方法大步长}\label{chebyshev-ux4f20ux64adux4e0eux8c31ux65b9ux6cd5ux5927ux6b65ux957f}}

对不含时哈密顿量，可利用 Chebyshev 多项式展开演化算符：

\[
e^{-i\hat H t/\hbar}
=
\sum_{n=0}^{\infty}
a_n(t)T_n(\tilde{\hat H}),
\]

其中 \(\tilde{\hat H}\) 是缩放到 \([-1,1]\) 的哈密顿量。若哈密顿量谱满足

\[
E_{\min}\le \hat H\le E_{\max},
\]

则令

\[
\tilde{\hat H}
=
\frac{2\hat H-(E_{\max}+E_{\min})I}
{E_{\max}-E_{\min}}.
\]

Chebyshev 展开系数为

\[
a_n(t)
=
(2-\delta_{n0})(-i)^n
J_n(\Delta E\,t/2\hbar),
\]

其中 \(J_n\) 是第一类 Bessel
函数，\(\Delta E=E_{\max}-E_{\min}\)。所需项数大约为

\[
N_{\mathrm{Cheb}}
\sim
\frac{\Delta E\,t}{2\hbar}
+
c\ln\epsilon^{-1}.
\]

Chebyshev
方法适合大时间步和长时传播，但要求谱界估计准确。若谱外态存在，可能出现非物理放大。

\hypertarget{ux5bf9ux5e94ux6e90ux7801-6}{%
\subsubsection{对应源码}\label{ux5bf9ux5e94ux6e90ux7801-6}}

\begin{itemize}
\tightlist
\item
  \href{https://github.com/l1Ha/QuantumGeneralModule/blob/main/src/mod_chebyshev_propagator.f90}{src/mod\_chebyshev\_propagator.f90}
\item
  二维码：\texttt{qr/src\_\_mod\_chebyshev\_propagator.f90.png}
\end{itemize}

\includegraphics[width=2cm,height=\textheight]{qr/src__mod_chebyshev_propagator.f90.png}

\hypertarget{ux5bc6ux5ea6ux77e9ux9635ux4e0eux5f00ux653eux91cfux5b50ux7cfbux7edf}{%
\subsection{3.8　密度矩阵与开放量子系统}\label{ux5bc6ux5ea6ux77e9ux9635ux4e0eux5f00ux653eux91cfux5b50ux7cfbux7edf}}

封闭系统的态可用纯态 \(|\psi\rangle\) 描述，对应密度矩阵

\[
\hat\rho=|\psi\rangle\langle\psi|.
\]

一般混合态为

\[
\hat\rho
=
\sum_\alpha p_\alpha
|\psi_\alpha\rangle\langle\psi_\alpha|,
\qquad
p_\alpha\ge0,
\qquad
\sum_\alpha p_\alpha=1.
\]

密度矩阵满足

\[
\hat\rho=\hat\rho^{\dagger},
\qquad
\operatorname{Tr}\hat\rho=1,
\qquad
\hat\rho\ge0.
\]

封闭系统的刘维尔方程为

\[
\frac{d\hat\rho}{dt}
=
-\frac{i}{\hbar}[\hat H,\hat\rho].
\]

与环境耦合后，马尔可夫近似下的 Lindblad 主方程为

\[
\frac{d\hat\rho}{dt}
=
-\frac{i}{\hbar}[\hat H,\hat\rho]
+
\sum_k\gamma_k
\left(
\hat L_k\hat\rho\hat L_k^{\dagger}
-
\frac12
\{\hat L_k^{\dagger}\hat L_k,\hat\rho\}
\right).
\]

该形式由 Gorini--Kossakowski--Sudarshan 和 Lindblad 给出【Gorini,
Kossakowski \& Sudarshan 1976；Lindblad 1976】。Breuer 和 Petruccione
对开放量子系统理论作了系统论述【Breuer \& Petruccione 2002】。

\hypertarget{ux5b8cux5168ux6b63ux5b9aux6027ux4e0eux8ff9ux5b88ux6052}{%
\subsubsection{完全正定性与迹守恒}\label{ux5b8cux5168ux6b63ux5b9aux6027ux4e0eux8ff9ux5b88ux6052}}

Lindblad 形式保证

\[
\frac{d}{dt}\operatorname{Tr}\hat\rho=0,
\]

且演化映射完全正定。数值实现必须检查：

\begin{enumerate}
\def\labelenumi{\arabic{enumi}.}
\tightlist
\item
  迹守恒： \(|\operatorname{Tr}\hat\rho-1|<\epsilon\)；
\item
  厄米性： \(\|\hat\rho-\hat\rho^{\dagger}\|<\epsilon\)；
\item
  正定性： 最小本征值不显著小于 \(-\epsilon\)；
\item
  纯度单调性或物理耗散趋势。
\end{enumerate}

\hypertarget{ux5e38ux89c1ux8017ux6563ux901aux9053}{%
\subsubsection{常见耗散通道}\label{ux5e38ux89c1ux8017ux6563ux901aux9053}}

自发辐射跃迁算符为

\[
\hat L_{\mathrm{rel}}=\sqrt{\gamma}|g\rangle\langle e|.
\]

纯退相位算符为

\[
\hat L_{\phi}=\sqrt{\gamma_\phi}|e\rangle\langle e|.
\]

二能级弛豫满足

\[
\frac{d\rho_{ee}}{dt}
=
-\gamma\rho_{ee},
\]

\[
\frac{d\rho_{eg}}{dt}
=
-\left(i\omega_0+\frac{\gamma}{2}+\gamma_\phi\right)\rho_{eg}.
\]

因此，布居弛豫时间为 \(1/\gamma\)，相干衰减时间为

\[
T_2^{-1}
=
\frac{1}{2T_1}+\gamma_\phi.
\]

\hypertarget{ux7eafux5ea6ux76f8ux5e72ux5ea6ux4e0eux71b5}{%
\subsubsection{纯度、相干度与熵}\label{ux7eafux5ea6ux76f8ux5e72ux5ea6ux4e0eux71b5}}

纯度为

\[
P=\operatorname{Tr}(\hat\rho^2).
\]

纯态满足 \(P=1\)；最大混合态满足 \(P=1/d\)。冯·诺依曼熵为

\[
S=-\operatorname{Tr}(\hat\rho\ln\hat\rho).
\]

\(l_1\) 范数相干度为

\[
C_{l1}(\hat\rho)
=
\sum_{i\ne j}|\rho_{ij}|.
\]

这三个量分别描述混合程度、非对角相干和信息含量。

\hypertarget{rk4-ux6570ux503cux63a8ux8fdb}{%
\subsubsection{RK4 数值推进}\label{rk4-ux6570ux503cux63a8ux8fdb}}

将 Lindblad 方程写成矩阵值常微分方程

\[
\dot{\hat\rho}=F(\hat\rho),
\]

则 RK4 步为

\[
\hat\rho_{n+1}
=
\hat\rho_n
+
\frac{\Delta t}{6}
(K_1+2K_2+2K_3+K_4).
\]

每一步之后可检查厄米性和迹，但不应在参数研究开始前就强制归一化，否则会掩盖真实时间步长误差。

\hypertarget{ux5bf9ux5e94ux6e90ux7801-7}{%
\subsubsection{对应源码}\label{ux5bf9ux5e94ux6e90ux7801-7}}

\begin{itemize}
\tightlist
\item
  \href{https://github.com/l1Ha/QuantumGeneralModule/blob/main/src/mod_open_quantum.f90}{src/mod\_open\_quantum.f90}
\item
  二维码：\texttt{qr/src\_\_mod\_open\_quantum.f90.png}
\end{itemize}

\includegraphics[width=2cm,height=\textheight]{qr/src__mod_open_quantum.f90.png}

\hypertarget{ux91cfux5b50ux6700ux4f18ux63a7ux5236ux4e0e-krotov-ux65b9ux6cd5}{%
\subsection{3.9　量子最优控制与 Krotov
方法}\label{ux91cfux5b50ux6700ux4f18ux63a7ux5236ux4e0e-krotov-ux65b9ux6cd5}}

量子控制的目标是设计外场
\(\epsilon(t)\)，使演化态在终端时刻逼近目标态。保真度为

\[
F
=
|\langle\phi_{\mathrm{target}}|\psi(T)\rangle|^2.
\]

最优控制问题可写成泛函

\[
J
=
F
-
\int_0^T
\lambda(t)\epsilon^2(t)\,dt.
\]

Krotov 方法引入伴随态 \(|\chi(t)\rangle\)，满足反向方程

\[
i\hbar\frac{\partial}{\partial t}|\chi(t)\rangle
=
\hat H(t)|\chi(t)\rangle,
\]

终端条件为

\[
|\chi(T)\rangle
=
|\phi_{\mathrm{target}}\rangle.
\]

场的更新形式可写为

\[
\epsilon^{(k+1)}(t)
=
\epsilon^{(k)}(t)
+
\frac{S(t)}{\alpha}
\operatorname{Im}
\left[
\langle\chi^{(k)}(t)|\hat\mu|\psi^{(k+1)}(t)\rangle
\right],
\]

其中 \(S(t)\) 是脉冲开关函数，\(\alpha\) 是场强惩罚参数。Somlói、Kazakov
和 Tannor 给出了分子运动的广义弛豫控制方法【Somlói, Kazakov \& Tannor
1993】；Reich、Ndong 和 Koch 讨论了单调收敛条件【Reich, Ndong \& Koch
2012】。

\hypertarget{krotov-ux5355ux6b21ux8fedux4ee3ux6d41ux7a0b}{%
\subsubsection{Krotov
单次迭代流程}\label{krotov-ux5355ux6b21ux8fedux4ee3ux6d41ux7a0b}}

\begin{enumerate}
\def\labelenumi{\arabic{enumi}.}
\tightlist
\item
  以当前场 \(\epsilon^{(k)}(t)\) 正向传播 \(|\psi(t)\rangle\)；
\item
  设置伴随态终端条件 \(|\chi(T)\rangle\)；
\item
  反向传播 \(|\chi(t)\rangle\)；
\item
  在每一时刻计算耦合项
  \(\operatorname{Im}\langle\chi|\hat\mu|\psi\rangle\)；
\item
  更新 \(\epsilon(t)\)；
\item
  用新场开始下一轮迭代。
\end{enumerate}

收敛判据应同时包括保真度增量、场变化和物理约束。

\hypertarget{ux5bf9ux5e94ux6e90ux7801-8}{%
\subsubsection{对应源码}\label{ux5bf9ux5e94ux6e90ux7801-8}}

\begin{itemize}
\tightlist
\item
  \href{https://github.com/l1Ha/QuantumGeneralModule/blob/main/src/mod_optimal_control.f90}{src/mod\_optimal\_control.f90}
\item
  二维码：\texttt{qr/src\_\_mod\_optimal\_control.f90.png}
\end{itemize}

\includegraphics[width=2cm,height=\textheight]{qr/src__mod_optimal_control.f90.png}

\hypertarget{ux9a8cux8bc1ux4e0eux6536ux655bux5b9eux9a8c}{%
\subsection{3.10　验证与收敛实验}\label{ux9a8cux8bc1ux4e0eux6536ux655bux5b9eux9a8c}}

含时量子动力学的验证应包含以下层次：

\begin{longtable}[]{@{}lll@{}}
\toprule
验证层次 & 检查量 & 目标\tabularnewline
\midrule
\endhead
解析极限 & 自由粒子、谐振子、\(\pi\) 脉冲 &
算法和单位正确\tabularnewline
范数守恒 & \(\|\psi(t)\|\) & 幺正性与 FFT 一致\tabularnewline
可观测守恒 & \(\langle H\rangle\)、\(\langle J\rangle\) &
对称性未被破坏\tabularnewline
收敛性 & \(E\)、\(T(E)\)、\(F\) 随 \(N,\Delta t\) 变化 &
误差可控\tabularnewline
CAP 一致性 & 内部范数损失与流出通量 &
吸收边界不污染物理区\tabularnewline
开放系统 & 迹、厄米性、正定性 & Lindblad 结构保持\tabularnewline
控制收敛 & \(F\)、\(\epsilon(t)\) 变化 & 最优控制收敛\tabularnewline
\bottomrule
\end{longtable}

一个基本数值实验是自由粒子波包。设初始高斯波包为

\[
\psi(x,0)
=
\left(\frac{1}{2\pi\sigma^2}\right)^{1/4}
\exp\left[
-\frac{(x-x_0)^2}{4\sigma^2}
+ik_0x
\right].
\]

无外场时，

\[
\langle x\rangle(t)
=
x_0+\frac{\hbar k_0}{m}t.
\]

数值传播应再现该线性运动和动量守恒。

\hypertarget{ux5bf9ux5e94ux6d4bux8bd5}{%
\subsubsection{对应测试}\label{ux5bf9ux5e94ux6d4bux8bd5}}

\begin{itemize}
\tightlist
\item
  \href{https://github.com/l1Ha/QuantumGeneralModule/blob/main/tests/test_propagators.f90}{tests/test\_propagators.f90}
\item
  \href{https://github.com/l1Ha/QuantumGeneralModule/blob/main/tests/test_open_quantum_opt.f90}{tests/test\_open\_quantum\_opt.f90}
\item
  二维码：\texttt{qr/tests\_\_test\_propagators.f90.png}
\end{itemize}

\includegraphics[width=2cm,height=\textheight]{qr/tests__test_propagators.f90.png}

\hypertarget{ux53c2ux8003ux6587ux732e-1}{%
\subsection{3.11　参考文献}\label{ux53c2ux8003ux6587ux732e-1}}

\begin{enumerate}
\def\labelenumi{\arabic{enumi}.}
\tightlist
\item
  M. D. Feit, J. A. Fleck, Jr., and A. Steiger, ``Solution of the
  Schrödinger equation by a spectral method'', \emph{J. Comput. Phys.}
  \textbf{47}, 412 (1982). DOI: 10.1016/0021-9991(82)90091-2.
\item
  R. Kosloff, ``Time-dependent quantum-mechanical methods for molecular
  dynamics'', \emph{J. Phys. Chem.} \textbf{92}, 2087 (1988). DOI:
  10.1021/j100319a003.
\item
  V. Gorini, A. Kossakowski, and E. C. G. Sudarshan, ``Completely
  positive dynamical semigroups of N-level systems'', \emph{J. Math.
  Phys.} \textbf{17}, 821 (1976). DOI: 10.1063/1.522979.
\item
  G. Lindblad, ``On the generators of quantum dynamical semigroups'',
  \emph{Commun. Math. Phys.} \textbf{48}, 119 (1976). DOI:
  10.1007/BF01608499.
\item
  H.-P. Breuer and F. Petruccione, \emph{The Theory of Open Quantum
  Systems}, Oxford University Press, Oxford, 2002.
\item
  R. Somlói, J. Kazakov, and D. J. Tannor, ``A generalized relaxation
  method for optimal control of molecular motion'', \emph{Chem. Phys.}
  \textbf{172}, 85 (1993). DOI: 10.1016/0301-0104(93)80108-L.
\item
  D. M. Reich, M. Ndong, and C. P. Koch, ``Monotonically convergent
  optimal control theory of quantum systems'', \emph{J. Chem. Phys.}
  \textbf{136}, 104103 (2012). DOI: 10.1063/1.3691827.
\item
  D. J. Tannor, \emph{Introduction to Quantum Mechanics: A
  Time-Dependent Perspective}, University Science Books, Sausalito,
  2007.
\item
  GeneralModule
  源码与文档：\url{https://github.com/l1Ha/QuantumGeneralModule} \#
  第四部分　散射、分子动力学与复杂体系
\end{enumerate}

\hypertarget{ux5b9aux6001ux6563ux5c04ux7684ux57faux672cux95eeux9898}{%
\subsection{4.1　定态散射的基本问题}\label{ux5b9aux6001ux6563ux5c04ux7684ux57faux672cux95eeux9898}}

对中心势 \(V(r)\)，径向波函数 \(u_l(r)\) 满足

\[
-\frac{\hbar^2}{2\mu}
\frac{d^2u_l}{dr^2}
+
\left[
V(r)
+
\frac{\hbar^2l(l+1)}{2\mu r^2}
\right]u_l(r)
=
Eu_l(r).
\]

在原子单位下，

\[
-\frac{1}{2\mu}u_l''(r)
+
\left[
V(r)
+
\frac{l(l+1)}{2\mu r^2}
\right]u_l(r)
=
Eu_l(r).
\]

在远区 \(V(r)\rightarrow0\) 时，

\[
u_l(r)
\rightarrow
\frac{1}{k}
\sin\left(kr-\frac{l\pi}{2}+\delta_l\right).
\]

散射矩阵元为

\[
S_l=e^{2i\delta_l}.
\]

反应矩阵为

\[
K_l=\tan\delta_l.
\]

跃迁矩阵为

\[
T_l=S_l-I.
\]

分波弹性截面为

\[
\sigma_l
=
\frac{4\pi}{k^2}(2l+1)\sin^2\delta_l.
\]

总弹性截面为

\[
\sigma_{\mathrm{el}}
=
\sum_l\sigma_l.
\]

\hypertarget{numerov-ux65b9ux6cd5ux4e0eux6563ux5c04ux957fux5ea6}{%
\subsection{4.2　Numerov
方法与散射长度}\label{numerov-ux65b9ux6cd5ux4e0eux6563ux5c04ux957fux5ea6}}

Numerov 方法适用于二阶常微分方程

\[
u''(r)=f(r)u(r).
\]

其三点递推为

\[
\left(1-\frac{\Delta r^2}{12}f_{i+1}\right)u_{i+1}
=
2
\left(1+\frac{5\Delta r^2}{12}f_i\right)u_i
-
\left(1-\frac{\Delta r^2}{12}f_{i-1}\right)u_{i-1}.
\]

对零能散射问题，

\[
-\frac{\hbar^2}{2\mu}u''(r)+V(r)u(r)=0.
\]

若势能在远区趋于零，则渐近形式为

\[
u(r)\rightarrow C(r-a_s).
\]

因此

\[
a_s=r-\frac{u(r)}{u'(r)}.
\]

GeneralModule 的 \texttt{calc\_scattering\_length\_numerov}
实现该流程，并在远区用三点公式计算对数导数：

\[
u'(r_N)
\approx
\frac{3u_N-4u_{N-1}+u_{N-2}}{2\Delta r}.
\]

\hypertarget{ux5bf9ux5e94ux6e90ux7801-9}{%
\subsubsection{对应源码}\label{ux5bf9ux5e94ux6e90ux7801-9}}

\begin{itemize}
\tightlist
\item
  \href{https://github.com/l1Ha/QuantumGeneralModule/blob/main/src/mod_ti_scattering.f90}{src/mod\_ti\_scattering.f90}
\item
  二维码：\texttt{qr/src\_\_mod\_ti\_scattering.f90.png}
\end{itemize}

\includegraphics[width=2cm,height=\textheight]{qr/src__mod_ti_scattering.f90.png}

\hypertarget{johnson-ux5bf9ux6570ux5bfcux6570ux65b9ux6cd5}{%
\subsection{4.3　Johnson
对数导数方法}\label{johnson-ux5bf9ux6570ux5bfcux6570ux65b9ux6cd5}}

直接传播 \(u(r)\) 在深势阱中容易指数溢出。Johnson 方法传播相邻点比值

\[
R_i=\frac{u_i}{u_{i-1}}.
\]

Numerov 递推可改写为

\[
R_{i+1}
=
\frac{
2(1-5w_i)-(1+w_{i-1})/R_i
}
{1+w_{i+1}},
\]

其中

\[
w_i=\frac{\Delta r^2}{12}f_i.
\]

渐近截距可由

\[
a_s
=
\frac{R_Nr_{N-1}-r_N}{R_N-1}
\]

得到。Johnson 方法在深势阱和长程势中数值稳定性更好。

\hypertarget{ux76f8ux79fbux4e0eux5149ux5b66ux5b9aux7406}{%
\subsection{4.4　相移与光学定理}\label{ux76f8ux79fbux4e0eux5149ux5b66ux5b9aux7406}}

散射波函数可展开为

\[
\psi(\mathbf r)
=
\sum_{lm}
\frac{u_l(r)}{r}
Y_{lm}(\hat{\mathbf r}).
\]

在远区，

\[
u_l(r)
\rightarrow
\frac{1}{k}
\sin\left(kr-\frac{l\pi}{2}+\delta_l\right).
\]

散射矩阵元为

\[
S_l=e^{2i\delta_l}.
\]

总截面与向前散射振幅满足光学定理

\[
\sigma_{\mathrm{tot}}
=
\frac{4\pi}{k}
\operatorname{Im}f(0).
\]

该关系可用于检验散射振幅、相移和截面计算的一致性。

\hypertarget{ux6709ux6548ux529bux7a0bux5c55ux5f00}{%
\subsection{4.5　有效力程展开}\label{ux6709ux6548ux529bux7a0bux5c55ux5f00}}

低能散射的有效力程展开为

\[
k\cot\delta_0(k)
=
-\frac{1}{a_s}
+
\frac12r_0k^2
+
O(k^4).
\]

零能截面为

\[
\sigma(0)=4\pi a_s^2.
\]

对范德华长程势，Gribakin--Flambaum 半经典平均散射长度为

\[
\bar a
=
\frac{2}{\sqrt\pi}
\frac{
\Gamma(5/6)
}{
\Gamma(2/3)
}
\left(
\frac{2\mu C_6}{\hbar^2}
\right)^{1/4}.
\]

GeneralModule 的 \texttt{fit\_effective\_range\_expansion} 和
\texttt{gribakin\_flambaum\_length} 分别实现这两类低能分析。

\hypertarget{ux591aux901aux9053ux8026ux5408ux4e0e-johnson-ux77e9ux9635ux5bf9ux6570ux5bfcux6570}{%
\subsection{4.6　多通道耦合与 Johnson
矩阵对数导数}\label{ux591aux901aux9053ux8026ux5408ux4e0e-johnson-ux77e9ux9635ux5bf9ux6570ux5bfcux6570}}

设通道波函数为 \(\mathbf F(r)\)，则耦合径向方程为

\[
-\frac{\hbar^2}{2\mu}
\mathbf F''(r)
+
\left[
\mathbf V(r)
+
\mathbf E_{\mathrm{th}}
+
\frac{\hbar^2\mathbf L^2}{2\mu r^2}
\right]
\mathbf F(r)
=
E\mathbf F(r).
\]

在原子单位下，

\[
-\frac{1}{2\mu}
\mathbf F''(r)
+
\left[
\mathbf V(r)
+
\mathbf E_{\mathrm{th}}
+
\frac{\mathbf L^2}{2\mu r^2}
\right]
\mathbf F(r)
=
E\mathbf F(r).
\]

多通道 Log-Derivative 方法传播矩阵

\[
\mathbf Y(r)
=
\mathbf F'(r)\mathbf F^{-1}(r).
\]

在匹配半径处，由 \(\mathbf Y\) 构造 \(\mathbf K\) 矩阵，再由 Cayley
变换得到幺正散射矩阵：

\[
\mathbf S
=
(\mathbf I+i\mathbf K)(\mathbf I-i\mathbf K)^{-1}.
\]

跃迁矩阵为

\[
\mathbf T=\mathbf S-\mathbf I.
\]

态到态跃迁概率为

\[
P_{i\rightarrow j}=|S_{ij}|^2.
\]

\hypertarget{ux5206ux6bb5ux7f51ux683cux4e0eux5c40ux90e8ux6247ux533a}{%
\subsection{4.7　分段网格与局部扇区}\label{ux5206ux6bb5ux7f51ux683cux4e0eux5c40ux90e8ux6247ux533a}}

在长程势或强局域势问题中，单一均匀网格往往效率低。分段网格将径向区间划分为多个局部扇区：

\[
[r_{\min},r_{\max}]
=
\bigcup_{s=1}^{S}
[r_s,r_{s+1}].
\]

每个扇区可使用不同步长
\(\Delta r_s\)。这种方法在近区加密、远区稀疏，可以显著减少总网格数，同时保持误差可控。

GeneralModule 的 \texttt{segmented\_grid\_t}
存储每个扇区的起点、终点、步长、点数和全局索引。相关函数包括：

\begin{itemize}
\tightlist
\item
  \texttt{create\_segmented\_grid}
\item
  \texttt{calc\_scattering\_length\_segmented\_numerov}
\item
  \texttt{calc\_scattering\_wavefunction\_segmented\_ti}
\item
  \texttt{calc\_phase\_shift\_segmented}
\item
  \texttt{calc\_multichannel\_close\_coupling\_segmented\_logder}
\end{itemize}

\hypertarget{feshbach-ux5171ux632f}{%
\subsection{4.8　Feshbach 共振}\label{feshbach-ux5171ux632f}}

闭通道束缚态与开通道连续谱耦合会产生 Feshbach
共振。散射长度随磁场变化可写为

\[
a_s(B)
=
a_{\mathrm{bg}}
\left[
1-\frac{\Delta B}{B-B_0}
\right].
\]

共振宽度为 \(\Delta B\)，共振位置为 \(B_0\)。对超冷原子，Feshbach
共振是调控相互作用强度的核心手段。

GeneralModule 的 \texttt{calc\_feshbach\_resonance\_scan}
在磁场网格上计算 \(a_s(B)\)，并拟合共振参数。

\hypertarget{ux5bf9ux5e94ux6e90ux7801-10}{%
\subsubsection{对应源码}\label{ux5bf9ux5e94ux6e90ux7801-10}}

\begin{itemize}
\tightlist
\item
  \href{https://github.com/l1Ha/QuantumGeneralModule/blob/main/src/mod_field_scattering.f90}{src/mod\_field\_scattering.f90}
\item
  二维码：\texttt{qr/src\_\_mod\_field\_scattering.f90.png}
\end{itemize}

\includegraphics[width=2cm,height=\textheight]{qr/src__mod_field_scattering.f90.png}

\hypertarget{ux542bux65f6ux6ce2ux5305ux6563ux5c04}{%
\subsection{4.9　含时波包散射}\label{ux542bux65f6ux6ce2ux5305ux6563ux5c04}}

初始高斯波包为

\[
\psi(x,0)
=
\left(\frac{1}{2\pi\sigma^2}\right)^{1/4}
\exp\left[
-\frac{(x-x_0)^2}{4\sigma^2}
+ik_0x
\right].
\]

其动量空间振幅为

\[
\tilde\psi(k,0)
=
\left(\frac{2\sigma^2}{\pi}\right)^{1/4}
\exp\left[
-\sigma^2(k-k_0)^2
-ikx_0
\right].
\]

定义左、右通量积分：

\[
T(t)
=
\int_{x>x_m}|\psi(x,t)|^2\,dx,
\]

\[
R(t)
=
\int_{x<x_l}|\psi(x,t)|^2\,dx.
\]

当 \(t\rightarrow\infty\) 时，\(T+R\) 应接近
1，除非存在吸收或非弹性通道。

定态散射给出能量分辨的
\(S(E)\)；含时波包通过傅里叶变换也可得到能量分辨量。两者应满足：

\begin{itemize}
\tightlist
\item
  透射峰位置一致；
\item
  共振宽度一致；
\item
  相移或延迟一致；
\item
  总通量守恒。
\end{itemize}

GeneralModule 的 \texttt{ex07}
示例专门用于定态和含时散射波函数交叉验证。

\hypertarget{ux5bf9ux5e94ux6e90ux7801-11}{%
\subsubsection{对应源码}\label{ux5bf9ux5e94ux6e90ux7801-11}}

\begin{itemize}
\tightlist
\item
  \href{https://github.com/l1Ha/QuantumGeneralModule/blob/main/src/mod_td_scattering.f90}{src/mod\_td\_scattering.f90}
\item
  二维码：\texttt{qr/src\_\_mod\_td\_scattering.f90.png}
\end{itemize}

\includegraphics[width=2cm,height=\textheight]{qr/src__mod_td_scattering.f90.png}

\hypertarget{ux5206ux5b50ux8f6cux632fux54c8ux5bc6ux987fux91cfux4e0eux5149ux8c31}{%
\subsection{4.10　分子转振哈密顿量与光谱}\label{ux5206ux5b50ux8f6cux632fux54c8ux5bc6ux987fux91cfux4e0eux5149ux8c31}}

双分子振转哈密顿量可近似为

\[
\hat H
=
-\frac{\hbar^2}{2\mu}
\frac{d^2}{dR^2}
+
\frac{\hat{\mathbf J}^2}{2\mu R^2}
+
V(R).
\]

在原子单位下，

\[
\hat H
=
-\frac{1}{2\mu}
\frac{d^2}{dR^2}
+
\frac{\hat{\mathbf J}^2}{2\mu R^2}
+
V(R).
\]

Morse 势为

\[
V(R)
=
D_e
\left[
1-e^{-\beta(R-R_e)}
\right]^2.
\]

解析能级为

\[
E_v
=
\hbar\omega_e\left(v+\frac12\right)
-
\hbar\omega_ex_e\left(v+\frac12\right)^2,
\]

其中

\[
\omega_e=\beta\sqrt{\frac{2D_e}{\mu}},
\]

\[
\omega_ex_e=\frac{\omega_e^2}{4D_e}.
\]

电偶极跃迁矩阵元为

\[
M_{v'v''J'J''}
=
\int
\chi_{v'J'}^{*}(R)
\mu(R)
\chi_{v''J''}(R)\,dR.
\]

Franck--Condon 因子为

\[
F_{v'v''}
=
\left|
\int
\chi_{v'}^{*}(R)\chi_{v''}(R)\,dR
\right|^2.
\]

\hypertarget{ux5bf9ux5e94ux6e90ux7801-12}{%
\subsubsection{对应源码}\label{ux5bf9ux5e94ux6e90ux7801-12}}

\begin{itemize}
\tightlist
\item
  \href{https://github.com/l1Ha/QuantumGeneralModule/blob/main/src/mod_rovibrational.f90}{src/mod\_rovibrational.f90}
\item
  二维码：\texttt{qr/src\_\_mod\_rovibrational.f90.png}
\end{itemize}

\includegraphics[width=2cm,height=\textheight]{qr/src__mod_rovibrational.f90.png}

\hypertarget{ux5149ux788eux7247ux52a8ux80fdux91caux653eux8c31}{%
\subsection{4.11　光碎片动能释放谱}\label{ux5149ux788eux7247ux52a8ux80fdux91caux653eux8c31}}

自相关函数为

\[
C(t)=\langle\psi(0)|\psi(t)\rangle.
\]

吸收谱为

\[
\sigma(\omega)
\propto
\omega
\operatorname{Re}
\int_0^\infty
C(t)e^{i\omega t}W(t)\,dt,
\]

其中 \(W(t)\) 是时间窗函数。动能释放谱为

\[
P(E_k)
=
\sqrt{\frac{2E_k}{m}}
|A(E_k)|^2.
\]

\hypertarget{ux5bf9ux5e94ux6e90ux7801-13}{%
\subsubsection{对应源码}\label{ux5bf9ux5e94ux6e90ux7801-13}}

\begin{itemize}
\tightlist
\item
  \href{https://github.com/l1Ha/QuantumGeneralModule/blob/main/src/mod_photofragment_flux.f90}{src/mod\_photofragment\_flux.f90}
\item
  二维码：\texttt{qr/src\_\_mod\_photofragment\_flux.f90.png}
\end{itemize}

\includegraphics[width=2cm,height=\textheight]{qr/src__mod_photofragment_flux.f90.png}

\hypertarget{ux975eux7eddux70edux52a8ux529bux5b66ux4e0e-fssh}{%
\subsection{4.12　非绝热动力学与
FSSH}\label{ux975eux7eddux70edux52a8ux529bux5b66ux4e0e-fssh}}

设电子态为 \(|\phi_1\rangle,|\phi_2\rangle\)，核坐标为
\(R\)，非绝热耦合为

\[
d_{12}(R)
=
\langle\phi_1|\partial_R|\phi_2\rangle.
\]

核速度为 \(v\) 时，电子振幅方程为

\[
i\hbar\dot c_1
=
E_1c_1
-
i\hbar v d_{12}c_2,
\]

\[
i\hbar\dot c_2
=
E_2c_2
+
i\hbar v d_{12}c_1.
\]

FSSH 使用经典核轨迹和量子电子振幅。在活性态 \(j\) 上，向态 \(k\)
的跳跃概率为

\[
g_{j\rightarrow k}
=
\max
\left[
0,
-\frac{2\Delta t}{|c_j|^2}
\operatorname{Re}
\left(
c_j^{*}c_k\,\mathbf v\cdot\mathbf d_{jk}
\right)
\right].
\]

若随机数小于
\(g_{j\rightarrow k}\)，则尝试跳跃。跳跃后需要沿非绝热耦合方向重标动量以保持总能量：

\[
\mathbf p'\cdot\hat{\mathbf e}_{jk}
=
\operatorname{sign}
\left(
\mathbf p\cdot\hat{\mathbf e}_{jk}
\right)
\sqrt{
(\mathbf p\cdot\hat{\mathbf e}_{jk})^2
-
2M\Delta E
}.
\]

若动量不足，则为 frustrated hop。Ehrenfest 方法使用平均场力：

\[
\mathbf F
=
-\left\langle\Psi\left|\frac{\partial\hat H}{\partial\mathbf R}\right|\Psi\right\rangle.
\]

核仍沿经典轨迹运动，电子态保持相干演化。与 FSSH 相比，Ehrenfest
不能自然描述分支化过程。

\begin{itemize}
\tightlist
\item
  FSSH 源码：
\end{itemize}

\href{https://github.com/l1Ha/QuantumGeneralModule/blob/main/src/mod_surface_hopping_fssh.f90}{src/mod\_surface\_hopping\_fssh.f90}

\begin{itemize}
\tightlist
\item
  二维码：
\end{itemize}

\includegraphics[width=2cm,height=\textheight]{qr/src__mod_surface_hopping_fssh.f90.png}

\hypertarget{ux9ad8ux7ea7ux4e13ux9898ux7684ux7edfux4e00ux7814ux7a76ux6846ux67b6}{%
\subsection{4.13　高级专题的统一研究框架}\label{ux9ad8ux7ea7ux4e13ux9898ux7684ux7edfux4e00ux7814ux7a76ux6846ux67b6}}

GeneralModule 还包含多个研究级专题模块。它们不是孤立
API，而是同一套数学结构的延伸。

\begin{longtable}[]{@{}lll@{}}
\toprule
专题 & 代表模块 & 核心问题\tabularnewline
\midrule
\endhead
Efimov 物理 & \texttt{mod\_three\_body\_recombination} &
三体共振、离散标度不变性\tabularnewline
CIR & \texttt{mod\_confined\_scattering} &
低维受限诱导共振\tabularnewline
Fano/CCR & \texttt{mod\_autoionization\_fano} &
自电离线型与复标度\tabularnewline
强场 NSDI & \texttt{mod\_strong\_field\_nsdi} &
隧穿、重碰撞、双电离\tabularnewline
ATAS & \texttt{mod\_attosecond\_transient\_absorption} &
阿秒瞬态吸收\tabularnewline
PECD & \texttt{mod\_bicircular\_pecd} &
手性光电子圆二色性\tabularnewline
Rydberg blockade & \texttt{mod\_rydberg\_blockade} &
里德堡阻塞与多体\tabularnewline
表面散射 & \texttt{mod\_surface\_scattering} &
表面衍射与吸附\tabularnewline
Eley--Rideal & \texttt{mod\_surface\_reaction\_er} &
气---固反应\tabularnewline
电子摩擦 & \texttt{mod\_surface\_electronic\_friction} &
非绝热耗散\tabularnewline
GIFAD & \texttt{mod\_grazing\_fast\_atom\_diffraction} &
掠入射衍射\tabularnewline
Dirac 原子 & \texttt{mod\_relativistic\_atomic} &
相对论精细结构\tabularnewline
RIXS & \texttt{mod\_resonant\_xray\_scattering} & 共振非弹性 X
射线散射\tabularnewline
Penning 电离 & \texttt{mod\_penning\_associative\_ionization} &
亚稳态碰撞电离\tabularnewline
\bottomrule
\end{longtable}

这些模块的共同验证策略是：先在解析极限或低维模型中验证，再逐步加入耦合、外场、耗散和多通道效应。

\hypertarget{ux53c2ux8003ux6587ux732e-2}{%
\subsection{4.14　参考文献}\label{ux53c2ux8003ux6587ux732e-2}}

\begin{enumerate}
\def\labelenumi{\arabic{enumi}.}
\tightlist
\item
  B. R. Johnson, ``The multichannel log-derivative method for treating
  reactive collisions'', \emph{J. Comput. Phys.} \textbf{13}, 445
  (1973). DOI: 10.1016/0021-9991(73)90049-1.
\item
  D. E. Manolopoulos, ``An improved log-derivative method for solving
  the radial Schrödinger equation'', \emph{J. Chem. Phys.} \textbf{85},
  6425 (1986). DOI: 10.1063/1.451472.
\item
  H. Feshbach, ``Unified theory of nuclear reactions'', \emph{Ann. Phys.
  (N.Y.)} \textbf{5}, 357 (1958). DOI: 10.1016/0003-4916(58)90007-1.
\item
  C. Chin, R. Grimm, P. S. Julienne, and E. Tiesinga, ``Feshbach
  resonances in ultracold gases'', \emph{Rev.~Mod. Phys.} \textbf{82},
  1225 (2010). DOI: 10.1103/RevModPhys.82.1225.
\item
  J. C. Tully, ``Molecular dynamics with electronic transitions'',
  \emph{J. Chem. Phys.} \textbf{93}, 1061 (1990). DOI: 10.1063/1.459170.
\item
  GeneralModule
  源码与文档：\url{https://github.com/l1Ha/QuantumGeneralModule}
\end{enumerate}

\hypertarget{ux7b2cux4e94ux90e8ux5206-ux5de5ux7a0bux5b9eux73b0ux9a8cux8bc1ux4e0eux590dux73b0}{%
\section{第五部分　工程实现、验证与复现}\label{ux7b2cux4e94ux90e8ux5206-ux5de5ux7a0bux5b9eux73b0ux9a8cux8bc1ux4e0eux590dux73b0}}

\hypertarget{ux4ee3ux7801ux67b6ux6784ux7684ux5206ux5c42ux539fux5219}{%
\subsection{5.1　代码架构的分层原则}\label{ux4ee3ux7801ux67b6ux6784ux7684ux5206ux5c42ux539fux5219}}

GeneralModule 的工程结构不是按物理主题随机堆叠，而是按依赖方向分层：

\begin{enumerate}
\def\labelenumi{\arabic{enumi}.}
\tightlist
\item
  物理常数与单位：\texttt{mod\_constants}；
\item
  数学基础：\texttt{mod\_special\_functions}、\texttt{mod\_linear\_algebra}；
\item
  表示与离散：\texttt{mod\_dvr\_grid}、\texttt{mod\_interpolation}、\texttt{mod\_io\_utils}；
\item
  通用动力学：波包传播、开放系统、最优控制；
\item
  物理专题：散射、表面、相对论、RIXS、Penning 等；
\item
  聚合门面：\texttt{general\_module}。
\end{enumerate}

低层模块不应反向依赖高层模块，也不应依赖聚合模块
\texttt{general\_module}。这保证单元测试可以独立编译，也避免循环依赖导致编译顺序不可维护。

模块接口通常遵循以下规范：

\begin{itemize}
\tightlist
\item
  所有公共实体显式 \texttt{public}；
\item
  数值类型统一使用 \texttt{dp\ =\textgreater{}\ real64}；
\item
  输入输出通过
  \texttt{intent(in)}、\texttt{intent(out)}、\texttt{intent(inout)}
  说明；
\item
  大多数无副作用计算标记为 \texttt{pure}；
\item
  动态数组使用 \texttt{allocatable}；
\item
  错误通过状态码 \texttt{stat} 返回，而不是直接终止进程。
\end{itemize}

这种接口风格便于编译器检查，也便于 AI 辅助生成代码时减少隐式错误。

\hypertarget{ux5bf9ux5e94ux6e90ux7801-14}{%
\subsubsection{对应源码}\label{ux5bf9ux5e94ux6e90ux7801-14}}

\begin{itemize}
\tightlist
\item
  \href{https://github.com/l1Ha/QuantumGeneralModule/blob/main/src/mod_constants.f90}{src/mod\_constants.f90}
\item
  二维码：
\end{itemize}

\includegraphics[width=2cm,height=\textheight]{qr/src__mod_constants.f90.png}

\begin{itemize}
\tightlist
\item
  \href{https://github.com/l1Ha/QuantumGeneralModule/blob/main/src/general_module.f90}{src/general\_module.f90}
\end{itemize}

\hypertarget{ux4e09ux5957ux6784ux5efaux4f53ux7cfbux7684ux5206ux5de5}{%
\subsection{5.2　三套构建体系的分工}\label{ux4e09ux5957ux6784ux5efaux4f53ux7cfbux7684ux5206ux5de5}}

GeneralModule 同时提供 \texttt{Makefile}、\texttt{fpm.toml} 和
\texttt{CMakeLists.txt}。三者不是重复实现，而是服务于不同工作流。

\begin{longtable}[]{@{}llll@{}}
\toprule
构建入口 & 主要用途 & 优点 & 局限\tabularnewline
\midrule
\endhead
Makefile & 快速本地构建 & 明确拓扑顺序 & 平台相关性较强\tabularnewline
fpm.toml & Fortran 包管理 & 标准化、依赖清晰 &
自动发现关闭\tabularnewline
CMakeLists.txt & 跨平台与 IDE & 对象库、静态库、共享库、CTest &
配置复杂\tabularnewline
\bottomrule
\end{longtable}

\texttt{Makefile} 使用显式模块顺序编译，保证 \texttt{.mod}
文件依赖正确。\texttt{fpm.toml} 显式列出 41 个测试和 36
个示例，避免隐式自动发现。\texttt{CMakeLists.txt}
使用对象库组织源码，并生成静态库、共享库和可执行文件。

\begin{itemize}
\tightlist
\item
  Makefile：
\end{itemize}

\href{https://github.com/l1Ha/QuantumGeneralModule/blob/main/Makefile}{Makefile}

\begin{itemize}
\tightlist
\item
  二维码：
\end{itemize}

\includegraphics[width=2cm,height=\textheight]{qr/Makefile.png}

\begin{itemize}
\tightlist
\item
  fpm 配置：
\end{itemize}

\href{https://github.com/l1Ha/QuantumGeneralModule/blob/main/fpm.toml}{fpm.toml}

\begin{itemize}
\tightlist
\item
  二维码：
\end{itemize}

\includegraphics[width=2cm,height=\textheight]{qr/fpm.toml.png}

\begin{itemize}
\tightlist
\item
  CMake：
\end{itemize}

\href{https://github.com/l1Ha/QuantumGeneralModule/blob/main/CMakeLists.txt}{CMakeLists.txt}

\begin{itemize}
\tightlist
\item
  二维码：
\end{itemize}

\includegraphics[width=2cm,height=\textheight]{qr/CMakeLists.txt.png}

\hypertarget{ux6d4bux8bd5ux4f53ux7cfbux7684ux9a8cux8bc1ux5c42ux6b21}{%
\subsection{5.3　测试体系的验证层次}\label{ux6d4bux8bd5ux4f53ux7cfbux7684ux9a8cux8bc1ux5c42ux6b21}}

GeneralModule 当前包含 41 个 Fortran 测试程序和 36
个物理示例。测试框架由项目自实现，每个测试程序维护计数器和断言逻辑。

测试应覆盖以下层次：

\begin{longtable}[]{@{}lll@{}}
\toprule
层次 & 检查内容 & 典型失败模式\tabularnewline
\midrule
\endhead
单位与常数 & 往返转换误差 & 错误量纲、数量级错\tabularnewline
特殊函数 & 选择定则、归一化、对称性 & 相位约定错误\tabularnewline
本征问题 & 正交性、解析谱、条件数 & 非厄米离散\tabularnewline
DVR & 网格收敛、边界效应 & 网格欠采样\tabularnewline
传播 & 范数、能量、动量守恒 & 时间步过大\tabularnewline
散射 & 光学定理、幺正性、解析极限 & 匹配半径不当\tabularnewline
开放系统 & 迹、厄米性、正定性 & 强制归一化掩盖误差\tabularnewline
非绝热 & 分支比、能量守恒、统计收敛 & 轨迹数不足\tabularnewline
工程接口 & 编译、链接、可执行性 & 模块顺序错误\tabularnewline
\bottomrule
\end{longtable}

\hypertarget{ux5355ux5143ux6d4bux8bd5ux793aux4f8b}{%
\subsubsection{单元测试示例}\label{ux5355ux5143ux6d4bux8bd5ux793aux4f8b}}

常数和单位转换测试应包含

\[
\left|
\operatorname{from}_{au}(\operatorname{to}_{au}(x,u),u)-x
\right|
\le
\epsilon|x|.
\]

DVR 测试应检查

\[
\sum_i w_i\psi_i^{(k)}\psi_i^{(l)}
=
\delta_{kl}.
\]

散射测试应检查

\[
\sigma_{\mathrm{tot}}
=
\frac{4\pi}{k}
\operatorname{Im}f(0).
\]

开放系统测试应检查

\[
\left|\operatorname{Tr}\rho(t)-1\right|<\epsilon,
\qquad
\|\rho(t)-\rho^{\dagger}(t)\|<\epsilon.
\]

\hypertarget{ux5bf9ux5e94ux6e90ux7801-15}{%
\subsubsection{对应源码}\label{ux5bf9ux5e94ux6e90ux7801-15}}

\begin{itemize}
\tightlist
\item
  \href{https://github.com/l1Ha/QuantumGeneralModule/blob/main/tests/run_all_tests.sh}{tests/run\_all\_tests.sh}
\item
  二维码：
\end{itemize}

\includegraphics[width=2cm,height=\textheight]{qr/tests__run_all_tests.sh.png}

\begin{itemize}
\item
  \href{https://github.com/l1Ha/QuantumGeneralModule/blob/main/tests/test_dvr_grid.f90}{tests/test\_dvr\_grid.f90}
\item
  二维码：
\end{itemize}

\includegraphics[width=2cm,height=\textheight]{qr/tests__test_dvr_grid.f90.png}

\begin{itemize}
\item
  \href{https://github.com/l1Ha/QuantumGeneralModule/blob/main/tests/test_ti_scattering.f90}{tests/test\_ti\_scattering.f90}
\item
  二维码：
\end{itemize}

\includegraphics[width=2cm,height=\textheight]{qr/tests__test_ti_scattering.f90.png}

\hypertarget{ux6301ux7eedux96c6ux6210ux4e0eux5e73ux53f0ux9a8cux8bc1}{%
\subsection{5.4　持续集成与平台验证}\label{ux6301ux7eedux96c6ux6210ux4e0eux5e73ux53f0ux9a8cux8bc1}}

GitHub Actions 配置位于：

\href{https://github.com/l1Ha/QuantumGeneralModule/blob/main/.github/workflows/ci.yml}{.github/workflows/ci.yml}

其作用是自动化编译、测试和示例运行，防止平台相关错误进入主干。CI
应至少检查：

\begin{enumerate}
\def\labelenumi{\arabic{enumi}.}
\tightlist
\item
  源码能否编译；
\item
  静态库和共享库能否生成；
\item
  单元测试是否通过；
\item
  物理示例是否运行；
\item
  Python 测试是否通过；
\item
  输出文件是否存在且非空。
\end{enumerate}

CI 不能替代物理收敛研究，但可以快速暴露接口和平台相关问题。

\hypertarget{python-ux4ea4ux53c9ux9a8cux8bc1ux4e0eux53efux89c6ux5316}{%
\subsection{5.5　Python
交叉验证与可视化}\label{python-ux4ea4ux53c9ux9a8cux8bc1ux4e0eux53efux89c6ux5316}}

Python 包 \texttt{pygenmod} 使用 NumPy 和 Matplotlib
对部分算法进行独立实现、交叉验证和可视化。它不是 Fortran
库的替代品，而是验证与绘图工具。

\begin{itemize}
\tightlist
\item
  NumPy 源码：
\end{itemize}

\href{https://github.com/l1Ha/QuantumGeneralModule/blob/main/python/pygenmod/visualizer.py}{python/pygenmod/visualizer.py}

\begin{itemize}
\tightlist
\item
  二维码：
\end{itemize}

\includegraphics[width=2cm,height=\textheight]{qr/python__pygenmod__visualizer.py.png}

Python 测试位于：

\href{https://github.com/l1Ha/QuantumGeneralModule/blob/main/python/test_pygenmod.py}{python/test\_pygenmod.py}

\begin{itemize}
\tightlist
\item
  二维码：
\end{itemize}

\includegraphics[width=2cm,height=\textheight]{qr/python__test_pygenmod.py.png}

推荐流程是：

\begin{enumerate}
\def\labelenumi{\arabic{enumi}.}
\tightlist
\item
  Fortran 输出原始数据；
\item
  Python 独立计算参考值；
\item
  绘制对比图；
\item
  保存输入参数、随机种子、版本和命令；
\item
  将图和数据放入同一实验目录。
\end{enumerate}

\hypertarget{ux53efux590dux73b0ux5b9eux9a8cux8bb0ux5f55}{%
\subsection{5.6　可复现实验记录}\label{ux53efux590dux73b0ux5b9eux9a8cux8bb0ux5f55}}

数值实验应至少记录：

\begin{longtable}[]{@{}ll@{}}
\toprule
字段 & 内容\tabularnewline
\midrule
\endhead
实验名称 & 物理问题和目标\tabularnewline
源码版本 & commit、文件时间或 tag\tabularnewline
环境 & 编译器、系统、CPU、Python 版本\tabularnewline
参数 & 单位、网格、初态、时间步\tabularnewline
随机性 & 种子或确定性声明\tabularnewline
理论基准 & 解析解、文献值、守恒量\tabularnewline
输出 & 能量、截面、范数、分支比\tabularnewline
收敛判据 & 相对误差和阈值\tabularnewline
复现命令 & 完整脚本\tabularnewline
\bottomrule
\end{longtable}

没有上述记录，即使结果看起来正确，也难以复现或定位误差。

\hypertarget{ux5f15ux7528ux4e0eux5f00ux6e90ux89c4ux8303}{%
\subsection{5.7　引用与开源规范}\label{ux5f15ux7528ux4e0eux5f00ux6e90ux89c4ux8303}}

本书引用的文献用于说明理论来源和算法背景。GeneralModule
的实现是独立编写的；引用不表示复制第三方源码。若未来引入第三方代码，必须：

\begin{enumerate}
\def\labelenumi{\arabic{enumi}.}
\tightlist
\item
  保留其许可证；
\item
  在源码头注明来源；
\item
  在 \texttt{CHANGELOG} 中说明修改范围；
\item
  在文档中给出原始仓库或论文链接。
\end{enumerate}

\hypertarget{ux7b2cux516dux90e8ux5206-ux9ad8ux7ea7ux4e13ux9898ux4e0eux4ee3ux7801ux6620ux5c04}{%
\section{第六部分　高级专题与代码映射}\label{ux7b2cux516dux90e8ux5206-ux9ad8ux7ea7ux4e13ux9898ux4e0eux4ee3ux7801ux6620ux5c04}}

\hypertarget{ux4e09ux4f53ux590dux5408ux4e0e-efimov-ux7269ux7406}{%
\subsection{6.1　三体复合与 Efimov
物理}\label{ux4e09ux4f53ux590dux5408ux4e0e-efimov-ux7269ux7406}}

在共振两体相互作用极限下（\(|a|\gg r_0\)），全同玻色子三体超径向波函数方程导出特征超越方程

\[
s\cosh\left(\frac{\pi s}{2}\right)
-
\frac{8}{\sqrt{3}}\sinh\left(\frac{\pi s}{6}\right)
=
0.
\]

该方程在实轴上具有唯一正根

\[
s_0\approx1.0062378.
\]

它决定了有效吸引超径向势

\[
V_{\mathrm{eff}}(R)
=
-\frac{s_0^2+1/4}{2\mu R^2}.
\]

三体体系破缺连续标度对称性，展现离散标度不变性：

\[
\lambda=e^{\pi/s_0}\approx22.69438.
\]

Braaten 和 Hammer 对少体普适性作了系统综述【Braaten \& Hammer 2006】。

\hypertarget{ux5bf9ux5e94ux6e90ux7801-16}{%
\subsubsection{对应源码}\label{ux5bf9ux5e94ux6e90ux7801-16}}

\begin{itemize}
\tightlist
\item
  \href{https://github.com/l1Ha/QuantumGeneralModule/blob/main/src/mod_three_body_recombination.f90}{src/mod\_three\_body\_recombination.f90}
\end{itemize}

\hypertarget{ux4f4eux7ef4ux53d7ux9650ux6563ux5c04ux4e0e-cir}{%
\subsection{6.2　低维受限散射与
CIR}\label{ux4f4eux7ef4ux53d7ux9650ux6563ux5c04ux4e0e-cir}}

在横向谐振子紧束缚光波导中，

\[
V_\perp(\rho)=\frac12\mu\omega_\perp^2\rho^2,
\]

振子特征长度为

\[
a_\perp=\sqrt{\frac{\hbar}{\mu\omega_\perp}}.
\]

准一维有效接触势耦合常数由横向模态多体格林函数重整化：

\[
g_{\mathrm{1D}}(a_s)
=
\frac{2\hbar^2a_s}{\mu a_\perp^2}
\frac{1}{1-C a_s/a_\perp},
\qquad
C=\frac{|\zeta(1/2)|}{\sqrt2}\approx1.0326.
\]

当三维散射长度逼近

\[
a_{\mathrm{CIR}}=\frac{a_\perp}{C}
\approx0.96843a_\perp
\]

时，\(g_{\mathrm{1D}}\rightarrow\pm\infty\)，发生约束诱导共振。

Olshanii 给出了准一维约束诱导共振理论【Olshanii 1998】。

\hypertarget{ux5bf9ux5e94ux6e90ux7801-17}{%
\subsubsection{对应源码}\label{ux5bf9ux5e94ux6e90ux7801-17}}

\begin{itemize}
\tightlist
\item
  \href{https://github.com/l1Ha/QuantumGeneralModule/blob/main/src/mod_confined_scattering.f90}{src/mod\_confined\_scattering.f90}
\end{itemize}

\hypertarget{fano-ux5171ux632fux4e0eux590dux5750ux6807ux65cbux8f6c}{%
\subsection{6.3　Fano
共振与复坐标旋转}\label{fano-ux5171ux632fux4e0eux590dux5750ux6807ux65cbux8f6c}}

当离散准束缚态 \(|\phi\rangle\) 与连续态 \(|\psi_E\rangle\)
耦合时，吸收截面表现为非对称 Fano 线型：

\[
\sigma(\epsilon)
=
\sigma_0\frac{(q+\epsilon)^2}{1+\epsilon^2},
\qquad
\epsilon=\frac{E-E_0}{\Gamma/2}.
\]

自电离寿命为

\[
\tau=\frac{\hbar}{\Gamma}.
\]

复坐标旋转法通过非厄米变换

\[
U(\theta)=\exp(i\theta r\partial_r)
\]

将坐标旋转为

\[
r\rightarrow re^{i\theta}.
\]

连续谱旋转后，共振极点

\[
E_{\mathrm{res}}=E_R-i\frac{\Gamma}{2}
\]

可从第二黎曼叶暴露出来。

Fano 给出了组态相互作用理论【Fano 1961】；Reinhardt 和 Moiseyev
讨论了复坐标旋转方法【Reinhardt 1982；Moiseyev 1998】。

\hypertarget{ux5bf9ux5e94ux6e90ux7801-18}{%
\subsubsection{对应源码}\label{ux5bf9ux5e94ux6e90ux7801-18}}

\begin{itemize}
\tightlist
\item
  \href{https://github.com/l1Ha/QuantumGeneralModule/blob/main/src/mod_autoionization_fano.f90}{src/mod\_autoionization\_fano.f90}
\end{itemize}

\hypertarget{ux4ea4ux53c9ux7535ux78c1ux573aux4e0eux5206ux5b50ux53d6ux5411}{%
\subsection{6.4　交叉电磁场与分子取向}\label{ux4ea4ux53c9ux7535ux78c1ux573aux4e0eux5206ux5b50ux53d6ux5411}}

极性开壳层分子在非共线静电场和磁场中的哈密顿量可写为

\[
\hat H
=
B_{\mathrm{rot}}\hat{\mathbf J}^2
-
d\mathcal E\cos\theta_R
+
g_S\mu_B
\left(
B_z\hat S_z+B_x\hat S_x
\right).
\]

横向磁场分量 \(B_x=B\sin\theta_{EB}\) 产生非对角算符

\[
\hat S_x=\frac{\hat S_++\hat S_-}{2},
\]

驱动 \(\Delta M_S=\pm1\)
耦合，破坏沿电场轴的柱对称性，形成可由倾角连续调控的 Stark--Zeeman
反交叉谱。

\hypertarget{ux5bf9ux5e94ux6e90ux7801-19}{%
\subsubsection{对应源码}\label{ux5bf9ux5e94ux6e90ux7801-19}}

\begin{itemize}
\tightlist
\item
  \href{https://github.com/l1Ha/QuantumGeneralModule/blob/main/src/mod_crossed_field_scattering.f90}{src/mod\_crossed\_field\_scattering.f90}
\end{itemize}

\hypertarget{ux4e09ux539fux5b50ux53cdux5e94ux4e0e-berry-ux76f8ux4f4d}{%
\subsection{6.5　三原子反应与 Berry
相位}\label{ux4e09ux539fux5b50ux53cdux5e94ux4e0e-berry-ux76f8ux4f4d}}

在 \(A+BC\) 反应中，Jacobi 坐标 \((r,R,\gamma)\) 与三原子核间距
\((r_{12},r_{23},r_{31})\) 存在解析映射。London--Eyring--Polanyi--Sato
势能面为

\[
V(r_{12},r_{23},r_{31})
=
Q_1+Q_2+Q_3
-
\sqrt{
\frac12
\left[
(J_1-J_2)^2
+
(J_2-J_3)^2
+
(J_3-J_1)^2
\right]
}.
\]

在锥形交叉附近，绝热波函数环绕退化点一周获得 Berry 相位

\[
\Phi_B=\oint\mathbf A\cdot d\mathbf R=\pi.
\]

Longuet-Higgins 等讨论了 Jahn--Teller 与几何相位【Longuet-Higgins et
al.~1958】；Berry 给出了量子几何相位的一般表述【Berry 1984】。

\hypertarget{ux5bf9ux5e94ux6e90ux7801-20}{%
\subsubsection{对应源码}\label{ux5bf9ux5e94ux6e90ux7801-20}}

\begin{itemize}
\tightlist
\item
  \href{https://github.com/l1Ha/QuantumGeneralModule/blob/main/src/mod_triatomic_geometry.f90}{src/mod\_triatomic\_geometry.f90}
\end{itemize}

\hypertarget{ux65cbux91cf-bec-ux4e0eux81eaux65cbux6df7ux5408}{%
\subsection{6.6　旋量 BEC
与自旋混合}\label{ux65cbux91cf-bec-ux4e0eux81eaux65cbux6df7ux5408}}

自旋 \(F=1\) 玻色气体的低能碰撞由总自旋 \(F_{\mathrm{tot}}=0,2\)
两个通道的散射长度决定：

\[
c_0=\frac{4\pi\hbar^2(a_0+2a_2)}{3m},
\]

\[
c_2=\frac{4\pi\hbar^2(a_2-a_0)}{3m}.
\]

在单模近似下，旋量分量满足非线性方程

\[
i\hbar\frac{d\zeta_{\pm1}}{dt}
=
\left[
q_Z+c_2n(\rho_0+\rho_{\pm1}-\rho_{\mp1})
\right]\zeta_{\pm1}
+
c_2n\zeta_0^2\zeta_{\mp1}^{*},
\]

\[
i\hbar\frac{d\zeta_0}{dt}
=
c_2n(\rho_{+1}+\rho_{-1})\zeta_0
+
2c_2n\zeta_{+1}\zeta_{-1}\zeta_0^{*}.
\]

系统守恒总几率与磁化强度

\[
m_z=|\zeta_{+1}|^2-|\zeta_{-1}|^2.
\]

Ho 和 Ohmi--Machida 给出了旋量 BEC 理论【Ho 1998；Ohmi \& Machida
1998】。

\hypertarget{ux5bf9ux5e94ux6e90ux7801-21}{%
\subsubsection{对应源码}\label{ux5bf9ux5e94ux6e90ux7801-21}}

\begin{itemize}
\tightlist
\item
  \href{https://github.com/l1Ha/QuantumGeneralModule/blob/main/src/mod_spinor_bec.f90}{src/mod\_spinor\_bec.f90}
\end{itemize}

\hypertarget{ux8868ux9762ux6563ux5c04ux4e0eux7535ux5b50ux6469ux64e6}{%
\subsection{6.7　表面散射与电子摩擦}\label{ux8868ux9762ux6563ux5c04ux4e0eux7535ux5b50ux6469ux64e6}}

表面散射需要同时处理表面周期势、入射动量和吸附通道。Eley--Rideal
机理描述气相原子或分子直接撞击表面吸附物并发生反应。电子摩擦则将金属表面电子---空穴对激发等效为阻尼力：

\[
\mathbf F_{\mathrm{fric}}=-\eta(\mathbf R)\dot{\mathbf R}.
\]

这类问题通常需要结合势能面、非绝热耦合和经典轨迹。

\hypertarget{ux5bf9ux5e94ux6e90ux7801-22}{%
\subsubsection{对应源码}\label{ux5bf9ux5e94ux6e90ux7801-22}}

\begin{itemize}
\tightlist
\item
  \href{https://github.com/l1Ha/QuantumGeneralModule/blob/main/src/mod_surface_scattering.f90}{src/mod\_surface\_scattering.f90}
\item
  \href{https://github.com/l1Ha/QuantumGeneralModule/blob/main/src/mod_surface_reaction_er.f90}{src/mod\_surface\_reaction\_er.f90}
\item
  \href{https://github.com/l1Ha/QuantumGeneralModule/blob/main/src/mod_surface_electronic_friction.f90}{src/mod\_surface\_electronic\_friction.f90}
\end{itemize}

\hypertarget{ux76f8ux5bf9ux8bbaux539fux5b50ux4e0e-rixs}{%
\subsection{6.8　相对论原子与
RIXS}\label{ux76f8ux5bf9ux8bbaux539fux5b50ux4e0e-rixs}}

相对论径向 Dirac 方程为

\[
\left[
c\boldsymbol{\alpha}\cdot\mathbf p
+
\beta mc^2
+
V(r)
\right]
\psi
=
E\psi.
\]

它自然给出自旋轨道分裂和精细结构。RIXS 使用 Kramers--Heisenberg
二阶散射截面：

\[
I_{fi}(\omega)
\propto
\left|
\sum_n
\frac{
\langle f|\hat D^{\dagger}|n\rangle
\langle n|\hat D|i\rangle
}
{E_i+\hbar\omega-E_n+i\Gamma_n/2}
\right|^2.
\]

这类模块主要用于内壳层激发、共振非弹性散射和精细结构分析。

\hypertarget{ux5bf9ux5e94ux6e90ux7801-23}{%
\subsubsection{对应源码}\label{ux5bf9ux5e94ux6e90ux7801-23}}

\begin{itemize}
\tightlist
\item
  \href{https://github.com/l1Ha/QuantumGeneralModule/blob/main/src/mod_relativistic_atomic.f90}{src/mod\_relativistic\_atomic.f90}
\item
  \href{https://github.com/l1Ha/QuantumGeneralModule/blob/main/src/mod_resonant_xray_scattering.f90}{src/mod\_resonant\_xray\_scattering.f90}
\end{itemize}

\hypertarget{penning-ux4e0eux7f14ux5408ux7535ux79bb}{%
\subsection{6.9　Penning
与缔合电离}\label{penning-ux4e0eux7f14ux5408ux7535ux79bb}}

Penning 电离发生在亚稳态原子与中性原子碰撞时：

\[
A^{*}+B\rightarrow A+B^{+}+e^{-}.
\]

缔合电离则为

\[
A^{*}+B\rightarrow AB^{+}+e^{-}.
\]

GeneralModule 的 \texttt{mod\_penning\_associative\_ionization}
使用入口中性态 Morse 势、离子态 Morse 势和距离依赖自电离宽度
\(\Gamma(R)\)，并计算经典存活概率

\[
S(b)
=
\exp\left[
-2\int_{R_0}^{\infty}
\frac{\Gamma(R)}{\hbar v_R(R)}\,dR
\right].
\]

\hypertarget{ux5bf9ux5e94ux6e90ux7801-24}{%
\subsubsection{对应源码}\label{ux5bf9ux5e94ux6e90ux7801-24}}

\begin{itemize}
\tightlist
\item
  \href{https://github.com/l1Ha/QuantumGeneralModule/blob/main/src/mod_penning_associative_ionization.f90}{src/mod\_penning\_associative\_ionization.f90}
\end{itemize}

\hypertarget{ux5f3aux573aux7269ux7406ux4e0e-hhg}{%
\subsection{6.10　强场物理与
HHG}\label{ux5f3aux573aux7269ux7406ux4e0e-hhg}}

强场物理中常用 Keldysh 参数

\[
\gamma=\frac{\omega\sqrt{2I_p}}{E_0},
\]

和有质动力能

\[
U_p=\frac{E_0^2}{4\omega^2}.
\]

HHG 截止能量近似为

\[
E_{\mathrm{cut}}
=
I_p+3.17U_p.
\]

Corkum 提出三步模型【Corkum 1993】；Lewenstein
等给出强场近似下的偶极响应理论【Lewenstein et al.~1994】。

\hypertarget{ux5bf9ux5e94ux6e90ux7801-25}{%
\subsubsection{对应源码}\label{ux5bf9ux5e94ux6e90ux7801-25}}

\begin{itemize}
\tightlist
\item
  \href{https://github.com/l1Ha/QuantumGeneralModule/blob/main/src/mod_coulomb_atomic.f90}{src/mod\_coulomb\_atomic.f90}
\item
  \href{https://github.com/l1Ha/QuantumGeneralModule/blob/main/src/mod_hhg_spectra.f90}{src/mod\_hhg\_spectra.f90}
\item
  \href{https://github.com/l1Ha/QuantumGeneralModule/blob/main/src/mod_strong_field_nsdi.f90}{src/mod\_strong\_field\_nsdi.f90}
\end{itemize}

\hypertarget{ux91ccux5fb7ux5821ux963bux585eux4e0eux91cfux5b50ux591aux4f53}{%
\subsection{6.11　里德堡阻塞与量子多体}\label{ux91ccux5fb7ux5821ux963bux585eux4e0eux91cfux5b50ux591aux4f53}}

里德堡阻塞源于强偶极相互作用引起的能级移位：

\[
V(R)\sim\frac{C_6}{R^6}
\quad\text{或}\quad
V(R)\sim\frac{C_3}{R^3}.
\]

当相互作用能大于激发激光线宽时，一个里德堡激发会阻止邻近原子再被激发。该机制用于量子门、多体动力学和量子多体疤痕。

\hypertarget{ux5bf9ux5e94ux6e90ux7801-26}{%
\subsubsection{对应源码}\label{ux5bf9ux5e94ux6e90ux7801-26}}

\begin{itemize}
\tightlist
\item
  \href{https://github.com/l1Ha/QuantumGeneralModule/blob/main/src/mod_rydberg_blockade.f90}{src/mod\_rydberg\_blockade.f90}
\end{itemize}

\hypertarget{ux53c2ux8003ux6587ux732e-3}{%
\subsection{6.12　参考文献}\label{ux53c2ux8003ux6587ux732e-3}}

\begin{enumerate}
\def\labelenumi{\arabic{enumi}.}
\tightlist
\item
  E. Braaten and H.-W. Hammer, ``Universality in few-body systems with
  large scattering length'', \emph{Phys. Rep.} \textbf{428}, 259 (2006).
  DOI: 10.1016/j.physrep.2006.03.001.
\item
  M. Olshanii, ``Atomic scattering in the presence of an external
  confinement and a gas of hard-core bosons'', \emph{Phys. Rev.~Lett.}
  \textbf{81}, 938 (1998). DOI: 10.1103/PhysRevLett.81.938.
\item
  U. Fano, ``Effects of configuration interaction on intensities and
  phase shifts'', \emph{Phys. Rev.} \textbf{124}, 1866 (1961). DOI:
  10.1103/PhysRev.124.1866.
\item
  W. P. Reinhardt, ``Complex coordinates in the theory of atomic and
  molecular structure and dynamics'', \emph{Annu. Rev.~Phys. Chem.}
  \textbf{33}, 223 (1982). DOI: 10.1146/annurev.pc.33.100182.001255.
\item
  N. Moiseyev, ``Quantum theory of resonances: calculating energies,
  widths and cross-sections by complex scaling'', \emph{Phys. Rep.}
  \textbf{302}, 212 (1998). DOI: 10.1016/S0370-1573(98)00002-7.
\item
  H. C. Longuet-Higgins et al., ``Studies of the Jahn--Teller effect.
  II. The dynamical problem'', \emph{Proc. R. Soc. Lond. A}
  \textbf{244}, 1 (1958). DOI: 10.1098/rspa.1958.0022.
\item
  M. V. Berry, ``Quantal phase factors accompanying adiabatic changes'',
  \emph{Proc. R. Soc. Lond. A} \textbf{392}, 45 (1984). DOI:
  10.1098/rspa.1984.0023.
\item
  T.-L. Ho, ``Spinor Bose condensates in optical traps'', \emph{Phys.
  Rev.~Lett.} \textbf{81}, 742 (1998). DOI: 10.1103/PhysRevLett.81.742.
\item
  T. Ohmi and K. Machida, ``Bose--Einstein condensation with internal
  degrees of freedom in alkali atom gases'', \emph{J. Phys. Soc. Jpn.}
  \textbf{67}, 1822 (1998). DOI: 10.1143/JPSJ.67.1822.
\item
  P. B. Corkum, ``Plasma perspective on strong field multiphoton
  ionization'', \emph{Phys. Rev.~Lett.} \textbf{71}, 1994 (1993). DOI:
  10.1103/PhysRevLett.71.1994.
\item
  M. Lewenstein, P. Balcou, M. Y. Ivanov, A. L'Huillier, and P. B.
  Corkum, ``Theory of high-harmonic generation by low-frequency laser
  fields'', \emph{Phys. Rev.~A} \textbf{49}, 2117 (1994). DOI:
  10.1103/PhysRevA.49.2117.
\item
  GeneralModule
  源码与文档：\url{https://github.com/l1Ha/QuantumGeneralModule} \#
  第七部分　典型数值实验与代码映射
\end{enumerate}

\hypertarget{morse-ux52bfux80fdux9762ux7684-fgh-ux672cux5f81ux6001}{%
\subsection{7.1　Morse 势能面的 FGH
本征态}\label{morse-ux52bfux80fdux9762ux7684-fgh-ux672cux5f81ux6001}}

\hypertarget{ux7269ux7406ux6a21ux578b}{%
\subsubsection{物理模型}\label{ux7269ux7406ux6a21ux578b}}

双原子分子 Morse 势为

\[
V(R)=D_e\left[1-e^{-\beta(R-R_e)}\right]^2.
\]

一维径向哈密顿量为

\[
\hat H
=
-\frac{\hbar^2}{2\mu}
\frac{d^2}{dR^2}
+
V(R).
\]

在原子单位下，

\[
\hat H
=
-\frac{1}{2\mu}
\frac{d^2}{dR^2}
+
V(R).
\]

\hypertarget{ux6570ux503cux5b9eux9a8c-1}{%
\subsubsection{数值实验}\label{ux6570ux503cux5b9eux9a8c-1}}

\begin{enumerate}
\def\labelenumi{\arabic{enumi}.}
\tightlist
\item
  选择计算域 \([R_{\min},R_{\max}]\)，例如 \([0.8,6.0]\) Bohr。
\item
  选择网格数 \(N\)，例如 \(128,256,512,1024\)。
\item
  用 Sinc-DVR 构造动能矩阵。
\item
  在网格点上计算 \(V(R_i)\)。
\item
  组装 \(H_{ij}=T_{ij}+V(R_i)\delta_{ij}\)。
\item
  调用实对称本征求解器。
\item
  将本征向量除以 \(\sqrt{\Delta x}\) 得到连续归一化波函数。
\item
  将前若干能级与 Morse 解析式比较。
\end{enumerate}

\hypertarget{ux89e3ux6790ux80fdux7ea7}{%
\subsubsection{解析能级}\label{ux89e3ux6790ux80fdux7ea7}}

Morse 势的解析能级为

\[
E_v
=
\hbar\omega_e\left(v+\frac12\right)
-
\hbar\omega_ex_e\left(v+\frac12\right)^2,
\]

其中

\[
\omega_e=\beta\sqrt{\frac{2D_e}{\mu}},
\qquad
\omega_ex_e=\frac{\omega_e^2}{4D_e}.
\]

\hypertarget{ux9a8cux8bc1ux6307ux6807}{%
\subsubsection{验证指标}\label{ux9a8cux8bc1ux6307ux6807}}

\begin{itemize}
\tightlist
\item
  能级误差随 \(N\) 收敛；
\item
  波函数在边界处趋于零；
\item
  离散内积满足 \(\sum_i\Delta x\,|\psi_i|^2=1\)；
\item
  期望值 \(\langle R\rangle\) 和 \(\langle R^2\rangle\) 稳定。
\end{itemize}

\hypertarget{ux5bf9ux5e94ux6e90ux7801-27}{%
\subsubsection{对应源码}\label{ux5bf9ux5e94ux6e90ux7801-27}}

\begin{itemize}
\tightlist
\item
  \href{https://github.com/l1Ha/QuantumGeneralModule/blob/main/examples/ex01_fgh_diatomic_bound_states.f90}{examples/ex01\_fgh\_diatomic\_bound\_states.f90}
\item
  二维码：
\end{itemize}

\includegraphics[width=2cm,height=\textheight]{qr/examples__ex01_fgh_diatomic_bound_states.f90.png}

\begin{itemize}
\item
  \href{https://github.com/l1Ha/QuantumGeneralModule/blob/main/src/mod_dvr_grid.f90}{src/mod\_dvr\_grid.f90}
\item
  二维码：
\end{itemize}

\includegraphics[width=2cm,height=\textheight]{qr/src__mod_dvr_grid.f90.png}

\hypertarget{ux4e00ux7ef4ux6ce2ux5305ux7684-split-operator-ux4f20ux64ad}{%
\subsection{7.2　一维波包的 Split-Operator
传播}\label{ux4e00ux7ef4ux6ce2ux5305ux7684-split-operator-ux4f20ux64ad}}

\hypertarget{ux7269ux7406ux6a21ux578b-1}{%
\subsubsection{物理模型}\label{ux7269ux7406ux6a21ux578b-1}}

考虑自由粒子或一维势阱中的高斯波包：

\[
\psi(x,0)
=
\left(\frac{1}{2\pi\sigma^2}\right)^{1/4}
\exp\left[
-\frac{(x-x_0)^2}{4\sigma^2}
+ik_0x
\right].
\]

无外场时，

\[
\langle x\rangle(t)=x_0+\frac{\hbar k_0}{m}t.
\]

\hypertarget{ux6570ux503cux5b9eux9a8c-2}{%
\subsubsection{数值实验}\label{ux6570ux503cux5b9eux9a8c-2}}

\begin{enumerate}
\def\labelenumi{\arabic{enumi}.}
\tightlist
\item
  在均匀网格上初始化高斯波包。
\item
  设置势能数组 \(V(x)\)。
\item
  使用 Split-Operator 传播多个时间步。
\item
  每步后计算
  \(\langle x\rangle\)、\(\langle p\rangle\)、\(\langle H\rangle\) 和
  \(\|\psi\|\)。
\item
  改变 \(\Delta t\) 和 \(N\)，绘制误差曲线。
\end{enumerate}

\hypertarget{ux9a8cux8bc1ux6307ux6807-1}{%
\subsubsection{验证指标}\label{ux9a8cux8bc1ux6307ux6807-1}}

\begin{itemize}
\tightlist
\item
  自由粒子时 \(\langle x\rangle(t)\) 线性；
\item
  动量 \(\langle p\rangle\) 守恒；
\item
  能量 \(\langle H\rangle\) 守恒；
\item
  范数 \(\|\psi\|\) 守恒；
\item
  边界处波函数无明显反射。
\end{itemize}

\hypertarget{ux5bf9ux5e94ux6e90ux7801-28}{%
\subsubsection{对应源码}\label{ux5bf9ux5e94ux6e90ux7801-28}}

\begin{itemize}
\tightlist
\item
  \href{https://github.com/l1Ha/QuantumGeneralModule/blob/main/examples/ex03_split_operator_1d.f90}{examples/ex03\_split\_operator\_1d.f90}
\item
  二维码：
\end{itemize}

\includegraphics[width=2cm,height=\textheight]{qr/examples__ex03_split_operator_1d.f90.png}

\begin{itemize}
\item
  \href{https://github.com/l1Ha/QuantumGeneralModule/blob/main/src/mod_wavepacket_propagator.f90}{src/mod\_wavepacket\_propagator.f90}
\item
  二维码：
\end{itemize}

\includegraphics[width=2cm,height=\textheight]{qr/src__mod_wavepacket_propagator.f90.png}

\hypertarget{ux5b9aux6001ux4e0eux542bux65f6ux6563ux5c04ux7684ux4ea4ux53c9ux9a8cux8bc1}{%
\subsection{7.3　定态与含时散射的交叉验证}\label{ux5b9aux6001ux4e0eux542bux65f6ux6563ux5c04ux7684ux4ea4ux53c9ux9a8cux8bc1}}

\hypertarget{ux7269ux7406ux6a21ux578b-2}{%
\subsubsection{物理模型}\label{ux7269ux7406ux6a21ux578b-2}}

一维散射问题的定态形式为

\[
-\frac{\hbar^2}{2m}\psi''(x)+V(x)\psi(x)=E\psi(x).
\]

含时形式为

\[
i\hbar\frac{\partial\psi}{\partial t}
=
-\frac{\hbar^2}{2m}\psi''(x)+V(x)\psi(x).
\]

\hypertarget{ux6570ux503cux5b9eux9a8c-3}{%
\subsubsection{数值实验}\label{ux6570ux503cux5b9eux9a8c-3}}

\begin{enumerate}
\def\labelenumi{\arabic{enumi}.}
\tightlist
\item
  在定态方法中扫描能量 \(E\)，得到透射率 \(T(E)\)。
\item
  在含时方法中用高斯波包入射，通过通量积分得到 \(T(E)\)。
\item
  比较共振峰位置、宽度和幅度。
\item
  改变波包宽度 \(\sigma\)，检查能量分辨率。
\item
  改变 CAP 长度和强度，检查吸收误差。
\end{enumerate}

\hypertarget{ux9a8cux8bc1ux6307ux6807-2}{%
\subsubsection{验证指标}\label{ux9a8cux8bc1ux6307ux6807-2}}

\begin{itemize}
\tightlist
\item
  峰位一致；
\item
  峰宽一致；
\item
  透射与反射之和接近 1；
\item
  不同 \(\sigma\) 下结果稳定。
\end{itemize}

\hypertarget{ux5bf9ux5e94ux6e90ux7801-29}{%
\subsubsection{对应源码}\label{ux5bf9ux5e94ux6e90ux7801-29}}

\begin{itemize}
\tightlist
\item
  \href{https://github.com/l1Ha/QuantumGeneralModule/blob/main/examples/ex07_scattering_wavefunctions_ti_td.f90}{examples/ex07\_scattering\_wavefunctions\_ti\_td.f90}
\item
  二维码：
\end{itemize}

\includegraphics[width=2cm,height=\textheight]{qr/examples__ex07_scattering_wavefunctions_ti_td.f90.png}

\begin{itemize}
\item
  \href{https://github.com/l1Ha/QuantumGeneralModule/blob/main/src/mod_td_scattering.f90}{src/mod\_td\_scattering.f90}
\item
  二维码：
\end{itemize}

\includegraphics[width=2cm,height=\textheight]{qr/src__mod_td_scattering.f90.png}

\hypertarget{ux591aux901aux9053-feshbach-ux626bux63cf}{%
\subsection{7.4　多通道 Feshbach
扫描}\label{ux591aux901aux9053-feshbach-ux626bux63cf}}

\hypertarget{ux7269ux7406ux6a21ux578b-3}{%
\subsubsection{物理模型}\label{ux7269ux7406ux6a21ux578b-3}}

多通道散射矩阵由开通道 K 矩阵构造：

\[
\mathbf S=(\mathbf I+i\mathbf K)(\mathbf I-i\mathbf K)^{-1}.
\]

Feshbach 共振的散射长度随磁场变化：

\[
a_s(B)
=
a_{\mathrm{bg}}
\left[
1-\frac{\Delta B}{B-B_0}
\right].
\]

\hypertarget{ux6570ux503cux5b9eux9a8c-4}{%
\subsubsection{数值实验}\label{ux6570ux503cux5b9eux9a8c-4}}

\begin{enumerate}
\def\labelenumi{\arabic{enumi}.}
\tightlist
\item
  选择磁场范围 \([B_{\min},B_{\max}]\)。
\item
  对每个 \(B\) 构造耦合势矩阵。
\item
  调用多通道 Log-Derivative 求解器。
\item
  提取入射道散射长度 \(a_s(B)\)。
\item
  拟合 \(a_{\mathrm{bg}}\)、\(\Delta B\) 和 \(B_0\)。
\item
  改变磁场步长和径向网格，检查拟合参数稳定。
\end{enumerate}

\hypertarget{ux9a8cux8bc1ux6307ux6807-3}{%
\subsubsection{验证指标}\label{ux9a8cux8bc1ux6307ux6807-3}}

\begin{itemize}
\tightlist
\item
  \(S\) 矩阵幺正；
\item
  开通道与闭通道划分正确；
\item
  共振位置对网格不敏感；
\item
  拟合残差可控。
\end{itemize}

\hypertarget{ux5bf9ux5e94ux6e90ux7801-30}{%
\subsubsection{对应源码}\label{ux5bf9ux5e94ux6e90ux7801-30}}

\begin{itemize}
\tightlist
\item
  \href{https://github.com/l1Ha/QuantumGeneralModule/blob/main/examples/ex08_ultracold_feshbach_segmented.f90}{examples/ex08\_ultracold\_feshbach\_segmented.f90}
\item
  二维码：
\end{itemize}

\includegraphics[width=2cm,height=\textheight]{qr/examples__ex08_ultracold_feshbach_segmented.f90.png}

\begin{itemize}
\item
  \href{https://github.com/l1Ha/QuantumGeneralModule/blob/main/src/mod_field_scattering.f90}{src/mod\_field\_scattering.f90}
\item
  二维码：
\end{itemize}

\includegraphics[width=2cm,height=\textheight]{qr/src__mod_field_scattering.f90.png}

\hypertarget{lindblad-ux5f1bux8c6bux4e0eux9000ux76f8ux4f4d}{%
\subsection{7.5　Lindblad
弛豫与退相位}\label{lindblad-ux5f1bux8c6bux4e0eux9000ux76f8ux4f4d}}

\hypertarget{ux7269ux7406ux6a21ux578b-4}{%
\subsubsection{物理模型}\label{ux7269ux7406ux6a21ux578b-4}}

二能级 Lindblad 主方程为

\[
\dot\rho
=
-\frac{i}{\hbar}[H,\rho]
+
\gamma\left(
\sigma_-\rho\sigma_+
-
\frac12\{\sigma_+\sigma_-,\rho\}
\right)
+
\gamma_\phi
\left(
\sigma_z\rho\sigma_z
-
\rho
\right).
\]

\hypertarget{ux6570ux503cux5b9eux9a8c-5}{%
\subsubsection{数值实验}\label{ux6570ux503cux5b9eux9a8c-5}}

\begin{enumerate}
\def\labelenumi{\arabic{enumi}.}
\tightlist
\item
  初始化激发态 \(\rho_{ee}(0)=1\)。
\item
  设置弛豫率 \(\gamma\) 和退相位率 \(\gamma_\phi\)。
\item
  用 RK4 传播若干时间步。
\item
  记录
  \(\rho_{ee}(t)\)、\(\rho_{eg}(t)\)、\(\operatorname{Tr}\rho\)、纯度
  \(P\)。
\item
  改变 \(\Delta t\)，检查收敛。
\end{enumerate}

\hypertarget{ux9a8cux8bc1ux6307ux6807-4}{%
\subsubsection{验证指标}\label{ux9a8cux8bc1ux6307ux6807-4}}

\begin{itemize}
\tightlist
\item
  \(\rho_{ee}(t)\) 指数衰减；
\item
  \(|\rho_{eg}(t)|\) 按 \(e^{-(\gamma/2+\gamma_\phi)t}\) 衰减；
\item
  迹守恒；
\item
  纯度单调下降。
\end{itemize}

\hypertarget{ux5bf9ux5e94ux6e90ux7801-31}{%
\subsubsection{对应源码}\label{ux5bf9ux5e94ux6e90ux7801-31}}

\begin{itemize}
\tightlist
\item
  \href{https://github.com/l1Ha/QuantumGeneralModule/blob/main/tests/test_open_quantum_opt.f90}{tests/test\_open\_quantum\_opt.f90}
\item
  二维码：
\end{itemize}

\includegraphics[width=2cm,height=\textheight]{qr/tests__test_open_quantum_opt.f90.png}

\begin{itemize}
\item
  \href{https://github.com/l1Ha/QuantumGeneralModule/blob/main/src/mod_open_quantum.f90}{src/mod\_open\_quantum.f90}
\item
  二维码：
\end{itemize}

\includegraphics[width=2cm,height=\textheight]{qr/src__mod_open_quantum.f90.png}

\hypertarget{tully-fssh-ux5206ux652fux6bd4}{%
\subsection{7.6　Tully FSSH
分支比}\label{tully-fssh-ux5206ux652fux6bd4}}

\hypertarget{ux7269ux7406ux6a21ux578b-5}{%
\subsubsection{物理模型}\label{ux7269ux7406ux6a21ux578b-5}}

FSSH 在活性态 \(j\) 上向态 \(k\) 的跳跃概率为

\[
g_{j\rightarrow k}
=
\max
\left[
0,
-\frac{2\Delta t}{|c_j|^2}
\operatorname{Re}
\left(
c_j^{*}c_k\,\mathbf v\cdot\mathbf d_{jk}
\right)
\right].
\]

\hypertarget{ux6570ux503cux5b9eux9a8c-6}{%
\subsubsection{数值实验}\label{ux6570ux503cux5b9eux9a8c-6}}

\begin{enumerate}
\def\labelenumi{\arabic{enumi}.}
\tightlist
\item
  初始化核坐标和动量。
\item
  选择初始电子态。
\item
  对每条轨迹使用 Velocity-Verlet 推进核。
\item
  使用 FSSH 规则判断跳跃。
\item
  统计透射和反射分支比。
\item
  改变轨迹数 \(N_{\mathrm{traj}}\)，估计统计误差。
\end{enumerate}

\hypertarget{ux9a8cux8bc1ux6307ux6807-5}{%
\subsubsection{验证指标}\label{ux9a8cux8bc1ux6307ux6807-5}}

\begin{itemize}
\tightlist
\item
  分支比随 \(N_{\mathrm{traj}}\) 收敛；
\item
  能量守恒；
\item
  frustrated hop 比例合理；
\item
  不同随机种子结果稳定。
\end{itemize}

\hypertarget{ux5bf9ux5e94ux6e90ux7801-32}{%
\subsubsection{对应源码}\label{ux5bf9ux5e94ux6e90ux7801-32}}

\begin{itemize}
\tightlist
\item
  \href{https://github.com/l1Ha/QuantumGeneralModule/blob/main/examples/ex30_tully_surface_hopping.f90}{examples/ex30\_tully\_surface\_hopping.f90}
\item
  二维码：
\end{itemize}

\includegraphics[width=2cm,height=\textheight]{qr/examples__ex30_tully_surface_hopping.f90.png}

\begin{itemize}
\item
  \href{https://github.com/l1Ha/QuantumGeneralModule/blob/main/src/mod_surface_hopping_fssh.f90}{src/mod\_surface\_hopping\_fssh.f90}
\item
  二维码：
\end{itemize}

\includegraphics[width=2cm,height=\textheight]{qr/src__mod_surface_hopping_fssh.f90.png}

\hypertarget{python-ux53efux89c6ux5316ux4e0eux4ea4ux53c9ux9a8cux8bc1}{%
\subsection{7.7　Python
可视化与交叉验证}\label{python-ux53efux89c6ux5316ux4e0eux4ea4ux53c9ux9a8cux8bc1}}

Fortran 输出通常为文本或二进制数据。Python 可用于：

\begin{enumerate}
\def\labelenumi{\arabic{enumi}.}
\tightlist
\item
  读取数据；
\item
  独立复算关键量；
\item
  绘制能级、波函数、截面和误差曲线；
\item
  生成发表级图像。
\end{enumerate}

推荐使用 NumPy 处理数组，Matplotlib 绘图【Harris et al.~2020；Hunter
2007】。

\hypertarget{ux5bf9ux5e94ux6e90ux7801-33}{%
\subsubsection{对应源码}\label{ux5bf9ux5e94ux6e90ux7801-33}}

\begin{itemize}
\tightlist
\item
  \href{https://github.com/l1Ha/QuantumGeneralModule/blob/main/python/pygenmod/visualizer.py}{python/pygenmod/visualizer.py}
\item
  二维码：
\end{itemize}

\includegraphics[width=2cm,height=\textheight]{qr/python__pygenmod__visualizer.py.png}

\hypertarget{ux5b9eux9a8cux8bb0ux5f55ux6a21ux677f}{%
\subsection{7.8　实验记录模板}\label{ux5b9eux9a8cux8bb0ux5f55ux6a21ux677f}}

\begin{longtable}[]{@{}ll@{}}
\toprule
字段 & 建议内容\tabularnewline
\midrule
\endhead
实验名称 & 物理问题和目标\tabularnewline
源码版本 & commit、tag 或文件时间\tabularnewline
编译器 & gfortran/ifx/flang 版本\tabularnewline
系统 & OS、CPU、内存\tabularnewline
参数 & 单位、网格、初态、时间步\tabularnewline
随机性 & 种子或确定性声明\tabularnewline
理论基准 & 解析解、文献值、守恒量\tabularnewline
输出 & 能量、截面、范数、分支比\tabularnewline
收敛判据 & 相对误差和阈值\tabularnewline
复现命令 & 完整脚本\tabularnewline
\bottomrule
\end{longtable}

\hypertarget{ux7b2cux516bux90e8ux5206-ux6e90ux7801ux9605ux8bfbux6a21ux5757ux6269ux5c55ux4e0eux5de5ux7a0bux5b9eux8df5}{%
\section{第八部分　源码阅读、模块扩展与工程实践}\label{ux7b2cux516bux90e8ux5206-ux6e90ux7801ux9605ux8bfbux6a21ux5757ux6269ux5c55ux4e0eux5de5ux7a0bux5b9eux8df5}}

\hypertarget{ux9605ux8bfb-generalmodule-ux6e90ux7801ux7684ux987aux5e8f}{%
\subsection{8.1　阅读 GeneralModule
源码的顺序}\label{ux9605ux8bfb-generalmodule-ux6e90ux7801ux7684ux987aux5e8f}}

阅读大型科学计算库时，建议按以下顺序：

\begin{enumerate}
\def\labelenumi{\arabic{enumi}.}
\tightlist
\item
  先看 \texttt{README.md}，确认项目定位、模块清单和示例结构。
\item
  再看 \texttt{LITERATURE.md}，了解每个模块对应的理论来源。
\item
  进入 \texttt{src/mod\_constants.f90}，掌握类型别名和单位转换。
\item
  阅读 \texttt{src/mod\_linear\_algebra.f90}，理解本征、求逆和 FFT。
\item
  阅读 \texttt{src/mod\_dvr\_grid.f90}，掌握 DVR 与 FGH。
\item
  阅读 \texttt{src/mod\_wavepacket\_propagator.f90}，理解
  Split-Operator、RK4、ABM4 和 Bloch 方程。
\item
  进入散射和专题模块，结合测试与示例阅读。
\item
  最后阅读 \texttt{general\_module.f90}，了解公共接口如何聚合。
\end{enumerate}

不要从物理专题模块直接入手。若缺少底层数学和表示方法的基础，很容易把实现细节误解为物理结构。

\begin{itemize}
\tightlist
\item
  README：
\end{itemize}

\href{https://github.com/l1Ha/QuantumGeneralModule/blob/main/README.md}{README.md}

\begin{itemize}
\tightlist
\item
  二维码：
\end{itemize}

\includegraphics[width=2cm,height=\textheight]{qr/README.md.png}

\begin{itemize}
\tightlist
\item
  LITERATURE：
\end{itemize}

\href{https://github.com/l1Ha/QuantumGeneralModule/blob/main/LITERATURE.md}{LITERATURE.md}

\begin{itemize}
\tightlist
\item
  二维码：
\end{itemize}

\includegraphics[width=2cm,height=\textheight]{qr/LITERATURE.md.png}

\hypertarget{ux6a21ux5757ux6a21ux677fux4e0eux63a5ux53e3ux8bbeux8ba1}{%
\subsection{8.2　模块模板与接口设计}\label{ux6a21ux5757ux6a21ux677fux4e0eux63a5ux53e3ux8bbeux8ba1}}

一个新模块应遵循以下结构：

\begin{Shaded}
\begin{Highlighting}[]
\KeywordTok{module}\NormalTok{ mod\_example}
    \KeywordTok{use}\NormalTok{, }\KeywordTok{intrinsic} \DataTypeTok{::}\NormalTok{ iso\_fortran\_env, }\KeywordTok{only}\NormalTok{: dp }\KeywordTok{=}\OperatorTok{\textgreater{}}\NormalTok{ real64}
    \KeywordTok{implicit} \KeywordTok{none}
    \DataTypeTok{private}

    \DataTypeTok{public} \DataTypeTok{::}\NormalTok{ example\_config\_t}
    \DataTypeTok{public} \DataTypeTok{::}\NormalTok{ example\_result\_t}
    \DataTypeTok{public} \DataTypeTok{::}\NormalTok{ solve\_example}

    \DataTypeTok{type} \DataTypeTok{::}\NormalTok{ example\_config\_t}
\DataTypeTok{        real(dp)} \DataTypeTok{::}\NormalTok{ mass }\KeywordTok{=} \FloatTok{1.0\_dp}
\DataTypeTok{        real(dp)} \DataTypeTok{::}\NormalTok{ energy }\KeywordTok{=} \FloatTok{0.0\_dp}
        \DataTypeTok{integer}  \DataTypeTok{::}\NormalTok{ n\_grid }\KeywordTok{=} \DecValTok{256}
    \DataTypeTok{end type}\NormalTok{ example\_config\_t}

    \DataTypeTok{type} \DataTypeTok{::}\NormalTok{ example\_result\_t}
\DataTypeTok{        real(dp)}\NormalTok{, }\DataTypeTok{allocatable} \DataTypeTok{::}\NormalTok{ energies(:)}
\DataTypeTok{        real(dp)}\NormalTok{, }\DataTypeTok{allocatable} \DataTypeTok{::}\NormalTok{ wavefunctions(:, :)}
        \DataTypeTok{integer} \DataTypeTok{::}\NormalTok{ stat }\KeywordTok{=} \DecValTok{0}
    \DataTypeTok{end type}\NormalTok{ example\_result\_t}

\CommentTok{contains}

    \KeywordTok{subroutine}\NormalTok{ solve\_example(cfg, res)}
        \DataTypeTok{type(example\_config\_t)}\NormalTok{, }\DataTypeTok{intent(in)}  \DataTypeTok{::}\NormalTok{ cfg}
        \DataTypeTok{type(example\_result\_t)}\NormalTok{, }\DataTypeTok{intent(out)} \DataTypeTok{::}\NormalTok{ res}

\NormalTok{        res\%stat }\KeywordTok{=} \DecValTok{0}
        \KeywordTok{if}\NormalTok{ (cfg\%n\_grid }\OperatorTok{\textless{}} \DecValTok{8}\NormalTok{) }\KeywordTok{then}
\NormalTok{            res\%stat }\KeywordTok{=} \KeywordTok{{-}}\DecValTok{1}
            \KeywordTok{return}
        \KeywordTok{end if}

        \KeywordTok{allocate}\NormalTok{(res\%energies(cfg\%n\_grid))}
        \KeywordTok{allocate}\NormalTok{(res\%wavefunctions(cfg\%n\_grid, cfg\%n\_grid))}
    \KeywordTok{end subroutine}\NormalTok{ solve\_example}

\KeywordTok{end module}\NormalTok{ mod\_example}
\end{Highlighting}
\end{Shaded}

关键要求：

\begin{itemize}
\tightlist
\item
  不要在库内部直接写 \texttt{stop}；
\item
  不要把调试输出写死在核心算法中；
\item
  不要让模块依赖 \texttt{general\_module}；
\item
  不要用全局变量传递状态；
\item
  不要在接口中隐式转换单位；
\item
  所有可能失败的分配或求解都应返回状态码。
\end{itemize}

\hypertarget{ux6570ux503cux7b97ux6cd5ux7684ux9a8cux8bc1ux6a21ux677f}{%
\subsection{8.3　数值算法的验证模板}\label{ux6570ux503cux7b97ux6cd5ux7684ux9a8cux8bc1ux6a21ux677f}}

每个新算法至少应准备三类测试：

\begin{enumerate}
\def\labelenumi{\arabic{enumi}.}
\tightlist
\item
  解析极限；
\item
  守恒量；
\item
  收敛阶。
\end{enumerate}

例如对一个新 DVR，应验证：

\[
\sum_i w_i\psi_i^{(k)}\psi_i^{(l)}=\delta_{kl},
\]

\[
E_N-E
=
CN^{-p}+O(N^{-(p+1)}),
\]

以及对称矩阵的厄米性

\[
H=H^T.
\]

对传播器，应验证：

\[
\|\psi(t+\Delta t)\|_2=\|\psi(t)\|_2,
\]

\[
\langle H\rangle(t)=\langle H\rangle(0).
\]

对散射，应验证：

\[
\sigma_{\mathrm{tot}}=\frac{4\pi}{k}\operatorname{Im}f(0).
\]

\hypertarget{ux6d4bux8bd5ux6a21ux677f}{%
\subsection{8.4　测试模板}\label{ux6d4bux8bd5ux6a21ux677f}}

一个 Fortran 测试程序应包含：

\begin{itemize}
\tightlist
\item
  测试计数器；
\item
  断言函数；
\item
  明确的失败输出；
\item
  返回码；
\item
  可重复运行的参数。
\end{itemize}

\begin{Shaded}
\begin{Highlighting}[]
\KeywordTok{program}\NormalTok{ test\_example}
    \KeywordTok{use}\NormalTok{, }\KeywordTok{intrinsic} \DataTypeTok{::}\NormalTok{ iso\_fortran\_env, }\KeywordTok{only}\NormalTok{: dp }\KeywordTok{=}\OperatorTok{\textgreater{}}\NormalTok{ real64}
    \KeywordTok{implicit} \KeywordTok{none}

    \DataTypeTok{integer} \DataTypeTok{::}\NormalTok{ n\_pass, n\_total}
\DataTypeTok{    real(dp)} \DataTypeTok{::}\NormalTok{ x, y}

\NormalTok{    n\_pass }\KeywordTok{=} \DecValTok{0}
\NormalTok{    n\_total }\KeywordTok{=} \DecValTok{0}

\NormalTok{    x }\KeywordTok{=} \FloatTok{1.0\_dp}
\NormalTok{    y }\KeywordTok{=}\NormalTok{ x}\KeywordTok{*}\NormalTok{x}
\NormalTok{    n\_total }\KeywordTok{=}\NormalTok{ n\_total }\KeywordTok{+} \DecValTok{1}
    \KeywordTok{if}\NormalTok{ (}\BuiltInTok{abs}\NormalTok{(y }\KeywordTok{{-}} \FloatTok{1.0\_dp}\NormalTok{) }\OperatorTok{\textless{}} \FloatTok{1.0e{-}12\_dp}\NormalTok{) }\KeywordTok{then}
\NormalTok{        n\_pass }\KeywordTok{=}\NormalTok{ n\_pass }\KeywordTok{+} \DecValTok{1}
    \KeywordTok{else}
        \FunctionTok{print} \KeywordTok{*}\NormalTok{, }\StringTok{\textquotesingle{}FAIL: square\textquotesingle{}}
    \KeywordTok{end if}

    \FunctionTok{print} \StringTok{\textquotesingle{}(A,I0,A,I0,A)\textquotesingle{}}\NormalTok{, }\StringTok{\textquotesingle{}Passed \textquotesingle{}}\NormalTok{, n\_pass, }\StringTok{\textquotesingle{}/\textquotesingle{}}\NormalTok{, n\_total, }\StringTok{\textquotesingle{} tests.\textquotesingle{}}
    \KeywordTok{if}\NormalTok{ (n\_pass }\OperatorTok{/=}\NormalTok{ n\_total) }\KeywordTok{stop} \DecValTok{1}

\KeywordTok{end program}\NormalTok{ test\_example}
\end{Highlighting}
\end{Shaded}

测试应避免只检查``程序能运行''。应检查具体数值、守恒量和误差阶数。

\hypertarget{ci-ux4e2dux7684ux68c0ux67e5ux5c42ux6b21}{%
\subsection{8.5　CI
中的检查层次}\label{ci-ux4e2dux7684ux68c0ux67e5ux5c42ux6b21}}

持续集成应至少包含：

\begin{longtable}[]{@{}ll@{}}
\toprule
层次 & 内容\tabularnewline
\midrule
\endhead
编译 & 静态库、共享库、可执行文件\tabularnewline
单元测试 & 数学函数、本征、DVR、传播\tabularnewline
物理示例 & 输出存在且非空\tabularnewline
Python 测试 & NumPy 交叉验证\tabularnewline
文档 & README、API、公式可读\tabularnewline
复现 & 关键参数和随机种子\tabularnewline
\bottomrule
\end{longtable}

GitHub Actions 配置位于：

\href{https://github.com/l1Ha/QuantumGeneralModule/blob/main/.github/workflows/ci.yml}{.github/workflows/ci.yml}

\hypertarget{ux5982ux4f55ux6269ux5c55ux65b0ux7269ux7406ux6a21ux5757}{%
\subsection{8.6　如何扩展新物理模块}\label{ux5982ux4f55ux6269ux5c55ux65b0ux7269ux7406ux6a21ux5757}}

推荐步骤：

\begin{enumerate}
\def\labelenumi{\arabic{enumi}.}
\tightlist
\item
  明确物理模型和哈密顿量；
\item
  明确坐标、表象和边界条件；
\item
  写出核心公式和近似；
\item
  搜索文献并确认理论来源；
\item
  设计配置类型和结果类型；
\item
  实现最小可用算法；
\item
  编写解析极限测试；
\item
  加入收敛实验；
\item
  加入 Python 交叉验证；
\item
  更新 README、LITERATURE 和示例索引。
\end{enumerate}

\hypertarget{ux6587ux6863ux4e0eux5f15ux7528ux89c4ux8303}{%
\subsection{8.7　文档与引用规范}\label{ux6587ux6863ux4e0eux5f15ux7528ux89c4ux8303}}

每个模块应说明：

\begin{itemize}
\tightlist
\item
  物理问题；
\item
  关键假设；
\item
  输入输出；
\item
  单位制；
\item
  算法名称；
\item
  文献来源；
\item
  测试位置；
\item
  示例位置；
\item
  已知限制。
\end{itemize}

引用外部文献时，应给出作者、题目、期刊、卷、页码和
DOI。若引用开源软件，应给出项目地址和许可证。若未来引入第三方代码，必须在源码和变更记录中明确说明。

\hypertarget{ux7248ux672cux4e0eux53efux590dux73b0ux6027}{%
\subsection{8.8　版本与可复现性}\label{ux7248ux672cux4e0eux53efux590dux73b0ux6027}}

推荐使用语义化版本：

\begin{itemize}
\tightlist
\item
  \texttt{MAJOR}：接口不兼容变更；
\item
  \texttt{MINOR}：新增功能或模块；
\item
  \texttt{PATCH}：修复错误或改进文档。
\end{itemize}

每次重要数值变更都应记录：

\begin{enumerate}
\def\labelenumi{\arabic{enumi}.}
\tightlist
\item
  修改了哪些算法；
\item
  对哪些测试有影响；
\item
  是否改变默认精度；
\item
  是否改变输出格式；
\item
  是否需要重新标定物理参数。
\end{enumerate}

\hypertarget{ux7b2cux4e5dux90e8ux5206-ux8defux7ebfux56fe}{%
\section{第九部分　路线图}\label{ux7b2cux4e5dux90e8ux5206-ux8defux7ebfux56fe}}

\hypertarget{ux7b97ux6cd5ux65b9ux5411}{%
\subsection{9.1　算法方向}\label{ux7b97ux6cd5ux65b9ux5411}}

\begin{itemize}
\tightlist
\item
  更多正交多项式 DVR；
\item
  稀疏 Krylov 传播；
\item
  多维张量分解；
\item
  高阶分裂算符；
\item
  自适应时间步长；
\item
  并行散射计算。
\end{itemize}

\hypertarget{ux7269ux7406ux65b9ux5411}{%
\subsection{9.2　物理方向}\label{ux7269ux7406ux65b9ux5411}}

\begin{itemize}
\tightlist
\item
  更多表面反应模型；
\item
  非绝热耦合场；
\item
  开放系统谱方法；
\item
  高精度相对论结构；
\item
  多通道 Feshbach 系统；
\item
  少体与多体扩展。
\end{itemize}

\hypertarget{ux5de5ux7a0bux65b9ux5411}{%
\subsection{9.3　工程方向}\label{ux5de5ux7a0bux65b9ux5411}}

\begin{itemize}
\tightlist
\item
  统一测试断言库；
\item
  更多 CI 平台；
\item
  自动生成 API 文档；
\item
  自动化收敛报告；
\item
  数据格式标准化；
\item
  示例输出归档。
\end{itemize}

\hypertarget{ux6587ux6863ux65b9ux5411}{%
\subsection{9.4　文档方向}\label{ux6587ux6863ux65b9ux5411}}

\begin{itemize}
\tightlist
\item
  每个模块增加理论说明；
\item
  每个示例增加收敛数据；
\item
  将二维码索引扩展到全部源码；
\item
  增加中英双语 API 说明；
\item
  增加可复现实验清单。
\end{itemize}

\hypertarget{ux7b2cux5341ux90e8ux5206-ux6e90ux7801ux7d22ux5f15ux4e0eux4e8cux7ef4ux7801}{%
\section{第十部分　源码索引与二维码}\label{ux7b2cux5341ux90e8ux5206-ux6e90ux7801ux7d22ux5f15ux4e0eux4e8cux7ef4ux7801}}

以下二维码均指向 GeneralModule 仓库中的对应文件。二维码图像保存在
\texttt{docs/qr/}，文件名规则是把源码路径中的 \texttt{/} 替换为
\texttt{\_\_}。

\begin{longtable}[]{@{}lll@{}}
\toprule
\begin{minipage}[b]{0.30\columnwidth}\raggedright
内容\strut
\end{minipage} & \begin{minipage}[b]{0.30\columnwidth}\raggedright
路径\strut
\end{minipage} & \begin{minipage}[b]{0.30\columnwidth}\raggedright
二维码\strut
\end{minipage}\tabularnewline
\midrule
\endhead
\begin{minipage}[t]{0.30\columnwidth}\raggedright
仓库主页\strut
\end{minipage} & \begin{minipage}[t]{0.30\columnwidth}\raggedright
\url{https://github.com/l1Ha/QuantumGeneralModule}\strut
\end{minipage} & \begin{minipage}[t]{0.30\columnwidth}\raggedright
\texttt{docs/qr/repo\_root.png}\strut
\end{minipage}\tabularnewline
\begin{minipage}[t]{0.30\columnwidth}\raggedright
README\strut
\end{minipage} & \begin{minipage}[t]{0.30\columnwidth}\raggedright
\href{https://github.com/l1Ha/QuantumGeneralModule/blob/main/README.md}{README.md}\strut
\end{minipage} & \begin{minipage}[t]{0.30\columnwidth}\raggedright
\texttt{docs/qr/README.md.png}\strut
\end{minipage}\tabularnewline
\begin{minipage}[t]{0.30\columnwidth}\raggedright
配置指南\strut
\end{minipage} & \begin{minipage}[t]{0.30\columnwidth}\raggedright
\href{https://github.com/l1Ha/QuantumGeneralModule/blob/main/CONFIG_GUIDE.md}{CONFIG\_GUIDE.md}\strut
\end{minipage} & \begin{minipage}[t]{0.30\columnwidth}\raggedright
\texttt{docs/qr/CONFIG\_GUIDE.md.png}\strut
\end{minipage}\tabularnewline
\begin{minipage}[t]{0.30\columnwidth}\raggedright
文献映射\strut
\end{minipage} & \begin{minipage}[t]{0.30\columnwidth}\raggedright
\href{https://github.com/l1Ha/QuantumGeneralModule/blob/main/LITERATURE.md}{LITERATURE.md}\strut
\end{minipage} & \begin{minipage}[t]{0.30\columnwidth}\raggedright
\texttt{docs/qr/LITERATURE.md.png}\strut
\end{minipage}\tabularnewline
\begin{minipage}[t]{0.30\columnwidth}\raggedright
常数与单位\strut
\end{minipage} & \begin{minipage}[t]{0.30\columnwidth}\raggedright
\href{https://github.com/l1Ha/QuantumGeneralModule/blob/main/src/mod_constants.f90}{src/mod\_constants.f90}\strut
\end{minipage} & \begin{minipage}[t]{0.30\columnwidth}\raggedright
\texttt{docs/qr/src\_\_mod\_constants.f90.png}\strut
\end{minipage}\tabularnewline
\begin{minipage}[t]{0.30\columnwidth}\raggedright
特殊函数\strut
\end{minipage} & \begin{minipage}[t]{0.30\columnwidth}\raggedright
\href{https://github.com/l1Ha/QuantumGeneralModule/blob/main/src/mod_special_functions.f90}{src/mod\_special\_functions.f90}\strut
\end{minipage} & \begin{minipage}[t]{0.30\columnwidth}\raggedright
\texttt{docs/qr/src\_\_mod\_special\_functions.f90.png}\strut
\end{minipage}\tabularnewline
\begin{minipage}[t]{0.30\columnwidth}\raggedright
线性代数\strut
\end{minipage} & \begin{minipage}[t]{0.30\columnwidth}\raggedright
\href{https://github.com/l1Ha/QuantumGeneralModule/blob/main/src/mod_linear_algebra.f90}{src/mod\_linear\_algebra.f90}\strut
\end{minipage} & \begin{minipage}[t]{0.30\columnwidth}\raggedright
\texttt{docs/qr/src\_\_mod\_linear\_algebra.f90.png}\strut
\end{minipage}\tabularnewline
\begin{minipage}[t]{0.30\columnwidth}\raggedright
DVR\strut
\end{minipage} & \begin{minipage}[t]{0.30\columnwidth}\raggedright
\href{https://github.com/l1Ha/QuantumGeneralModule/blob/main/src/mod_dvr_grid.f90}{src/mod\_dvr\_grid.f90}\strut
\end{minipage} & \begin{minipage}[t]{0.30\columnwidth}\raggedright
\texttt{docs/qr/src\_\_mod\_dvr\_grid.f90.png}\strut
\end{minipage}\tabularnewline
\begin{minipage}[t]{0.30\columnwidth}\raggedright
波包传播\strut
\end{minipage} & \begin{minipage}[t]{0.30\columnwidth}\raggedright
\href{https://github.com/l1Ha/QuantumGeneralModule/blob/main/src/mod_wavepacket_propagator.f90}{src/mod\_wavepacket\_propagator.f90}\strut
\end{minipage} & \begin{minipage}[t]{0.30\columnwidth}\raggedright
\texttt{docs/qr/src\_\_mod\_wavepacket\_propagator.f90.png}\strut
\end{minipage}\tabularnewline
\begin{minipage}[t]{0.30\columnwidth}\raggedright
吸收边界\strut
\end{minipage} & \begin{minipage}[t]{0.30\columnwidth}\raggedright
\href{https://github.com/l1Ha/QuantumGeneralModule/blob/main/src/mod_absorbing_boundary.f90}{src/mod\_absorbing\_boundary.f90}\strut
\end{minipage} & \begin{minipage}[t]{0.30\columnwidth}\raggedright
\texttt{docs/qr/src\_\_mod\_absorbing\_boundary.f90.png}\strut
\end{minipage}\tabularnewline
\begin{minipage}[t]{0.30\columnwidth}\raggedright
Chebyshev\strut
\end{minipage} & \begin{minipage}[t]{0.30\columnwidth}\raggedright
\href{https://github.com/l1Ha/QuantumGeneralModule/blob/main/src/mod_chebyshev_propagator.f90}{src/mod\_chebyshev\_propagator.f90}\strut
\end{minipage} & \begin{minipage}[t]{0.30\columnwidth}\raggedright
\texttt{docs/qr/src\_\_mod\_chebyshev\_propagator.f90.png}\strut
\end{minipage}\tabularnewline
\begin{minipage}[t]{0.30\columnwidth}\raggedright
Lindblad\strut
\end{minipage} & \begin{minipage}[t]{0.30\columnwidth}\raggedright
\href{https://github.com/l1Ha/QuantumGeneralModule/blob/main/src/mod_open_quantum.f90}{src/mod\_open\_quantum.f90}\strut
\end{minipage} & \begin{minipage}[t]{0.30\columnwidth}\raggedright
\texttt{docs/qr/src\_\_mod\_open\_quantum.f90.png}\strut
\end{minipage}\tabularnewline
\begin{minipage}[t]{0.30\columnwidth}\raggedright
Krotov\strut
\end{minipage} & \begin{minipage}[t]{0.30\columnwidth}\raggedright
\href{https://github.com/l1Ha/QuantumGeneralModule/blob/main/src/mod_optimal_control.f90}{src/mod\_optimal\_control.f90}\strut
\end{minipage} & \begin{minipage}[t]{0.30\columnwidth}\raggedright
\texttt{docs/qr/src\_\_mod\_optimal\_control.f90.png}\strut
\end{minipage}\tabularnewline
\begin{minipage}[t]{0.30\columnwidth}\raggedright
定态散射\strut
\end{minipage} & \begin{minipage}[t]{0.30\columnwidth}\raggedright
\href{https://github.com/l1Ha/QuantumGeneralModule/blob/main/src/mod_ti_scattering.f90}{src/mod\_ti\_scattering.f90}\strut
\end{minipage} & \begin{minipage}[t]{0.30\columnwidth}\raggedright
\texttt{docs/qr/src\_\_mod\_ti\_scattering.f90.png}\strut
\end{minipage}\tabularnewline
\begin{minipage}[t]{0.30\columnwidth}\raggedright
含时散射\strut
\end{minipage} & \begin{minipage}[t]{0.30\columnwidth}\raggedright
\href{https://github.com/l1Ha/QuantumGeneralModule/blob/main/src/mod_td_scattering.f90}{src/mod\_td\_scattering.f90}\strut
\end{minipage} & \begin{minipage}[t]{0.30\columnwidth}\raggedright
\texttt{docs/qr/src\_\_mod\_td\_scattering.f90.png}\strut
\end{minipage}\tabularnewline
\begin{minipage}[t]{0.30\columnwidth}\raggedright
外场散射\strut
\end{minipage} & \begin{minipage}[t]{0.30\columnwidth}\raggedright
\href{https://github.com/l1Ha/QuantumGeneralModule/blob/main/src/mod_field_scattering.f90}{src/mod\_field\_scattering.f90}\strut
\end{minipage} & \begin{minipage}[t]{0.30\columnwidth}\raggedright
\texttt{docs/qr/src\_\_mod\_field\_scattering.f90.png}\strut
\end{minipage}\tabularnewline
\begin{minipage}[t]{0.30\columnwidth}\raggedright
转振光谱\strut
\end{minipage} & \begin{minipage}[t]{0.30\columnwidth}\raggedright
\href{https://github.com/l1Ha/QuantumGeneralModule/blob/main/src/mod_rovibrational.f90}{src/mod\_rovibrational.f90}\strut
\end{minipage} & \begin{minipage}[t]{0.30\columnwidth}\raggedright
\texttt{docs/qr/src\_\_mod\_rovibrational.f90.png}\strut
\end{minipage}\tabularnewline
\begin{minipage}[t]{0.30\columnwidth}\raggedright
光碎片\strut
\end{minipage} & \begin{minipage}[t]{0.30\columnwidth}\raggedright
\href{https://github.com/l1Ha/QuantumGeneralModule/blob/main/src/mod_photofragment_flux.f90}{src/mod\_photofragment\_flux.f90}\strut
\end{minipage} & \begin{minipage}[t]{0.30\columnwidth}\raggedright
\texttt{docs/qr/src\_\_mod\_photofragment\_flux.f90.png}\strut
\end{minipage}\tabularnewline
\begin{minipage}[t]{0.30\columnwidth}\raggedright
FSSH\strut
\end{minipage} & \begin{minipage}[t]{0.30\columnwidth}\raggedright
\href{https://github.com/l1Ha/QuantumGeneralModule/blob/main/src/mod_surface_hopping_fssh.f90}{src/mod\_surface\_hopping\_fssh.f90}\strut
\end{minipage} & \begin{minipage}[t]{0.30\columnwidth}\raggedright
\texttt{docs/qr/src\_\_mod\_surface\_hopping\_fssh.f90.png}\strut
\end{minipage}\tabularnewline
\begin{minipage}[t]{0.30\columnwidth}\raggedright
Jacobi 坐标\strut
\end{minipage} & \begin{minipage}[t]{0.30\columnwidth}\raggedright
\href{https://github.com/l1Ha/QuantumGeneralModule/blob/main/src/mod_triatomic_geometry.f90}{src/mod\_triatomic\_geometry.f90}\strut
\end{minipage} & \begin{minipage}[t]{0.30\columnwidth}\raggedright
\texttt{docs/qr/src\_\_mod\_triatomic\_geometry.f90.png}\strut
\end{minipage}\tabularnewline
\begin{minipage}[t]{0.30\columnwidth}\raggedright
超球反应\strut
\end{minipage} & \begin{minipage}[t]{0.30\columnwidth}\raggedright
\href{https://github.com/l1Ha/QuantumGeneralModule/blob/main/src/mod_hyperspherical_reactive.f90}{src/mod\_hyperspherical\_reactive.f90}\strut
\end{minipage} & \begin{minipage}[t]{0.30\columnwidth}\raggedright
\texttt{docs/qr/src\_\_mod\_hyperspherical\_reactive.f90.png}\strut
\end{minipage}\tabularnewline
\begin{minipage}[t]{0.30\columnwidth}\raggedright
三体复合\strut
\end{minipage} & \begin{minipage}[t]{0.30\columnwidth}\raggedright
\href{https://github.com/l1Ha/QuantumGeneralModule/blob/main/src/mod_three_body_recombination.f90}{src/mod\_three\_body\_recombination.f90}\strut
\end{minipage} & \begin{minipage}[t]{0.30\columnwidth}\raggedright
\texttt{docs/qr/src\_\_mod\_three\_body\_recombination.f90.png}\strut
\end{minipage}\tabularnewline
\begin{minipage}[t]{0.30\columnwidth}\raggedright
CIR\strut
\end{minipage} & \begin{minipage}[t]{0.30\columnwidth}\raggedright
\href{https://github.com/l1Ha/QuantumGeneralModule/blob/main/src/mod_confined_scattering.f90}{src/mod\_confined\_scattering.f90}\strut
\end{minipage} & \begin{minipage}[t]{0.30\columnwidth}\raggedright
\texttt{docs/qr/src\_\_mod\_confined\_scattering.f90.png}\strut
\end{minipage}\tabularnewline
\begin{minipage}[t]{0.30\columnwidth}\raggedright
Fano/CCR\strut
\end{minipage} & \begin{minipage}[t]{0.30\columnwidth}\raggedright
\href{https://github.com/l1Ha/QuantumGeneralModule/blob/main/src/mod_autoionization_fano.f90}{src/mod\_autoionization\_fano.f90}\strut
\end{minipage} & \begin{minipage}[t]{0.30\columnwidth}\raggedright
\texttt{docs/qr/src\_\_mod\_autoionization\_fano.f90.png}\strut
\end{minipage}\tabularnewline
\begin{minipage}[t]{0.30\columnwidth}\raggedright
交叉场\strut
\end{minipage} & \begin{minipage}[t]{0.30\columnwidth}\raggedright
\href{https://github.com/l1Ha/QuantumGeneralModule/blob/main/src/mod_crossed_field_scattering.f90}{src/mod\_crossed\_field\_scattering.f90}\strut
\end{minipage} & \begin{minipage}[t]{0.30\columnwidth}\raggedright
\texttt{docs/qr/src\_\_mod\_crossed\_field\_scattering.f90.png}\strut
\end{minipage}\tabularnewline
\begin{minipage}[t]{0.30\columnwidth}\raggedright
旋量 BEC\strut
\end{minipage} & \begin{minipage}[t]{0.30\columnwidth}\raggedright
\href{https://github.com/l1Ha/QuantumGeneralModule/blob/main/src/mod_spinor_bec.f90}{src/mod\_spinor\_bec.f90}\strut
\end{minipage} & \begin{minipage}[t]{0.30\columnwidth}\raggedright
\texttt{docs/qr/src\_\_mod\_spinor\_bec.f90.png}\strut
\end{minipage}\tabularnewline
\begin{minipage}[t]{0.30\columnwidth}\raggedright
表面散射\strut
\end{minipage} & \begin{minipage}[t]{0.30\columnwidth}\raggedright
\href{https://github.com/l1Ha/QuantumGeneralModule/blob/main/src/mod_surface_scattering.f90}{src/mod\_surface\_scattering.f90}\strut
\end{minipage} & \begin{minipage}[t]{0.30\columnwidth}\raggedright
\texttt{docs/qr/src\_\_mod\_surface\_scattering.f90.png}\strut
\end{minipage}\tabularnewline
\begin{minipage}[t]{0.30\columnwidth}\raggedright
Eley--Rideal\strut
\end{minipage} & \begin{minipage}[t]{0.30\columnwidth}\raggedright
\href{https://github.com/l1Ha/QuantumGeneralModule/blob/main/src/mod_surface_reaction_er.f90}{src/mod\_surface\_reaction\_er.f90}\strut
\end{minipage} & \begin{minipage}[t]{0.30\columnwidth}\raggedright
\texttt{docs/qr/src\_\_mod\_surface\_reaction\_er.f90.png}\strut
\end{minipage}\tabularnewline
\begin{minipage}[t]{0.30\columnwidth}\raggedright
电子摩擦\strut
\end{minipage} & \begin{minipage}[t]{0.30\columnwidth}\raggedright
\href{https://github.com/l1Ha/QuantumGeneralModule/blob/main/src/mod_surface_electronic_friction.f90}{src/mod\_surface\_electronic\_friction.f90}\strut
\end{minipage} & \begin{minipage}[t]{0.30\columnwidth}\raggedright
\texttt{docs/qr/src\_\_mod\_surface\_electronic\_friction.f90.png}\strut
\end{minipage}\tabularnewline
\begin{minipage}[t]{0.30\columnwidth}\raggedright
相对论原子\strut
\end{minipage} & \begin{minipage}[t]{0.30\columnwidth}\raggedright
\href{https://github.com/l1Ha/QuantumGeneralModule/blob/main/src/mod_relativistic_atomic.f90}{src/mod\_relativistic\_atomic.f90}\strut
\end{minipage} & \begin{minipage}[t]{0.30\columnwidth}\raggedright
\texttt{docs/qr/src\_\_mod\_relativistic\_atomic.f90.png}\strut
\end{minipage}\tabularnewline
\begin{minipage}[t]{0.30\columnwidth}\raggedright
RIXS\strut
\end{minipage} & \begin{minipage}[t]{0.30\columnwidth}\raggedright
\href{https://github.com/l1Ha/QuantumGeneralModule/blob/main/src/mod_resonant_xray_scattering.f90}{src/mod\_resonant\_xray\_scattering.f90}\strut
\end{minipage} & \begin{minipage}[t]{0.30\columnwidth}\raggedright
\texttt{docs/qr/src\_\_mod\_resonant\_xray\_scattering.f90.png}\strut
\end{minipage}\tabularnewline
\begin{minipage}[t]{0.30\columnwidth}\raggedright
Penning\strut
\end{minipage} & \begin{minipage}[t]{0.30\columnwidth}\raggedright
\href{https://github.com/l1Ha/QuantumGeneralModule/blob/main/src/mod_penning_associative_ionization.f90}{src/mod\_penning\_associative\_ionization.f90}\strut
\end{minipage} & \begin{minipage}[t]{0.30\columnwidth}\raggedright
\texttt{docs/qr/src\_\_mod\_penning\_associative\_ionization.f90.png}\strut
\end{minipage}\tabularnewline
\begin{minipage}[t]{0.30\columnwidth}\raggedright
强场 NSDI\strut
\end{minipage} & \begin{minipage}[t]{0.30\columnwidth}\raggedright
\href{https://github.com/l1Ha/QuantumGeneralModule/blob/main/src/mod_strong_field_nsdi.f90}{src/mod\_strong\_field\_nsdi.f90}\strut
\end{minipage} & \begin{minipage}[t]{0.30\columnwidth}\raggedright
\texttt{docs/qr/src\_\_mod\_strong\_field\_nsdi.f90.png}\strut
\end{minipage}\tabularnewline
\begin{minipage}[t]{0.30\columnwidth}\raggedright
里德堡阻塞\strut
\end{minipage} & \begin{minipage}[t]{0.30\columnwidth}\raggedright
\href{https://github.com/l1Ha/QuantumGeneralModule/blob/main/src/mod_rydberg_blockade.f90}{src/mod\_rydberg\_blockade.f90}\strut
\end{minipage} & \begin{minipage}[t]{0.30\columnwidth}\raggedright
\texttt{docs/qr/src\_\_mod\_rydberg\_blockade.f90.png}\strut
\end{minipage}\tabularnewline
\begin{minipage}[t]{0.30\columnwidth}\raggedright
FGH 示例\strut
\end{minipage} & \begin{minipage}[t]{0.30\columnwidth}\raggedright
\href{https://github.com/l1Ha/QuantumGeneralModule/blob/main/examples/ex01_fgh_diatomic_bound_states.f90}{examples/ex01\_fgh\_diatomic\_bound\_states.f90}\strut
\end{minipage} & \begin{minipage}[t]{0.30\columnwidth}\raggedright
\texttt{docs/qr/examples\_\_ex01\_fgh\_diatomic\_bound\_states.f90.png}\strut
\end{minipage}\tabularnewline
\begin{minipage}[t]{0.30\columnwidth}\raggedright
Split 示例\strut
\end{minipage} & \begin{minipage}[t]{0.30\columnwidth}\raggedright
\href{https://github.com/l1Ha/QuantumGeneralModule/blob/main/examples/ex03_split_operator_1d.f90}{examples/ex03\_split\_operator\_1d.f90}\strut
\end{minipage} & \begin{minipage}[t]{0.30\columnwidth}\raggedright
\texttt{docs/qr/examples\_\_ex03\_split\_operator\_1d.f90.png}\strut
\end{minipage}\tabularnewline
\begin{minipage}[t]{0.30\columnwidth}\raggedright
TI/TD 示例\strut
\end{minipage} & \begin{minipage}[t]{0.30\columnwidth}\raggedright
\href{https://github.com/l1Ha/QuantumGeneralModule/blob/main/examples/ex07_scattering_wavefunctions_ti_td.f90}{examples/ex07\_scattering\_wavefunctions\_ti\_td.f90}\strut
\end{minipage} & \begin{minipage}[t]{0.30\columnwidth}\raggedright
\texttt{docs/qr/examples\_\_ex07\_scattering\_wavefunctions\_ti\_td.f90.png}\strut
\end{minipage}\tabularnewline
\begin{minipage}[t]{0.30\columnwidth}\raggedright
Feshbach 示例\strut
\end{minipage} & \begin{minipage}[t]{0.30\columnwidth}\raggedright
\href{https://github.com/l1Ha/QuantumGeneralModule/blob/main/examples/ex08_ultracold_feshbach_segmented.f90}{examples/ex08\_ultracold\_feshbach\_segmented.f90}\strut
\end{minipage} & \begin{minipage}[t]{0.30\columnwidth}\raggedright
\texttt{docs/qr/examples\_\_ex08\_ultracold\_feshbach\_segmented.f90.png}\strut
\end{minipage}\tabularnewline
\begin{minipage}[t]{0.30\columnwidth}\raggedright
FSSH 示例\strut
\end{minipage} & \begin{minipage}[t]{0.30\columnwidth}\raggedright
\href{https://github.com/l1Ha/QuantumGeneralModule/blob/main/examples/ex30_tully_surface_hopping.f90}{examples/ex30\_tully\_surface\_hopping.f90}\strut
\end{minipage} & \begin{minipage}[t]{0.30\columnwidth}\raggedright
\texttt{docs/qr/examples\_\_ex30\_tully\_surface\_hopping.f90.png}\strut
\end{minipage}\tabularnewline
\begin{minipage}[t]{0.30\columnwidth}\raggedright
Python 可视化\strut
\end{minipage} & \begin{minipage}[t]{0.30\columnwidth}\raggedright
\href{https://github.com/l1Ha/QuantumGeneralModule/blob/main/python/pygenmod/visualizer.py}{python/pygenmod/visualizer.py}\strut
\end{minipage} & \begin{minipage}[t]{0.30\columnwidth}\raggedright
\texttt{docs/qr/python\_\_pygenmod\_\_visualizer.py.png}\strut
\end{minipage}\tabularnewline
\bottomrule
\end{longtable}

\hypertarget{ux7ed3ux8bed}{%
\section{结语}\label{ux7ed3ux8bed}}

本书的目标是把 GeneralModule
从``可编译的代码库''推进到``可验证、可复现、可教学、可扩展的科研资料''。后续版本应继续沿三个方向推进：

\begin{enumerate}
\def\labelenumi{\arabic{enumi}.}
\tightlist
\item
  理论上补齐更多严格推导和文献映射；
\item
  算法上扩展高阶方法和多维表示；
\item
  工程上完善测试、CI、可复现性和模块模板。
\end{enumerate}

若发现 Bug、公式错误、代码不一致或复现步骤缺失，请通过仓库 Issue
或以下联系方式反馈：

\begin{itemize}
\tightlist
\item
  维护者：李皓
\item
  电子邮箱：LIH\_ao@outlook.com
\item
  仓库：\url{https://github.com/l1Ha/QuantumGeneralModule}
\end{itemize}

\hypertarget{ux7b2cux5341ux4e00ux90e8ux5206-ux9ad8ux9636ux65b9ux6cd5ux4e0eux5b9eux73b0ux7ec6ux8282}{%
\section{第十一部分　高阶方法与实现细节}\label{ux7b2cux5341ux4e00ux90e8ux5206-ux9ad8ux9636ux65b9ux6cd5ux4e0eux5b9eux73b0ux7ec6ux8282}}

\hypertarget{krylov-ux5b50ux7a7aux95f4ux65b9ux6cd5}{%
\subsection{11.1　Krylov
子空间方法}\label{krylov-ux5b50ux7a7aux95f4ux65b9ux6cd5}}

当哈密顿矩阵过大而无法稠密对角化时，可以使用 Krylov 子空间方法。给定初态
\(|\psi_0\rangle\)，Krylov 子空间为

\[
\mathcal K_m(\hat H,|\psi_0\rangle)
=
\operatorname{span}
\{
|\psi_0\rangle,
\hat H|\psi_0\rangle,
\ldots,
\hat H^{m-1}|\psi_0\rangle
\}.
\]

Lanczos 算法构造正交基

\[
|\psi_0\rangle,|\psi_1\rangle,\ldots,|\psi_{m-1}\rangle,
\]

使得哈密顿量在该子空间内表示为三对角矩阵

\[
T_m
=
\begin{pmatrix}
\alpha_1 & \beta_1 &        &        \\
\beta_1 & \alpha_2 & \beta_2 &        \\
        & \beta_2 & \ddots & \ddots \\
        &        & \ddots & \alpha_m
\end{pmatrix}.
\]

于是

\[
|\psi(t)\rangle
\approx
V_m
\exp(-i T_m t/\hbar)
e_1,
\]

其中 \(V_m\) 的列向量是 Krylov 基。

Lanczos
方法适合求极端本征值和时间演化，但对重复本征值和近简并子空间需要块方法或重启策略。

\hypertarget{ux9ad8ux9636ux5206ux88c2ux7b97ux7b26}{%
\subsection{11.2　高阶分裂算符}\label{ux9ad8ux9636ux5206ux88c2ux7b97ux7b26}}

二阶 Strang 分裂为

\[
S_2(\Delta t)
=
e^{-i\hat V\Delta t/(2\hbar)}
e^{-i\hat T\Delta t/\hbar}
e^{-i\hat V\Delta t/(2\hbar)}.
\]

四阶方法可以使用 Suzuki 分数步长：

\[
S_4(\Delta t)
=
S_2(p_1\Delta t)
S_2(p_2\Delta t)
S_2(p_1\Delta t),
\]

其中

\[
p_1=\frac{1}{2-2^{1/3}},
\qquad
p_2=1-2p_1.
\]

更高阶方法虽然每步代价增加，但对长时间积分和严格要求幺正性的问题更有效。

\hypertarget{ux7a00ux758fux77e9ux9635ux4e0eux5b58ux50a8ux7b56ux7565}{%
\subsection{11.3　稀疏矩阵与存储策略}\label{ux7a00ux758fux77e9ux9635ux4e0eux5b58ux50a8ux7b56ux7565}}

在 DVR 中，若只保留非零矩阵元，可显著降低存储需求。设平均每行非零元数为
\(z\)，则稀疏存储需求为

\[
O(zN),
\]

而稠密存储需求为

\[
O(N^2).
\]

当 \(z\ll N\) 时，稀疏矩阵向量乘法复杂度为

\[
O(zN),
\]

远低于稠密矩阵向量乘法的 \(O(N^2)\)。

适合稀疏存储的结构包括：

\begin{itemize}
\tightlist
\item
  局域势能的对角矩阵；
\item
  近邻耦合动能矩阵；
\item
  分块对角角动量矩阵；
\item
  单电子耦合算符。
\end{itemize}

\hypertarget{ux5f20ux91cfux5206ux89e3ux4e0eux9ad8ux7ef4ux95eeux9898}{%
\subsection{11.4　张量分解与高维问题}\label{ux5f20ux91cfux5206ux89e3ux4e0eux9ad8ux7ef4ux95eeux9898}}

高维问题的网格数为

\[
N_{\mathrm{tot}}=\prod_{\alpha=1}^{d}N_\alpha.
\]

直接稠密存储代价为

\[
O(N^{2d}).
\]

张量分解将高阶张量表示为低秩因子乘积，例如 Tucker 分解：

\[
A_{i_1\ldots i_d}
=
\sum_{r_1,\ldots,r_d}
G_{r_1\ldots r_d}
U^{(1)}_{i_1r_1}
\cdots
U^{(d)}_{i_dr_d}.
\]

Tensor-train 表示则为

\[
A_{i_1\ldots i_d}
=
G_1(i_1)
G_2(i_2)
\cdots
G_d(i_d).
\]

若张量秩较低，存储和计算代价可从指数级降低到多项式级。

\hypertarget{ux591aux901aux9053ux6563ux5c04ux7684ux5339ux914dux4e0eux6e10ux8fd1ux5c55ux5f00}{%
\subsection{11.5　多通道散射的匹配与渐近展开}\label{ux591aux901aux9053ux6563ux5c04ux7684ux5339ux914dux4e0eux6e10ux8fd1ux5c55ux5f00}}

在匹配半径 \(R_m\) 处，径向波函数可写成入射波和出射波叠加：

\[
\mathbf F(R)
=
\mathbf J(R)\mathbf A
+
\mathbf N(R)\mathbf B,
\]

其中 \(\mathbf J\) 和 \(\mathbf N\)
分别是规则和不规则径向函数矩阵。定义对数导数

\[
\mathbf Y(R_m)
=
\mathbf F'(R_m)\mathbf F^{-1}(R_m).
\]

由 \(\mathbf Y\) 可以构造 \(\mathbf K\) 矩阵：

\[
\mathbf K
=
\left[
\mathbf J'(R_m)-\mathbf Y(R_m)\mathbf J(R_m)
\right]^{-1}
\left[
\mathbf Y(R_m)\mathbf N(R_m)-\mathbf N'(R_m)
\right].
\]

再由 Cayley 变换得到散射矩阵：

\[
\mathbf S
=
(\mathbf I+i\mathbf K)(\mathbf I-i\mathbf K)^{-1}.
\]

\hypertarget{ux5206ux6bb5ux7f51ux683cux7684ux8befux5deeux63a7ux5236}{%
\subsection{11.6　分段网格的误差控制}\label{ux5206ux6bb5ux7f51ux683cux7684ux8befux5deeux63a7ux5236}}

分段网格将径向域划分为若干扇区：

\[
[r_{\min},r_{\max}]
=
\bigcup_{s=1}^{S}[r_s,r_{s+1}].
\]

每个扇区可使用不同步长 \(\Delta r_s\)。若局部误差为

\[
\epsilon_s
\approx
C_s\Delta r_s^{p},
\]

则总误差可近似为

\[
\epsilon_{\mathrm{tot}}
\approx
\sum_s\epsilon_s.
\]

为了等误差分布，可令

\[
C_s\Delta r_s^p
=
\mathrm{const}.
\]

这样在强变化区域使用细网格，在平滑区域使用粗网格。

\hypertarget{fssh-ux7684ux7edfux8ba1ux8befux5dee}{%
\subsection{11.7　FSSH
的统计误差}\label{fssh-ux7684ux7edfux8ba1ux8befux5dee}}

FSSH 是随机轨迹方法。若透射概率为 \(p\)，轨迹数为 \(N\)，则标准误差为

\[
\sigma_p
=
\sqrt{\frac{p(1-p)}{N}}.
\]

更稳健的置信区间可使用 Wilson 区间估计：

\[
\frac{
\hat p+\frac{z^2}{2N}
\pm
z\sqrt{
\frac{\hat p(1-\hat p)}{N}
+
\frac{z^2}{4N^2}
}
}
{
1+\frac{z^2}{N}
},
\]

其中 \(z\) 是正态分位数。对 95\% 置信区间，\(z\approx1.96\)。

增加轨迹数时，统计误差按 \(N^{-1/2}\) 下降。

\hypertarget{ux5f00ux653eux7cfbux7edfux7684ux8c31ux65b9ux6cd5}{%
\subsection{11.8　开放系统的谱方法}\label{ux5f00ux653eux7cfbux7edfux7684ux8c31ux65b9ux6cd5}}

Lindblad 主方程为

\[
\dot\rho
=
\mathcal L\rho.
\]

若 \(\mathcal L\) 不含时，则形式解为

\[
\rho(t)=e^{\mathcal L t}\rho(0).
\]

可以把超算符 \(\mathcal L\) 表示为矩阵并做谱分解：

\[
\mathcal L
=
\sum_\lambda
\lambda |R_\lambda\rangle\langle L_\lambda|,
\]

其中 \(|R_\lambda\rangle\) 和 \(\langle L_\lambda|\)
分别是右本征算符和左本征算符。于是

\[
\rho(t)
=
\sum_\lambda
e^{\lambda t}
\langle L_\lambda|\rho(0)\rangle
|R_\lambda\rangle.
\]

实部为负的本征值对应衰减模式，纯虚本征值对应相干振荡。

\hypertarget{ux91cfux5b50ux63a7ux5236ux4e2dux7684ux7ea6ux675f}{%
\subsection{11.9　量子控制中的约束}\label{ux91cfux5b50ux63a7ux5236ux4e2dux7684ux7ea6ux675f}}

最优控制不仅要最大化保真度，还要满足物理约束：

\begin{itemize}
\tightlist
\item
  峰值场强 \(|\epsilon(t)|\le\epsilon_{\max}\)；
\item
  频谱带宽约束；
\item
  脉冲面积约束；
\item
  系统是否能实现目标态；
\item
  避免激发不可控态；
\item
  避免超过损伤阈值。
\end{itemize}

常见惩罚泛函为

\[
J_\epsilon
=
\int_0^T
\lambda(t)\epsilon^2(t)\,dt.
\]

若使用硬约束，则可采用投影或罚函数方法。

\hypertarget{ux6570ux503cux5b9eux9a8cux5efaux8bae}{%
\subsection{11.10　数值实验建议}\label{ux6570ux503cux5b9eux9a8cux5efaux8bae}}

\begin{longtable}[]{@{}ll@{}}
\toprule
主题 & 建议实验\tabularnewline
\midrule
\endhead
Krylov & 比较不同子空间维数下的传播误差\tabularnewline
高阶分裂 & 比较 \(S_2\) 与 \(S_4\) 的 \(\Delta t\) 标度\tabularnewline
稀疏矩阵 & 对比稠密与稀疏存储的内存和时间\tabularnewline
张量分解 & 在二维势能面上测试秩与误差关系\tabularnewline
多通道散射 & 改变匹配半径，检查 \(S\) 矩阵稳定性\tabularnewline
分段网格 & 改变局部步长，检查总误差\tabularnewline
FSSH & 增加轨迹数，绘制置信区间\tabularnewline
Lindblad & 改变耗散率，检查谱分解\tabularnewline
最优控制 & 改变惩罚参数，比较保真度和场强\tabularnewline
\bottomrule
\end{longtable}

\hypertarget{ux53c2ux8003ux6587ux732e-4}{%
\subsection{11.11　参考文献}\label{ux53c2ux8003ux6587ux732e-4}}

\begin{enumerate}
\def\labelenumi{\arabic{enumi}.}
\tightlist
\item
  G. H. Golub and C. F. Van Loan, \emph{Matrix Computations}, 4th ed.,
  Johns Hopkins University Press, Baltimore, 2013.
\item
  Y. Saad, \emph{Iterative Methods for Sparse Linear Systems}, 2nd ed.,
  SIAM, Philadelphia, 2003.
\item
  M. H. Kalos and P. A. Whitlock, \emph{Monte Carlo Methods}, 2nd ed.,
  Wiley-VCH, Weinheim, 2008.
\item
  D. J. Tannor, \emph{Introduction to Quantum Mechanics: A
  Time-Dependent Perspective}, University Science Books, Sausalito,
  2007.
\item
  GeneralModule
  源码与文档：\url{https://github.com/l1Ha/QuantumGeneralModule}
\end{enumerate}

\end{document}
